% concatenated from main.tex

\RequirePackage{rotating}
\documentclass[a4paper, english, 9pt]{amsart}

%PER STAMPARE
%
\headheight=0in
\headsep =-0.54in
\topmargin=0in
\textheight=10.52in
\textwidth=7.7in
\oddsidemargin=-0.63in
\evensidemargin=-0.63in
%%
\usepackage{listings}
\usepackage{rotating}

\usepackage{amsthm,amssymb,url,hyperref}
\usepackage{fullpage}
\usepackage{MnSymbol}
\usepackage[all]{xy}
\usepackage{tikz}
\usepackage{subfig} 
%\captionsetup[figure]{width=0.4\textwidth}
\usepackage{caption}
\usepackage{forest}
\usepackage{color}
\usepackage{amsaddr}
\usepackage{mathtools}
\usepackage{enumitem}
\usepackage{hyperref}
\usepackage{booktabs} % per \toprule, \midrule, \bottomrule

\theoremstyle{plain}
\newtheorem{theorem}{Theorem}[section]
\newtheorem{acknowledgement}{Acknowledgement}
\newtheorem{axiom}{Axiom}
\newtheorem{case}{Case}
\newtheorem{fact}{Fact}

\newtheorem*{theorem*}{Theorem}
\newtheorem{claim}{Claim}
\newtheorem{conclusion}{Conclusion}
\newtheorem{condition}[theorem]{Condition}
\newtheorem{conjecture}[theorem]{Conjecture}
\newtheorem{corollary}[theorem]{Corollary}
\newtheorem{criterion}{Criterion}
\newtheorem{exercise}{Exercise}
\newtheorem{lemma}[theorem]{Lemma}
\newtheorem{notation}[theorem]{Notation}
\newtheorem{proposition}[theorem]{Proposition}

\newtheorem{observation}[theorem]{Observation}
\newtheorem{solution}{Solution}
\newtheorem{summary}{Summary}
\newtheorem{construction}[theorem]{Construction}
\newtheorem{question}{Question}


\theoremstyle{definition}
\newtheorem{definition}[theorem]{Definition}
\newtheorem{problem}[theorem]{Problem}
\newtheorem{example}[theorem]{Example}
\newtheorem{remark}[theorem]{Remark}
\newtheorem{algorithm}[theorem]{Algorithm}

\def\ld{\backslash}
\def\img#1{\mathrm{Im}(#1)}
\def\ker#1{\mathrm{ker}(#1)}
\def\Ker#1{\mathrm{Ker}(#1)}
\def\aut#1{\mathrm{Aut}(#1)}
\def\Aut#1{\mathrm{Aut}(#1)}
\def\End#1{\mathrm{End}(#1)}
\def\aff#1{\mathrm{Aff}#1}
\def\Aff#1{\mathrm{Aff}#1}
\def\lmlt{\mathrm{LMlt}}
\def\dis{\mathrm{Dis}}
\def\fix{\mathrm{Fix}}
\def\sym{\mathrm{Sym}}
\def\syl{\mathrm{Syl}}
\def\Sym{\mathrm{Sym}}
\def\map{\mathrm{Map}}
\def\con{\mathrm{Con}}
\def\lcm{\mathrm{lcm}}
\def\n{\mathrm{N}}
\def\comment#1{{\color{red} #1}}
\def\c#1{\mathrm{con}_{#1}}
\def\Q{\mathcal{Q}_{\mathrm{Hom}}}


\def\setof#1#2{\{#1\,\colon\,#2\}}
\def\A{\mathbf A}
\def\Z{\mathbb Z}
\def\F{\mathbb F}
\def\N{\mathrm{Norm}}

\def\Ma{\mathcal Max}

\def\ma{\mathcal Min}

\def\Q{\mathcal Q}
\def\TC{\mathrm{TC}}
\def\LSS{\mathcal{LSS}}
\def\vv#1#2{\left(\begin{matrix} #1 \\ #2\end{matrix}\right)}
\def\m#1#2#3#4{\left(\begin{smallmatrix} #1 & #2 \\ #3 & #4\end{smallmatrix}\right)}
\def\mm#1#2#3#4{\left(\begin{matrix} #1 & #2 \\ #3 & #4\end{matrix}\right)}
%\def\cg#1{\cg\alpha }
\def\cg#1{\equiv_\alpha}
\newcommand*\xbar[1]{%
   \hbox{%
     \vbox{%
       \hrule height 0.5pt % The actual bar
       \kern0.5ex%         % Distance between bar and symbol
       \hbox{%
         \kern-0.1em%      % Shortening on the left side
         \ensuremath{#1}%
         \kern-0.1em%      % Shortening on the right side
       }%
     }%
   }%
} 
\def\fs#1{{\color{cyan} #1}}

\title{latin quandles of size $16p$}
\author{Marco Bonatto \and Filippo Spaggiari}
\email[1]{marco.bonatto.87@gmail.com}
\email[2]{spaggiari@karlin.mff.cuni.cz}




\begin{document}


\begin{abstract}
    In this paper we obtain the classification of latin quandles of size $16p$ where $p$ is an odd prime. In particular we have that such quandles are always subdirectly reducible but in some cases for small primes. Specifically, we have that latin quandles of size $16p$ are affine if $p\neq 1 \pmod{3}$ or $p\neq 3,5$. If $p=1\pmod 3$ there are $2$ subdirectly reducible non directly decomposable latin quandles of size $16p$ for every prime with $p=1\pmod{3}$ and there are one subdirectly irreducible latin quandle of size $16p$ for $p=3,5$. We provide explicit constructions as coset quandles for all the quandles mentioned above.
\end{abstract}
\maketitle




\section*{Introduction}

Quandles are binary algebras introduced in the framework of knot theory to construct knot invariants \cite{J, Matveev}. Beyond their origins, they have proven to be deeply connected to various other mathematical structures, such as solutions to the set-theoretical quantum Yang-Baxter equation \cite{etingof2001indecomposable, etingof1999set} and pointed Hopf algebras \cite{AG}. These connections underline the importance of quandles as versatile and foundational objects in both algebra and topology.

Quandles have been extensively studied through both group and module theory. In addition, universal algebra offers an important framework for investigating the complex relationship between properties of the congruence lattice and those of the so-called displacement group, which encodes many properties of the quandle itself. Indeed, there are two Galois connections between the congruence lattice of a quandle and the lattice of normal subgroups of the left multiplication group contained in the displacement group. This rich interplay between the two lattices has proven crucial in addressing classification problems for finite quandles, facilitating a deeper understanding of their structure and properties. Moreover, it has also been fundamental in developing commutator theory for quandles. Commutator theory, initially developed for general algebraic structures in \cite{comm}, has been specialized for quandles in \cite{CP}.

In this paper, we use this set of tools to obtain the classification of latin quandles of size $16p$, where $p$ is a prime, taking advantage of the classification of latin quandles of size $4,16,p$, and $4p$ already available \cite{ESS, LSS, RIG}, as well as the homogeneous representation of connected quandles over the displacement group \cite{J}. Using this machinery, we can translate the classification problem into a purely group-theoretical one. Universal algebraic techniques are primarily used to identify the groups of interest for obtaining a homogeneous representation, assuming that the congruence lattice of the structure is known.

The core of the paper lies in Section \ref{sec:latin16p}, where we also collect existing results on the classification of small latin quandles and prove the main classification results for latin quandles of size $16p$ (completed in Section \ref{computational} with some computational results concerning the cases $p=3,5$). This classification, achieved through a detailed case analysis, demonstrates the power of the aforementioned techniques. Specifically, the classification strategy proceeds as follows:
\begin{itemize}
\item[(i)] Determine the structure of the congruence lattice;
\item[(ii)] Use the Galois correspondence between the congruence lattice and the lattice of normal subgroups of the left multiplication group to identify the displacement group;
\item[(iii)] Establish whether the groups satisfying the required properties exist and, if they do, whether their automorphisms give rise to coset quandles with the desired properties (in this case, being latin of size $16p$).
\end{itemize}
In the paper, we take advantage of the fact that finite latin quandles are solvable (in the sense of \cite{comm}) and so the description of finite simple latin quandles is completely understood. In particular, such quandles are affine over elementary abelian $p$-groups. 

The algorithm for classification is used to investigate directly indecomposable quandles. On the other hand, directly decomposable latin quandles of size $16p$ are known since the quandles of size $4,16,p$ and $4p$ are classified. If $p\neq 1\pmod{3}$ the directly decomposable latin quandles of size $16p$ are affine, as they are the direct product of the $9$ latin quandles of size $16$ and one of the $p-2$ latin quandles of size $p$ and such quandles are affine. If $p= 1\pmod{3}$ we also have the direct product of the unique latin quandle of size $4$ and one of the two non-affine latin quandles of size $4p$ described in \cite{LSS}.



Combining the computational results from Section \ref{computational}, Proposition \ref{group K}, and Theorem \ref{sr quandles} (in which we also provide an explicit coset representation for the quandles obtained) with the enumeration mentioned above, we achieve our main classification results, summarized in Table \eqref{enum}.



%\begin{theorem*}[Classification of latin quandles of order $16p$]
%Let $p$ be a prime such that $p=1\pmod{3}$. There exist exactly two directly indecomposable latin quandles of order $16p$. Such quandles are subdirectly reducible. Moreover, there exists a unique subdirectly irreducible latin quandle of size $16p$ for $p=3,5$.



\begin{table}[ht]
  \centering
  \begin{tabular}{cccc} 
    \toprule
    $p$ & Subdirectly irreducible & Directly decomposable (DD) & Subdirectly reducible (not DD) \\
    \midrule
    $p=1\pmod{3}$ & 0 & $9(p-2)+2$ & $2$ \\
    $p\neq 1\pmod{3}, \, p\neq 3,5 $ & 0 & $9(p-2)$ & 0 \\
    3 & 1 & 9 & 0 \\
    5 & 1 & 27 & 0 \\
    \bottomrule
  \end{tabular}\caption{Enumeration of latin quandles of size $16p$.}\label{enum}
\end{table}




%
%There exist exactly two latin quandles of order $16p$ for $p=3,5$, constructed as coset quandles 
%\[
%Q_3 = \mathcal{Q}(G_3, \fix(f_3), f_3), \quad Q_5 = \mathcal{Q}(G_5, \fix(f_5), f_5),
%\]
%respectively over the finite groups
%\begin{align*}
%G_3 &= \langle a,b,c \mid a^{3} = c^{2} = b^{2} = (cb)^{2} = (ba)^{3} = (ca)^{3} = (a^{-1} c a b)^{2} = (c a^{-1} b a)^{2} = 1 \rangle \cong \mathbb{Z}_2^4 \rtimes \mathbb{Z}_3, \\
%G_5 &= \langle a,b \mid a^{5} = b^{2} = (b a^{-1} b a)^{2} = (b a^{-1})^{5} = (b a^{-2} b a^{2})^{2} = 1 \rangle \cong \mathbb{Z}_2^4 \rtimes \mathbb{Z}_5,
%\end{align*}
%where the automorphisms $f_3 \in \aut{G_3}$ and $f_5 \in \aut{G_5}$ act on the generators as follows:
%\begin{align*}
 %   f_3 &: a \mapsto a^2 b a^{-1} b a, \quad b \mapsto b a^{-1} c a, \quad c \mapsto b c, \\
  %  f_5 &: a \mapsto a^{2} (a^{2} b)^{2} a^{-3} b a, \quad b \mapsto b a^{2} b a^{-2}.
  %  \end{align*}
%Both quandles $Q_3$ and $Q_5$ are subdirectly irreducible and their congruence lattice is a 3-element chain.
%\end{theorem*} 

In detail, the structure of the paper is as follows. Section \ref{sec:preliminary} collects foundational concepts from group theory and universal algebra. These preliminaries establish the necessary background for subsequent sections and include results concerning the congruence lattice of specific structures. For additional information on universal algebra, the reader may refer to \cite{bergman2011universal}, which provides an accessible introduction to the subject.

Section \ref{sec:quandles} introduces essential notations and concepts from quandle theory, including related groups of automorphisms. This section also presents key theorems addressing isomorphism problems and connections between the normal subgroups of the displacement group and the congruences of the quandle, laying the groundwork for the main results of the paper.

Section \ref{sec:latin16p} contains all the classification results, including the computational ones.

Finally, Section \ref{sec:appendix} presents the purely group-theoretical results used throughout the paper. %For example, we provide an explicit analysis of the $q$-Sylow subgroups of $GL_n(p)$, based on classical results from \cite{carter1964sylow} and \cite{weir1955sylow}, as summarized in \cite{schmidtsylow}. This analysis, and the way it can be exploited to understand the representation of nilpotent groups, is key to certain group-theoretical arguments used in Section \ref{sec:latin16p}, where these results are applied to address specific aspects of the classification problem.



\section{Preliminary results} \label{sec:preliminary}
\subsection{Group Theory}


We denote by $\Sym_Q$ the symmetric group over a non-empty set $Q$. An \emph{action} of a group $G$ on a set $Q$ is a group homomorphism $$\rho:G\longrightarrow \Sym_Q,\quad g\mapsto \rho_g.$$ We say that an action is \emph{faithful} if the map $\rho$ is injective, i.e. $\ker{\rho}=1$. For $x\in Q$, we denote by $G_x=\{g\in G\,\colon\, \rho_g(x)=x\}$ the \emph{stabilizer} of $x$ in $G$, and by $x^G=\{\rho_g(x)\,\colon\ g\in G\}$ the \emph{orbit} of $x$ with respect to the action of $G$.


We denote by $\aut{G}$ the group of automorphisms of a group $G$. The \emph{inner automorphism} with respect to $g\in G$ is the automorphism $\widehat{g}$ defined by setting 
\[
\widehat{g}(x) = gxg^{-1}, \quad \text{for every } x \in G.
\]
 
 If $h\in \aut{G}$ and $N$ is a $h$-invariant normal subgroup of $G$ (i.e. $h(N)=N$) then $h$ induces an automorphism on the factor $G/N$.% as $h(xN)=h(x)N$ for every $x\in G$.

\begin{lemma}
    Let $N$ be a normal subgroup of a group $G$, and let $f\in\aut{G}$ be such that $f(N)=N$. Then the map $$f_N:G/N\longrightarrow G/N,\quad xN \mapsto f(x)N$$ is a well defined automorphism of $G/N$.
\end{lemma}
%\comment{Used? Move?}\fs{non mi sembra lo usiamo} \comment{lo usiamo per definire i quozienti dei quandles. Teniamo il lemma o lo mettiamo solo in plain text?}\fs{magari allora teniamo il lemma che è scritto chiaro, ed eventualmente per future references}
\begin{proof}
    The map $xN\mapsto f(x)N$ is well defined. In fact, if $x^{-1}y\in N$ then $f(x^{-1}y)=f(x)^{-1}f(y)\in N$. Clearly $f$ is onto. Moreover 
    \begin{align*}
    f_N((xN)^{-1} yN)&=f_N(x^{-1}yN)=f(x^{-1}y)N=f(x^{-1})f(y)N\\&=f(x)^{-1}N f(y)N=f_N(xN)^{-1}f_N(yN).        
    \end{align*}
 Therefore $f_N\in \End{G/N}$. If $f_N(xN)=f(x)N=N$ then $f(x)\in N$ and $x\in N$ accordingly. So $f_N$ is injective.
\end{proof}

Let $N$ and $H$ be two groups, and let $\rho: H \to \aut{N}$ be a group action. The \emph{semidirect product} of $N$ and $H$ induced by $\rho$, denoted $N \rtimes_\rho H$, is the group $(N \times H,\cdot)$ where the operation $\cdot$ is defined by setting
\[
(n_1, h_1) \cdot (n_2, h_2) = (n_1 \rho_{h_1}(n_2), h_1 h_2).
\]
for every $n_1,n_2\in N$ and $h_1,h_2\in H$.



\subsection{Universal Algebra}


An {\it operation} of finite arity $n$ on a set $A$ is a map from $A^n$ to $A$. A (finitary) \emph{algebra} is a pair $\langle A, F\rangle$ where $A$ is a non-empty set and $F=\left\langle f_i: i \in I\right\rangle$ is a family of operations on $A$, indexed by some set $I$. The operations $\setof{f_i}{i\in I}$ are the  {\it basic operations} of $(A,F)$. The function $\rho: I \rightarrow \mathbb{N}$ which assigns to each $i \in I$ the arity of $f_i$ is called \emph{similarity type} of A. %Two algebras are called \emph{similar} if they have the same similarity type. 
A {\it term} is any meaningful composition of basic operations in $F$. A {\it Malt'sev term} is a ternary term $m(x,y,z)$ satisfying $m(y,x,x)\approx m(x,x,y)\approx y$. An algebra with a Malt'sev term is called a {\it Malt'sev algebra}.

%\emph{Congruence permutability} is a property of an algebra $(A, F)$ where any two congruences on $A$ commute with each other. Formally, if $\alpha$ and $\beta$ are two congruences on $A$, congruence permutability implies that $\alpha \circ \beta = \beta \circ \alpha$. 

%Algebras with this property are sometimes called \emph{Maltsev algebras}, and the congruence permutability condition is equivalent to have a \emph{Maltsev term}, that is, a ternary term $t(x,y,z)$ satisfying $t(y,x,x)\approx t(x,x,y)\approx y$.


%In this paper, we deal with algebraic structures with binary operations (namely functions from $A^2$ to $A$). 


A \emph{subalgebra of $A$} is a subset $B\subseteq A$ that is closed with respect to every basic operation in $F$.

A \emph{congruence} on an algebra $(A,F)$ is an equivalence relation $\theta$ on $A$ such that for every basic operation $f \in F$ of arity $n$ and every $a_i, b_i \in A$ with $a_i\, \theta\, b_i$ for $i = 1, \ldots, n$, we have $f(a_1, \ldots, a_n)\,\theta\, f(b_1, \ldots, b_n)$. The congruence lattice of an algebra $(A,F)$ is denoted by $\con(A,F)$ (or simply by $\con(A)$). We denote the congruence class of an element $x$ with respect to $\alpha\in \con(A)$ by $[x]_\alpha$ (or simply by $[x]$) and the set of classes by $A/\alpha$. The top element of $\con(A)$ is $1_A=A\times A$ and the bottom element of $\con(A)$ is the identity relation, denoted by $0_A$.


An algebra is said to be:

\begin{enumerate}
    \item \emph{directly decomposable} if it can be expressed as a direct product of two non-trivial algebras. Equivalently, an algebra is directly decomposable if and only if there are two congruences $\alpha$ and $\beta$ such that $\alpha\wedge\beta=0_A$ and $\alpha\circ\beta=1_A$ ($\circ$ is the usual composition of binary relations). If an algebra is not directly decomposable, it is called \emph{directly indecomposable};




%Formally, $(A,F)$ is directly decomposable if there exist non-empty algebras $(A_1, F_1)$ and $(A_2, F_2)$ such that $(A, F) \cong (A_1, F_1) \times (A_2, F_2)$. 
\item \emph{subdirectly reducible} if there exists $\setof{\alpha_i}{i\in I}\subseteq \con(A)\setminus \{0_A\}$ such that $\wedge_{i\in I} \alpha_i= 0_A$. Otherwise, $A$ is said to be {\it subdirectly irreducible}. It is known that subdirectly irreducible algebras are exactly those algebra having a smallest non-trivial congruence. Subdirectly irreducible algebras are clearly directly indecomposable;

\item \emph{simple} if it has no non-trivial proper congruences, i.e., the only congruences of $A$ are $0_A$ and $1_A$. %Simple algebras are important in algebraic theory because they cannot be further decomposed by congruences.
\end{enumerate}


A \emph{homomorphism} between two algebras $(A, F)$ and $(B, G)$ of the same similarity type is a function $\varphi: A \rightarrow B$ that preserves all operations in $F$. Specifically, for each $n$-ary basic operation $f_i \in F$, the function $\varphi$ satisfies $\varphi(f_i(a_1, \ldots, a_n)) = g_i(\varphi(a_1), \ldots, \varphi(a_n))$ for all $a_1, \ldots, a_n \in A$. If $\varphi$ is a bijection, then $\varphi$ is called an \emph{isomorphism}.



An \emph{automorphism} of an algebra $(A, F)$ is an isomorphism from the algebra to itself. The set of all automorphisms of an algebra with the composition of maps is a group called \emph{automorphism group} of the algebra. We denote the group of automorphisms of $(A,F)$ by $\aut{A,F}$ (or simply by $\aut{A}$).


The commutator theory for general algebras was introduced in \cite{comm} in order to extend the well-known concept of commutator of subgroups to the notion of commutator of congruences of any algebraic structure. The notion of {\it central} and {\it abelian} congruences is needed in the broader framework of general algebra to define {\it solvability} and {\it nilpotence} through the existence of special series of congruences named \emph{central series} and \emph{derived series} as for groups. To do so, given an algebra $A$ and $\alpha,\beta$ being congruences of $A$, the {\it commutator} $[\alpha,\beta]$ is a congruence of $A$ satisfying a certain term condition that for groups reduces to the usual notions of the commutator of normal subgroups. We are not going to explicitly describe such conditions and we refer the reader to \cite{comm} for further details (and to \cite{CP} for the case of quandles). Based on the notion of commutator of congruences, we say that $\alpha$ is {\it abelian} (resp. {\it central}) if $[\alpha,\alpha]=0_A$ (resp. $[\alpha,1_A]=0_A$). The largest central congruence is called the {\it center} of $A$ and is denoted by $\zeta_A$. 

The derived and central series are defined, respectively, as 
\begin{align*}
    &\gamma^0(A)=\gamma_0(A)=1_A,\quad \gamma^{i+1}(A)=[\gamma^i(A),\gamma^i(A)], \quad \gamma_{i+1}(A)=[\gamma_i(A),1_A],\text{ for } i\geq 0.
\end{align*}

An algebra is {\it solvable} (resp. {\it nilpotent}) of length $n$ if $\gamma^n(A)=0_A$ (resp. $\gamma_n(A)=0_A$). If $1_A$ is abelian, then $A$ is said to be {\it abelian} (equivalently $\gamma_1(A)=0_A$ or $\zeta_A=1_A$). The congruence $\gamma_1(A)$, simply denoted by $\gamma_A$, is the smallest congruence of $A$ with the abelian factor. Thus, given a congruence $\alpha$, the factor algebra $A/\alpha$ is abelian if and only if $\gamma_A \leq \alpha$. In particular, solvable simple algebras are abelian (indeed, if $A$ is simple and solvable then $\gamma_A\neq 1_A$ and so necessarily $\gamma_A=0_A$).

\subsection{Congruence slim algebras}\label{slim}

Let $A$ be an algebra. Recall that a {\it maximal} (resp. {\it minimal}) congruence is a congruence $\alpha\neq 0_A,1_A$ such that if $\alpha< \beta$ (resp. $\beta<\alpha$) then $\beta=1_A$ (resp. $\beta=0_A$). We denote by $\Ma_A$ (resp. $\ma_A$) the set of the maximal (resp. minimal) congruences of $A$ and we define 
\begin{align*}
\mu_A =\bigwedge_{\alpha\in \Ma_A} \alpha, \qquad
\nu_A =\bigvee_{\alpha\in \ma_A} \alpha.
\end{align*}



\begin{lemma}\label{cong of directly irred}
Let $A$ be a finite, directly indecomposable Malt'sev algebra. Then: 
\begin{enumerate}
    \item[(i)] $\nu_A\leq \mu_A$. 
    \item[(ii)] $\Ma_A\cap \ma_A\neq \emptyset$ if and only if $\con(A)$ is the three-element chain.
\end{enumerate}

\end{lemma}

\begin{proof}
(i) Let $\alpha\in \ma_A$ and $\beta\in \Ma_A$. Recall that for Malt'sev algebra we have $\alpha\circ \beta=\alpha\vee\beta$. Then $\alpha\wedge\beta\leq \alpha$ and $\beta\leq\alpha\circ \beta$. Hence, if $\alpha$ is not contained in $\beta$ we have that $\alpha\wedge\beta=0_A$ and $\alpha\circ\beta=1_A$ implying that $A$ is directly decomposable. Thus $\alpha\leq \beta$ for all $\alpha\in \ma_A$ and $\beta\in \Ma_A$, then accordingly $\nu_A\leq\mu_A$. 

(ii) Clearly, if $\con(A)$ is the chain of length $3$, then the unique proper congruence is both maximal and minimal. On the other hand, assume that $\alpha\in \Ma_A\cap \ma_A$. Then we have that
$$\nu_A\leq \mu_A\leq \alpha\leq \nu_A\leq \mu_A$$
and therefore $\nu_A= \mu_A=\alpha$, i.e. $\Ma_A=\ma_A=\{\alpha\}$. If $\beta$ is a congruence of $A$ then it is contained in a maximal congruence, and so $\beta\leq \alpha$. Hence, either $\beta=0_A$ or $\beta=\alpha$. So we see that $\alpha$ is the unique proper non-trivial congruence of $A$. 
%
%Let $\alpha\in \Ma_A\cap \ma_A$ and $\beta\in \con(Q)$. Then $\beta\wedge\alpha\leq \alpha$ and $\alpha\leq \alpha\vee\beta$. If $A$ is subdirectly irreducible then $\alpha=\alpha\wedge \beta$, so $\beta=1_A$. 
%
%If $\alpha\neq \beta\in \con(A)$ then $\alpha\wedge \beta=0_Q$ and $\alpha\vee\beta=1_Q$ and therefore $A$ is directly reducible, contradiction.
\end{proof}


Let $A$ be an algebra $A$. The \emph{height} of $\con(A)$ is the maximum of the chains' lengths of $\con(A)$. If $\con(A)=\Ma_A\cup \ma_A\cup \{0_A,1_A\}$ we say that $A$ is \emph{congruence slim}. We have that congruence slim algebras are those for which the height of $\con(A)$ is less than or equal to $4$.


\begin{lemma}\label{llength 4}
Let $A$ be a directly indecomposable algebra. The following are equivalent: 
\begin{itemize}
    \item[(i)] $A$ is congruence slim. %$\con(A)=\Ma_A\cup \ma_A\cup \{0_A,1_A\}$; 
    \item[(ii)] $\con(A)$ has height less than or equal to $4$. 
\end{itemize}    
\end{lemma}
%\comment{need finite for (ii)? }

\begin{proof}

(i) $\Rightarrow$ (ii) Let $0_A< \alpha_1< \ldots \alpha_n< 1_A$ be a chain of distinct congruences of maximal length. Then $\alpha_1\in \ma_A, \alpha_{2}\in \Ma_A$, so $n=2$ and then the maximal length of a chain of congruences is $4$.

(ii) $\Rightarrow$ (i) Let $\alpha$ be a congruence that is not in $\Ma_A\cup \ma_A\cup \{0_A,1_A\}$. Then there exists $\beta,\delta\in \con(A)$ such that $0_A<\beta<\alpha<\delta<1_A$. Then we have a chain of length $5$, which is a contradiction. %given by $$
%The maximal length of a chain in $\con(A)$ is $4$. Then if $\alpha\neq 0_A,1_A$ then $\alpha$ is either minimal or maximal.
\end{proof}



 
\begin{proposition}\label{lattices}
Let $A$ be a directly indecomposable, congruence slim Malt'sev algebra. Then one of the following holds:
\begin{itemize}
\item[(i)] $\con(A)$ is a chain of length less than or equal to $4$.
\item[(ii)] $\con(A)$ is
\begin{center}
		$\xymatrixrowsep{0.15in}
		\xymatrixcolsep{0.15in}
		\xymatrix{ 
			& & 1_A  \ar@{-}[d]\ar@{-}[dr]\ar@{-}[drr] \ar@{-}[dl] \ar@{-}[dll] & &\\
			\alpha_1 \ar@{-}[drr] & \alpha_2\ar@{-}[dr] &\alpha_3\ar@{-}[d] & \ldots\ar@{-}[dl] & \alpha_n \ar@{-}[dll]\\
			& & \mu_A\ar@{-}[d] &  &\\
			& & 0_A & &
		}$
		\end{center}
		for some $n\in \mathbb{N}$.
		
		\item[(iii)] $\con(A)$ is
\begin{center}
		$\xymatrixrowsep{0.15in}
		\xymatrixcolsep{0.15in}
		\xymatrix{ 
	& & 1_A\ar@{-}[d] & &\\		
			& & \nu_A  \ar@{-}[d]\ar@{-}[dr]\ar@{-}[drr] \ar@{-}[dl] \ar@{-}[dll] & &\\
			\beta_1 \ar@{-}[drr] & \beta_2\ar@{-}[dr] &\beta_3\ar@{-}[d] & \ldots\ar@{-}[dl] & \beta_n \ar@{-}[dll]\\
			& & 0_A &  &\\
		}$
		\end{center}
		for some $n\in \mathbb{N}$.
\end{itemize}
\end{proposition}


\begin{proof}
Recall that for Malt'sev algebras it holds that $\alpha\circ\beta=\alpha\wedge\beta$ for every pair of congruences $\alpha,\beta$. Let us denote $L=\con(A)$. According to Lemma \ref{llength 4}, we have $L=\Ma_A\cup \ma_A\cup \{0_A,1_A\}$. 
\begin{itemize}
    \item  If $|\Ma_A|=|\ma_A|=1$ then $L$ is a chain of length less than or equal to $4$. 

\item Assume that $|\Ma_A|>1$ and let $\alpha,\beta\in\Ma_A$ with $\alpha\neq \beta$. Note that $\alpha\circ \beta=1_A$. If $\alpha\wedge\beta=0_A$ then $A$ is directly decomposable, thus $0_A\neq \alpha\wedge\beta$ and, moreover, $\alpha\wedge\beta\in \ma_A$. Therefore, $\mu_A\leq \alpha\wedge\beta \leq \nu_A\leq \mu_A$, where the last inequality follows from Lemma \ref{cong of directly irred}(i). Hence $\mu_A=\nu_A\in \ma_A$. Thus, $A$ has a unique minimal congruence, namely $\mu_A$ and it follows that $L$ is as in (ii).

\item Assume that $|\ma_A|>1$ and let $\alpha,\beta\in \ma_A$ with $\alpha\neq \beta$. A dual argument shows $\alpha\vee\beta=\nu_A=\mu_A\in \Ma_A$ for every pair of distinct minimal congruences $\alpha,\beta$. So, $A$ has a unique maximal congruence, namely $\nu_A$ and so $L$ is as in (iii).\qedhere
\end{itemize}

%
%\begin{itemize}
 %   \item[(i)] If $|\Ma_L|=|\ma_L|=1$ then $L$ is a chain of length at most $4$.
%\item[(ii)] If $|\Ma_L|\geq 2$ then according to Corollary \ref{cor mu and nu}, $\alpha\wedge\beta=\mu_A\in \ma_A$ and so $\mu_A\leq \nu_A\leq \mu_A\in \ma_A$. Thus $\mu_A=\nu_A$ and $\nu_A$ is the unique minimal congruence of $A$, since $A$ is directly idecomposable. Thus $L$ is as in (ii).
%\item[(iii)] If $|\Ma_L|\geq 2$ a dual argument shows that $\nu_A=\mu_A$ and $\nu_A$ is the unique maximal congruence of $A$ and so $L$ is as in (iii).\qedhere
%\end{itemize}
%
\end{proof}






%\begin{lemma}\label{DP}
%Let $Q$ be a finite solvable algebra and $\alpha,\beta\in \Ma_Q$. Then $Q/(\alpha\wedge\beta)\cong Q/\alpha\times Q/\beta$.
%\end{lemma}

%\comment{Cite the results for Maltsev algebras??} \comment{NEEDED?}

%\begin{proof}
%Clearly $Q/(\alpha\wedge\beta)$ embeds subdirectly into $Q/\alpha\times Q/\beta$. Both factors are strictly simple (they are simple and solvable \comment{add ref}), therefore the block $[a]_{\alpha/(\alpha\wedge \beta)}$ maps onto $Q/\beta$ and then $|Q/(\alpha\wedge\beta)|=|Q/\alpha||[a]_{\alpha/(\alpha\wedge\beta)}|\geq |Q/\alpha||Q/\beta|$. Hence $Q/(\alpha\wedge\beta)\cong Q/\alpha\times Q/\beta$.
%\end{proof}


%\subsection{Commutator theory}

%Commutator theory for general algebras was introduced in \cite{comm} in order to extend well-known concepts in group theory to any algebraic structure. In the context of a general algebra $A$, we need the notion of {\it central} and {\it abelian} congruences in order to define {\it solvability} and {\it nilpotence} as for groups through the existence of special series of congruences named \emph{central series} and \emph{derived series}. To do so, given an algebra $A$ and $\alpha,\beta$ be congruences of $A$, the {\it commutator} $[\alpha,\beta]$ is a congruence of $A$ satisfying a certain terms condition that reduces to the usual notions of the commutator of normal subgroups. We are not going to explicitely describe such conditions and we refer the reader to \cite{comm} for that. Based on the notion of commutator of congruences, we say that $\alpha$ is abelian (resp. central) if $[\alpha,\alpha]=0_Q$ (resp. $[\alpha,1_A]=0_Q$). The largest central series is called the {\it center} of $A$ and denoted by $\zeta_A$. 

%The derived and central series are defined respectively as 
%\begin{align*}
%    &\gamma^0(A)=1_A,\quad \gamma^{i+1}(A)=[\gamma^i(A),\gamma^i(A)],\\
%    &\gamma_0(A)=1_A,\quad \gamma_{i+1}(A)=[\gamma_i(A),1_A],
%\end{align*}
%for $i\geq 0$.

%An algebra is solvable (resp. nilpotent) of length $n$ if $\gamma^n(A)=0_A$ (resp. $\gamma_n(A)=0_Q$). If $1_A$ is abelian, then $A$ is said to be abelian (equivalently $\gamma_1(A)=0_A$ of $\zeta_A=1_A$). The congruence $\gamma_1(A)$, denoted simply by $\gamma_A$, is the smallest congruence of $A$ such that the algebra $A/\gamma_A$ is abelian. Thus, if $\alpha$ is a congruence such that $A/\alpha$ is abelian, then $\alpha \le \gamma_A$. In particular, solvable simple algebras are abelian (indeed if $A$ is simple and solvable then $\gamma_A\neq 1_A$ and so necessarily $\gamma_A=1_A$).
For solvable algebras, we can obtain further information about the congruence lattice.
\begin{lemma}\label{gamma and max} Let $A$ be a solvable algebra. Then $\gamma_A\leq \mu_A$.
\end{lemma}

\begin{proof}
Since $A$ is solvable, the simple factors of $A$ are abelian, therefore $A/\alpha$ is abelian for all $\alpha\in\Ma_A$. So,  $\gamma_A\leq\alpha$ for every $\alpha\in \Ma_A$ and thus $\gamma_A\leq \mu_A$. %the quotient algebra $A/\mu_A$ embeds naturally into the product of abelian algebras $\prod_{\alpha\in\Ma_A} A/\alpha$, and so it is also abelian. Accordingly, $\gamma_A\leq \mu_A$.
\end{proof}

\begin{corollary}\label{mu and gamma}
Let $A$ be a finite directly indecomposable solvable non-abelian Malt'sev algebra. If $A$ is congruence slim and $|\Ma_A|> 1$ then 
$\gamma_A	=\mu_A$ 
%for every pair of distinct maximal congruence $\alpha,\beta$.
\end{corollary}

\begin{proof}
By Lemma \ref{gamma and max} and Proposition \ref{lattices}(ii), we have $0_A\neq \gamma_A\leq \mu_A\in \ma_A$. Therefore, equality holds.
\end{proof}


\section{Quandles} \label{sec:quandles}

\subsection{Left quasigroups and quandles}
A {\it left quasigroup} is a set $Q$ endowed with a pair of binary operations $F=\{*,\backslash\}$, such that 
\begin{align}\label{LQG}
x*(x\backslash y)=y=x\backslash(x*y)
\end{align}
hold for every $x,y\in Q$. We usually denote $*$ just by juxtaposition. Given a left quasigroup $Q$ we can define the left and right multiplication maps as
$$L_x:y\mapsto x*y,\quad R_x:y\mapsto y*x$$
for every $x\in Q$. According to axioms \eqref{LQG}, we see that the map $L_x$ is a bijection of $Q$ for every $x\in Q$. So we have that $x\backslash y=L_x^{-1}(y)$ for every $x,y\in Q$ and the {\it left multiplication group} of $Q$ as $\lmlt(Q)=\langle L_x\colon\, x\in Q\rangle$.


The {\it Cayley kernel} of $Q$ is the equivalence relation $\lambda_Q$ defined by setting $x\, \lambda_Q\, y$ if and only if $L_x=L_y$. We say that $Q$ is:
\begin{enumerate}[label=(\roman*)]
    \item {\it faithful} if $\lambda_Q=0_Q$;
    \item {\it connected} if the natural action of $\lmlt(Q)$ is transitive on $Q$;
    \item {\it latin} if $R_x$ is bijective for every $x\in Q$.
\end{enumerate}
Note that latin left quasigroups are faithful and connected.

A {\it quandle} $Q$ is a left quasigroup such that the following axioms hold:
\begin{enumerate}[label=(\roman*)]
    \item $x*x=x$ for all $x\in Q$;
    \item $x*(y*z)=(x*y)*(x*z)$ for all $x,y,z\in Q$.
%    \item for all $x,y\in Q$ there is a unique $z\in Q$ such that $x*z=y$.
\end{enumerate}
 %An important normal subgroup is the {\it displacement group} of $Q$, defined by $\dis(Q)=\langle L_xL_y^{-1}\colon\, x,y\in Q\rangle$.

Note that (ii) is equivalent to having $L_x\in \aut{Q}$ for every $x\in Q$ and accordingly $\lmlt(Q)\leq \aut{Q}$.
Quandles can be constructed using groups and automorphisms as follows: let $G$ be a group, $f\in \aut{G}$ and $H\leq Fix(f)=\setof{x\in G}{f(x)=x}$. The set $G/H$ endowed with the operation 
$$xH*yH=xf(x^{-1}y)H$$ 
for every $x,y\in G$ is a quandle. We denote such quandle by $\mathcal{Q}(G,H,f)$ and refer to it as {\it coset quandles}. If $H=1$ we say that the coset quandle is {\it principal} and if $G$ is abelian, we say that it is {\it affine}. In the last case, we use the notation $\aff(G,f)$.


\subsection{Congruences and subgroups} 




Let $Q$ be a quandle, $\alpha\in\con(Q)$ and $N$ be a normal subgroup of $\lmlt(Q)$. As introduced in \cite{CP}, we define the following equivalence relations and subgroups of the left multiplication group of $Q$: 
\begin{align*}
    \c{N} &= \{(x,y)\in Q\times Q\,\colon\, L_xL_y^{-1}\in N\}\in \con(Q); \\
    \mathcal{O}_N &= \{(x,y)\in Q\times Q\,\colon\, n(x)=y\,\text{ for some }n\in N\}\in \con(Q);\\% \text{ and }\mathcal{O}_N\circ \alpha=\alpha\circ \mathcal{O}_N \text{ for every }\alpha\in \con(Q);\\
    \lmlt^\alpha &= \{h\in\lmlt(Q)\,\colon\, h(x)\,\alpha\, x\,\text{ for all } x\in Q\}\trianglelefteq \lmlt(Q); \\
    \dis_\alpha &= \langle L_xL_y^{-1}\colon\, x\,\alpha\, y\rangle \trianglelefteq \lmlt(Q).
\end{align*}
In particular, we define the {\it displacement group} of $Q$ as $\dis(Q)=\dis_{1_Q}$ which has the property that $\lmlt(Q)=\dis(Q)\langle L_x\rangle$ for every $x\in Q$ and consequently the orbits of $\lmlt(Q)$ and of $\dis(Q)$ coincide \cite[Proposition 2.1]{hsv}.

Note that $\lmlt^\alpha$ can be defined in the very same way for left quasigroups and $\dis_\alpha$ can also be defined as the normal closure of the subgroup generated by $\setof{L_x L_y^{-1}}{x\, \alpha \, y}$ in the left multiplication group. In particular, note that $\alpha\leq \lambda_Q$ if and only if $\dis_\alpha=1$. If so, we say that $Q$ is a {\it cover} of $Q/\alpha$ (see \cite{Eisermann, covering_paper, Involutory} for further details on quandle covers). We can also define $\dis^\alpha=\dis(Q)\cap \lmlt^\alpha$, and it is easy to verify that 
\begin{align*}
    \dis_\alpha\dis_\beta&=\dis_{\alpha\vee \beta} \\
    \dis^{\alpha\wedge\beta}&=\dis^\alpha\cap \dis^\beta\\
    \c{N\cap M}&=\c{N}\wedge\c{M}\\
    \mathcal{O}_{NM}&=\mathcal{O}_N\circ \mathcal{O}_M=\mathcal{O}_N\vee\mathcal{O}_M
\end{align*}
where $\alpha,\beta\in \con(Q)$ and $N,M$ are normal subgroups of $\lmlt(Q)$ (see also \cite[Proposition 3.2]{CP} for more details).

Moreover we have that $\mathcal{O}_N\leq \c{N}$ and according to \cite[Proposition 3.5(1)]{CP} we also  have that $$\dis_{\mathcal{O}_N}\leq \dis_{\c{N}}\leq N\leq \dis^{\mathcal{O}_N}\leq \dis^{\c{N}}$$  for every normal subgroup $N$ of $\lmlt(Q)$. In general we have
$$\mathcal{O}_{\dis_\alpha}\leq \mathcal{O}_{\dis^\alpha}\leq \alpha\leq \c{\dis_\alpha}\leq \c{\dis^\alpha}$$
for every congruence $\alpha$, and for finite latin quandles all such congruences are equal \cite[Lemma 1.8]{Galois}.

The pair of operators $(\dis_*,\c{*})$ (resp. $(\dis^*,\mathcal{O}_*)$) provides a monotone Galois connection between the congruence lattice of a quandle and the lattice of normal subgroups of its left multiplication group. We refer to \cite{CP} for more details on the properties of such Galois connections that we are going to exploit in the paper.

Given a left quasigroup $Q$ and $\alpha\in \con(Q)$, the map
$$\pi_\alpha:\lmlt(Q)\longrightarrow \lmlt(Q/\alpha),\quad L_x\mapsto L_{[x]}$$
can be extended to a surjective group homomorphism with kernel $\lmlt^\alpha$. Clearly $\pi_\alpha$ restricts and co-restricts to the displacement groups, and the kernel of this restriction is $\dis^\alpha$.

The displacement group plays a very important role for quandles. In fact, it can be used to build connected quandles and its factors with the coset quandle construction over the displacement group and its quotients (see also \cite{hsv}, Theorem 4.1, and \cite{J}). 

\begin{proposition}[\protect{\cite[Lemma 1.4]{GB}}]
    Let $Q$ be a connected quandle, $\alpha\in \con(Q)$ and $x\in Q$. Then $Q/\alpha\cong \mathcal{Q}(\dis(Q)/\dis^\alpha,\dis(Q)_{[x]}/\dis^\alpha ,\left(\widehat{L_x}\right)_{\dis^\alpha})$, where $\dis(Q)_{[x]}=\setof{h\in \dis(Q)}{h(x)\, \alpha\, x}$.
\end{proposition}


The coset quandle construction over the displacement group allows also to identify the isomorphism classes of quandles. Indeed, we have the following isomorphism theorem. 
\begin{theorem}[\protect{\cite[Theorem 5.12]{GB}\cite[Corollary 3.26]{GiuThe}}] \label{iso theo}
    Let $G$ be a group, $f_i\in \aut{G}$ and let $Q_i=\mathcal{Q}(G,Fix(f_i),f_i)$ be such that $\dis(Q_i)=G$ for $i=1,2$. Then $Q_1\cong Q_2$ if and only if $f_1$ and $f_2$ are conjugate in $\aut{G}$.
\end{theorem}



Sylow subgroups of nilpotent normal subgroups of the left multiplication group of a quandle provide trivially intersecting congruences through the operator $\mathcal{O}_*$.
\begin{lemma}\label{Sylow}
    Let $Q$ be a finite connected faithful quandle, $N$ be a normal nilpotent subgroup of $\lmlt(Q)$ such that $N=\prod_i P_i$ where $P_i$ are the $p$-Sylow subgroups of $N$. Then $ \mathcal{O}_{P_i}\wedge \mathcal{O}_{P_j}=0_Q$ for all $i\neq j$ and $\circ_{i} \mathcal{O}_{P_i}=\mathcal{O}_N$.
\end{lemma}

\begin{proof}

The Sylow subgroups of $N$ are characteristic in $N$, so they are normal in $\lmlt(Q)$. Thus, $\mathcal{O}_{P_i}$ is a congruence of $Q$ for every $i$. Moreover $\mathcal{O}_{P_i}\wedge \mathcal{O}_{P_j}\leq \c{P_i}\wedge \c{P_j}=\c{P_i\cap P_j}=\c{1}=\lambda_Q=0_Q$, since $P_i\cap P_j=1$. 
%The quandle $Q$ connected, so it is congruence uniform (see Lemma 2.3 of \cite{Maltsev_paper}). So, if $\alpha\leq \beta$ then $|[x]_\alpha|$ divides $|[x]_\beta|$.
%
     Moreover $N=\prod_i P_i$ and so $\mathcal{O}_N=\circ_i \mathcal{O}_{P_i}$.
\end{proof}



\begin{corollary}\label{center of dis^alpha}
    Let $Q$ be a finite faithful quandle, and $\alpha\in \ma_Q$. The following are equivalent:
    \begin{itemize}
        \item[(i)] $Z(\dis^\alpha)\neq 1$.
        \item[(ii)] $\dis^\alpha$ is a $p$-group of nilpotency class less than or equal to $2$ for some prime $p$.
    \end{itemize}
\end{corollary}
\begin{proof}
    (ii) $\Rightarrow$ (i) Clear.

    (i) $\Rightarrow$ (ii) Let $Z=Z(\dis^\alpha)\neq 1$. The subgroups $Z$ is characteristic in $\dis^\alpha$ and so is normal in $\lmlt(Q)$. Then $\beta=\mathcal{O}_Z\leq \alpha$ is a proper congruence, and so $\alpha=\beta$. Hence $\dis^\alpha\leq \dis_{\c{Z}}\leq Z$. According to \cite[Propositon 3.3(1)]{CP} we have $[\dis^\alpha,\dis^\alpha]\leq \dis_\alpha\leq Z$ and so $\dis^\alpha$ has nilpotency class $2$. According to Lemma \ref{Sylow} the Sylow subgroups of $\dis^\alpha$ provide trivially intersecting congruence below $\alpha$. Since $\alpha$ is minimal, we see that $\dis^\alpha$ is a $p$-group for some prime $p$.
\end{proof}

In the final part of this section, we examine how subdirect reducibility and direct decomposability of left quasigroups and quandles relate to the corresponding properties of their left multiplication group and displacement group.
\begin{lemma}\label{DD LQGs}
Let $Q_1$ and $Q_2$ be left quasigroups. Then $\dis(Q_1\times Q_2)\cong \dis(Q_1)\times \dis(Q_2)$.
\end{lemma}
\begin{proof}
    Let $p_i:Q\longrightarrow Q_i$ be the canonical projection, and $\alpha_i=\ker{p_i}$ for $i=1,2$. Then $\dis^{\alpha_1}\cap \dis^{\alpha_2}=\dis^{\alpha_1\wedge\alpha_2}=\dis^{0_Q}=1$ and $\dis(Q)=\dis_{\alpha_1\vee \alpha_2}=\dis_{\alpha_1}\dis_{\alpha_2}\leq \dis^\alpha_1\dis^\alpha_2$. Hence $\dis(Q)\cong \dis^\alpha_1\times \dis^\alpha_2$. Since $\dis(Q)/\dis^\alpha_1\cong \dis(Q/\alpha_1)\cong \dis^{\alpha_2}$ and $\dis(Q)/\dis^\alpha_2\cong \dis(Q/\alpha_2)\cong \dis^{\alpha_1}$ the statement follows.
\end{proof}
%\begin{proof}
%Let $p_i:Q\longrightarrow Q_i$ be the canonical projection and $\alpha_i=\ker{p_i}$ for $i=1,2$. The map $$\phi:\dis(Q)\longrightarrow \dis(Q_1)\times \dis(Q_2), \quad g\mapsto (\pi_{\alpha_1}(g),\pi_{\alpha_2}(g))$$
%is a well-defined isomorphism of groups.
% \end{proof}

%\begin{proof}
%    Let $Q=Q_1\times Q_2$, we define:
%\begin{align*}
 %   H_1&=\langle h L_{(x_1,y)} L_{(x_2,y)}^{-1}h^{-1},\,x_1,x_2\in Q_1, \, y\in Q_2, \,h \in \dis(Q)\rangle,\\
  %  H_2&=\langle h L_{(x,y_1)} L_{(x,y_2)}^{-1}h^{-1},\,y_1,y_2\in Q_2, \, x\in Q_1,\, h\in \dis(Q)\rangle.
%\end{align*}
 %It is easy to check that $H_1$ and $H_2$ are trivially intersecting normal subgroups of $\dis(Q)$ and that $H_1H_2=\dis(Q)$. Note that if $g\in \dis(Q)$ then $g(a,b)=(g_1(a),g_2(b))$ for some $g_i\in \dis(Q_i)$. Moreover, the map
% $$\phi_i:H_i\longrightarrow \dis(Q_i), \quad g\mapsto g_i$$
% for $i=1,2$ is an isomorphism since $h(a,b)=(h_1(a),b)$ and $h'(a,b)=(a,h_2(b))$ for every $h\in H_1$, $h'\in H_2$ and $(a,b)\in Q_1\times Q_2$.
% \begin{align*} 
  %   L_{(x_1,y)} L_{(x_2,y)}^{-1}(a,b)&=(L_{x_1} L_{x_2}^{-1}(a),b),\\
  %   L_{(x,y_1)} L_{(x,y_2)}^{-1}(a,b)&=(a,L_{y_1} L_{y_2}^{-1}(b)),\\
  %   L_{(z,t)} L_{(x,y_1)} L_{(x,y_2)}^{-1}L_{(z,t)}^{-1}(a,b)&=(a,L_t L_{y_1} L_{y_2}^{-1}L_t^{-1}(b))=(a, L_{L_t(y_1)} L_{L_t(y_2)}^{-1}(b)),\\
 %    L_{(z,t)} L_{(x_1,y)} L_{(x_2,y)}^{-1}L_{(z,t)}^{-1}(a,b)&=(L_t L_{x_1} L_{x_2}^{-1}L_t^{-1}(a),b)=(L_{L_t(x_1)} L_{L_t(x_2)}^{-1}(a),b),
 %\end{align*}
% for every $x_1,x_2,x,z,a\in Q_1$ $y_1,y_2,y,t,b\in Q_2$. Therefore $H_1,H_2$ are normal subgroups in $\dis(Q)$, they are isomorphic to $\dis(Q_1)$ and $\dis(Q_2)$ respectively and $H_1\cap H_2=1$. Since $$L_{(x_1,y_1)} L_{(x_2,y_2)}^{-1}(a,b)=(L_{x_1} L_{x_2}^{-1}(a),L_{y_1} L_{y_2}^{-1}(b))=L_{(x_1,y)} L_{(x_2,y)}^{-1}L_{(x,y_1)} L_{(x,y_2)}^{-1}(a,b)$$  for every $x_1,x_2,a\in Q_1$ $y_1,y_2,b\in Q_2$. we have that $\dis(Q)=H_1 H_2$, i.e. 
%So $\dis(Q)\cong H_1\times H_2\cong \dis(Q_1)\times \dis(Q_2)$.
%\end{proof}
\begin{corollary}
    Let $Q$ be a directly decomposable left quasigroup. Then $\dis(Q)$ is directly decomposable.
\end{corollary}


Note that if $\dis(Q)$ is directly decomposable, then $Q$ does not need to be directly decomposable, not even for quandles. In fact $Q={\tt SmallQuandle(32,5)}$ is directly indecomposable, but $\dis(Q)\cong \mathcal{Q}_8\times \Z_2^2$. 



\begin{lemma}
Let $Q$ be a connected faithful quandle. 
%\begin{itemize}
 %   \item[(i)] If $Q$ is faithful or connected and directly decomposable then $\dis(Q)$ is directly decomposable.
%\item[(ii)]  
If $\lmlt(Q)$ is directly decomposable, then $Q$ is directly decomposable.  
%\end{itemize}
\end{lemma}

    
\begin{proof}
%(i) Assume that $Q$ is directly decomposable, that is there exists a pair of congruence $\alpha,\beta\in \con(Q)$ such that $\alpha\wedge \beta= 0_Q$ and $\alpha\circ \beta=1_Q$. Therefore we have also that $\alpha\vee\beta=1_Q$. If $Q$ is faithful then $\dis_\alpha,\dis_\beta\neq 1$. If $Q$ is connected and $\dis_\alpha=1$ then $\dis_{\alpha\vee \beta}=\dis_\beta=\dis(Q)$ and so $\mathcal{O}_{\dis(Q)}=1_Q\leq \beta$. Therefore $\dis_\alpha,\dis_\beta\neq 1$. Thus, $\dis_\alpha\cap \dis_\beta\leq \dis^\alpha\cap \dis^\beta=1$ and $\dis_\alpha \dis_\beta=\dis^{\alpha\vee\beta}=\dis(Q)$ \cite[Proposition 3.2(2)]{CP}. Therefore $\dis(Q)\cong \dis_\alpha\times \dis_\beta$. 

%\comment{questo funziona anche senza faithful? per left quasigroups?}


%(ii) 
Assume that $\lmlt(Q)$ is directly decomposable, that is, there are normal subgroups $1 \neq N,M\trianglelefteq\lmlt(Q)$ such that $N\cap M=1$ and $\lmlt(Q)=NM$. Then $0_Q\neq \mathcal{O}_N\leq \c{N}$, $0_Q\neq \mathcal{O}_M\leq \c{M}$ and $\c{N}\wedge\c{M}=\c{N\cap M}=\c{1}= \lambda_Q=0_Q$. Moreover $\c{N}\circ \c{M}\geq \mathcal{O}_N\circ \mathcal{O}_M=\mathcal{O}_{NM}=\mathcal{O}_{\lmlt(Q)}=1_Q$. Then $Q$ is directly decomposable.
\end{proof}

For faithful left quasigroups, subdirect irreducibility can also be reduced to the same property for the left multiplication group.

\begin{lemma}\label{SI LQGs}
Let $Q$ be a faithful left quasigroup. 
\begin{itemize}
    \item[(i)] If $Q$ is subdirectly reducible, then $\lmlt(Q)$ is subdirectly reducible. 
\item[(ii)]  If $Q$ is a quandle and $\lmlt(Q)$ is subdirectly reducible, then $Q$ is subdirectly reducible.
\end{itemize}
\end{lemma}

%\comment{using $\dis_\alpha$ for left quasigroups}
\begin{proof}
(i) Assume that $Q$ is subdirectly reducible, that is, there exists a set $\setof{\alpha_i}{i\in I}\subseteq \con(Q)$ such that $\alpha_i\neq 0_Q$ for all $i\in I$ and $\beta=\bigwedge_{i\in I} \alpha_i=0_Q$. The left quasigroup $Q$ is faithful, so $\dis_{\alpha_i}\neq 1$ for every $i\in I$, so also the group $\dis^{\alpha_i}$ is not trivial for every $i\in I$. Then $\bigcap_{i\in I} \dis^{\alpha_i}=\dis^\beta=1$ and so $\lmlt(Q)$ is subdirectly reducible. 
%\end{proof}
%
%
%\begin{proposition}  
%Let $Q$ be a faithful quandle. Then $Q$ is subdirectly reducible if and only if $\lmlt(Q)$ is subdirectly reducible.
%\end{proposition}

%\begin{proof}

(ii) Assume that $\lmlt(Q)$ is subdirectly reducible, that is, there exist $N,M\neq 1$ be normal subgroups of $\lmlt(Q)$ such that $N\cap M=1$. Then $0_Q\neq \mathcal{O}_N\leq \c{N}$, $0_Q\neq \mathcal{O}_M\leq \c{M}$ and $\c{N}\wedge\c{M}=\c{N\cap M}=\c{1}= \lambda_Q=0_Q$. Then $Q$ is subdirectly reducible.
\end{proof}


%\begin{lemma}\label{DD LQGs}
%Let $Q_1,Q_2$ be quandles. Then $\dis(Q_1\times Q_2)\cong \dis(Q_1)\times \dis(Q_2)$.
%\end{lemma}
%\begin{proof}
 %   Let $Q=Q_1\times Q_2$, we define:
%\begin{align*}
 %   H_1&=\langle L_{(x_1,y)} L_{(x_2,y)}^{-1},\,x_1,x_2\in Q_1, \, y\in Q_2\rangle,\\
%    H_2&=\langle L_{(x,y_1)} L_{(x,y_2)}^{-1},\,y_1,y_2\in Q_2, \, x\in Q_1\rangle.
%\end{align*}
%  It is easy to check that $H_1$ and $H_2$ are trivially intersecting normal subgroups of $\dis(Q)$ generating $\dis(Q)$ and such that $\dis(Q_i)\cong H_i$ for $i=1,2$.
% \begin{align*} 
  %   L_{(x_1,y)} L_{(x_2,y)}^{-1}(a,b)&=(L_{x_1} L_{x_2}^{-1}(a),b),\\
  %   L_{(x,y_1)} L_{(x,y_2)}^{-1}(a,b)&=(a,L_{y_1} L_{y_2}^{-1}(b)),\\
  %   L_{(z,t)} L_{(x,y_1)} L_{(x,y_2)}^{-1}L_{(z,t)}^{-1}(a,b)&=(a,L_t L_{y_1} L_{y_2}^{-1}L_t^{-1}(b))=(a, L_{L_t(y_1)} L_{L_t(y_2)}^{-1}(b)),\\
 %    L_{(z,t)} L_{(x_1,y)} L_{(x_2,y)}^{-1}L_{(z,t)}^{-1}(a,b)&=(L_t L_{x_1} L_{x_2}^{-1}L_t^{-1}(a),b)=(L_{L_t(x_1)} L_{L_t(x_2)}^{-1}(a),b),
 %\end{align*}
% for every $x_1,x_2,x,z,a\in Q_1$ $y_1,y_2,y,t,b\in Q_2$. Therefore $H_1,H_2$ are normal subgroups in $\dis(Q)$, they are isomorphic to $\dis(Q_1)$ and $\dis(Q_2)$ respectively and $H_1\cap H_2=1$. Since $$L_{(x_1,y_1)} L_{(x_2,y_2)}^{-1}(a,b)=(L_{x_1} L_{x_2}^{-1}(a),L_{y_1} L_{y_2}^{-1}(b))=L_{(x_1,y)} L_{(x_2,y)}^{-1}L_{(x,y_1)} L_{(x,y_2)}^{-1}(a,b)$$  for every $x_1,x_2,a\in Q_1$ $y_1,y_2,b\in Q_2$. we have that $\dis(Q)=H_1 H_2$, i.e. 
%So $\dis(Q)\cong H_1\times H_2\cong \dis(Q_1)\times \dis(Q_2)$.
%\end{proof}



\subsection{Commutator theory for quandles}




Commutator theory for quandles can be understood in purely group-theoretical terms by looking at the properties of the displacement group and its subgroups. The main results are the following, and they will be freely used throughout the paper. Let $Q$ be a quandle:
\begin{itemize}
    \item 
    if $Q$ is faithful and $\alpha\in \con(Q)$ then $\alpha$ is abelian (resp. central) if and only if $\dis_\alpha$ is abelian (resp. central in $\dis(Q)$) \cite[Corollary 5.4]{CP}.

\item The quandle $Q$ is solvable (resp. nilpotent) if and only if $\dis(Q)$ is solvable (resp. nilpotent) \cite[Theorem 1.2]{CP}.

%\begin{proposition}[\cite{CP}, Corollary 5.4]
 %   Let $Q$ be a faithful quandle and $\alpha\in \con(Q)$. Then $\alpha$ is abelian (resp. central) if and only if $\dis_\alpha$ is abelian (resp. central in $\dis(Q)$).
%\end{proposition}

%\begin{proposition}[\cite{CP}, Theorem 1.2]
 %   Let $Q$ be a quandle. Then $Q$ is solvable (resp. nilpotent) if and only if $\dis(Q)$ is solvable (resp. nilpotent).
%\end{proposition}

\item Affine quandles are abelian. On the other hand, connected abelian quandles are affine \cite[Corollary 3.4]{Medial}.


%\begin{corollary}[\cite{Medial}, Proposition 3.3]
 %   Connected abelian quandles are affine.
%\end{corollary}

\item If $Q$ is finite and connected, then $\gamma_Q=\mathcal{O}_{\gamma_1(\dis(Q))}$ and $\dis^{\gamma_Q}=\gamma_1(\dis(Q))$ \cite[Proposition 1.7]{GB}.

%\begin{proposition}[\cite{GB}, Proposition 1.6] \label{gamma for quandles}
 %   Let $Q$ be a finite connected quandle and $G=\dis(Q)$. Then $\gamma_Q=\mathcal{O}_{\gamma_1(G)}$ and $\dis^{\gamma_Q}=\gamma_1(G)$.
%
%\item If $Q$ is latin then $\zeta_Q=\mathcal{O}_{Z(G)}$ and $\dis_{\zeta_Q}=Z(\dis(Q))$ \comment{add ref - explain better} \cite[Corollary 5.10]{CP}.
 
%\begin{proposition}[\cite{CP}, Corollary 5.10]
  %  Let $Q$ be a finite connected faithful quandle and $G=\dis(Q)$. Then $\zeta_Q=\mathcal{O}_{Z(G)}$.
%\end{proposition}



%\subsection{Abelian congruences and solvable quandles}
 

\end{itemize}

For finite latin quandles, we have the following characterization of the center and its relative displacement group.
\begin{lemma}\label{center for latin}
    Let $Q$ be a finite latin quandle. Then $\zeta_Q=\mathcal{O}_{Z(\dis(Q))}$ and $\dis_{\zeta_Q}=Z(\dis(Q))$.
\end{lemma}

\begin{proof}
Recall that for finite latin quandles we have that the operators $\mathcal{O}_*$ and $\c{*}$ coincide. According to \cite[Corollary 5.10]{CP} we have $\zeta_Q=\c{Z(\dis(Q))}$ and so it coincides with $\mathcal{O}_{Z(\dis(Q))}=\mathcal{O}_{\dis_{\zeta_Q}}$ and $\dis_{\zeta_Q}\leq Z(\dis(Q))$ accordingly. The subgroup $Z(\dis(Q))$ is a central subgroup of a transitive group, and so $Z(\dis(Q))_x=1$ for every $x\in Q$. Therefore, $Z(\dis(Q))=\dis_{\zeta_Q}Z(\dis(Q))_x=\dis_{\zeta_Q}$.
\end{proof}



The blocks of congruences of quandles are subquandles, and the blocks of abelian congruences are abelian quandles (this follows directly from the abstract definition of abelianness of a congruence as defined in \cite{comm}). We can use the description of abelian quandles and the close relationship between the properties of congruences and the properties of subgroups of the displacement group to understand the features of abelian congruences (or more generally congruences with abelian blocks) to build some tools to study solvable quandles.

Under some further assumptions, the blocks of abelian congruences are affine, so we have the following description of their structure.



\begin{lemma}\label{below abelian}
    Let $Q$ be a finite latin quandle, and let $\alpha$ be an abelian congruence. Then:
    \begin{itemize}
        \item[(i)] The block $[a]_\alpha$ is affine, i.e $[a]_\alpha\cong \aff(A,f)$ for some abelian group $A$ and some $f\in \aut{A}$. 
        \item[ (ii)] If $\beta\leq\alpha$ then $[0]_\beta$ is an $f$-invariant subgroup of $A$. In particular, $|\setof{\beta\in \con(Q)}{\beta\leq \alpha}|\leq |\setof{N\leq A}{f(N)=N}|$.
    \end{itemize}
    

\end{lemma}

%\comment{Using that the blocks of abelian congruences are abelian quandles and that latin abelian quandles are af}
\begin{proof}
%latin quandles are congruence regular, i.e. if $[x]_\beta=[x]_\delta$ then $\beta=\delta$.
%
(i) The blocks of $\alpha$ are abelian latin quandles. So, they are affine, that is, $[x]_\alpha\cong \aff(A,f)$ for some abelian group $A$ and some $f\in \aut{A}$. 

(ii) Given $\beta\leq \alpha$ then $[0]_\beta$ is a subquandle of $\aff(A,f)$ and so it is an $f$-invariant subgroup of $A$ \cite[Proposition 3.14]{Principal}. Connected quandles are congruence regular (see \cite[Lemma 2.4]{Maltsev_paper}), so the map $\beta \mapsto [0]_\beta $ from $\setof{\beta\in \con(Q)}{\beta\leq \alpha}$ to $\setof{N\leq A}{f(N)=N}$ is injective.
%On the other hand, according to... given $N\leq A$ such that $f(N)=N$ then $\sim_N$ defined as $x\,\sim_N\, y$ if and only if $[x]_\alpha=[y]_\alpha$ and $x-y\in N$.
\end{proof}



%\begin{lemma}\label{N cyclic 2}
 %   Let $Q$ be a finite quandle, $\alpha\in \con(Q)$ and $N\leq \dis^\alpha$. If $[x]_\alpha$ are simple latin affine quandles and $\gcd\{|N|,|[x]_\alpha|\}=1$ then $N$ is cyclic and $N$ embeds into $\aut{[x]}_x$ for every $x\in Q$.
%\end{lemma}
 %   \comment{fare un lemma con il primo paragrafo?}
 

%\begin{proof}
%The blocks of $\alpha$ are strictly simple, so $[x]\cong \aff(\Z_p^n,f)$ where $f$ acts irreducibly, $\aut{[x]}=\Z_p^n\rtimes \Z_{p^n-1}$, and $\aut{[x]}_x= C_{GL_n(p)}(f)$ \comment{add ref}. If $g\in \aut{[x]}$ has order relatively prime with $p$ then $g(y)=t+h(y)$ for $1\neq h\in C_{GL_n(p)}(f)$. Assume that $0\neq y\in Fix(h)$. Then $f(y)=fh(y)=h(f(y))$, so $Fix(h)$ is $f$-invariant, therefore $Fix(h)=0$ and so $1-h$ is invertible. Hence, there exists $z\in [x]$ such that $(1-h)(z)=t$, i.e. $g(z)=z$.

 %Let $h,h'\in N$. Then $h|_{[x]}\in\aut{[x]}$ has order relatively prime with $|[x]|$ and so for every $[x]\in Q/\alpha$ there exists $c\in [x]$ such that $h(c)=c$. Let $X$ be a set of representatives of the blocks of $\alpha$. Then $Q$ is generated by $X\cup [x]_\alpha$ for every $x\in Q$.
% 
 %Moreover $N$ embeds into a power of $\aut{[x]}_{c}$ that is abelian. So $N$ is abelian, and then $h h'(c)=h'h(c)=h'(c)$. Thus $h\in \aut{[x]}_{h'(c)}\cap \aut{[x]}_c=\aut{[x]}_{c,h'(c)}$. Since the quandle $[x]$ is generated by any pair of elements, then $\aut{[x]}_{a,b}=1$ whenever $a\neq b$. Thus, $h'(c)=c$. If $N|_{[x]}=1$ for every $[x]\in Q/\alpha$ then $N=1$. Otherwise, $N|_{[x]}\leq \aut{[x]}_c$ for some $x\in Q$. 
 %Let $\phi:h\mapsto h|_{[x]}$ mapping $N$ to $\aut{[x]}$. If $\phi(h)=1$ then $h$ fixes all the elements of the block $[x]$ and an element for every other block of $\alpha$. Then $h$ fixes a set of generators of $Q$ and so $N$ embeds into $\aut{[x]}$.

%he order of $N$ and of $\dis([x])$ are relatively prime, so $[\phi(N),\phi(N)]\leq N\cap \dis([x])=1$, so $N$ is abelian. So if $h,h'\in N$ and $h(x)=x$ then $h'(x)=h'h(x)=h(h'(x))$. So $h$ fixes $x$ and $h'(x)$. So $h$ fixes $Sg(x,h'(x))$. If $h'(x)\neq x$ then $S=[x]$ and $h=1$ accordingly. Therefore $h'(x)=x$, so we have that $N$ embeds into $\aut{([x])}_x$ that is cyclic. 
%\end{proof}



%
%\subsection*{8p}
%
%We know that there are no latin quandles of size $2$ and $2p$ for $p$ prime and that simple latin quandles have size a power of a prime.
%
%\begin{lemma}
%Let $Q$ be a latin quandle of size $8p$ then it has no factor of size $2,4,2p,4p$.
%\end{lemma}
%\begin{proof}
%The blocks have size $2p$.
%\end{proof}
%
%\begin{lemma}
%Let $Q$ be a subdirectly reducible latin quandle of size $8p$. Then $Q$ is a direct product of a quandle of size $8$ and a quandle of size $p$. In particular $Q$ is affine.
%\end{lemma}
%
%\begin{proof}
%Let $\alpha\wedge\beta=0_Q$. Then $Q$ embeds into $Q/\alpha\times Q/\beta$ and $|Q|$ divides $|Q/\alpha| |Q/\beta|$. According to the Lemma above the factors have size $8$ and $p$ respectively and they are simple (all latin quandles of size $8$ are simple. Hence $Q=Q/\alpha\times Q/\beta$ are $Q/\alpha$.
%\end{proof}
%
%
%\begin{theorem}
%There are no non-affine quandles of size $8p$ where $p$ is a prime.
%%Let $Q$ be a latin non-affine quandle of size $8p$, $p>5$. Then $Q\in \LSS(8,p,\gamma_Q)$ and $p=7,31$.
%\end{theorem}
%
%\begin{proof}
%If $p\leq 5$ a computer search on the RIG library of GAP reveals that there are no such quandles.\\ 
%Let $p>5$. The quandle $Q$ has no factor of size $2$, and it has also no factor of size $4$, since in such case the blocks of the congruence have size $2p$, but there is no connected quandle of size $2p$ for $p>5$. Simple factors of $Q$ have prime power size, so they have either size $8$ or $p$. If $Q$ has two different maximal congruence then they have either  it has a factor of size either $64$, $p^2$ or $8p$. Since the size of the factors divides the size of $Q$, the first two cases can not happen and in the last case $Q$ is a direct product of simple latin quandle and then affine. Let $\alpha$ be the unique maximal congruence of $Q$. If the blocks of $\alpha$ have size $p$ then it is also minimal, if the blocks have size $8$, since all the latin quandles of size $8$ are strictly simple then $\alpha$ is minimal aswell. Hence $Q$ has a unique proper congruence, which is $\gamma_Q$ since $Q$ is not affine and the factor is affine, and then $Q\in \LSS(p,8,\sigma_Q)$ or $Q\in \LSS(8,p,\sigma_Q)$. According to \comment{my results on $\LSS$} if $Q\in \LSS(p,2^n,\sigma_Q)$ then $n$ is even. Hence $Q\in \LSS(8,p,\sigma_Q)$ and therefore $\sigma_Q=\gamma_Q$ \comment{add ref}. \\
%In particular, $\dis_\alpha=\dis^\alpha=\gamma_1(\dis(Q)$ embeds into $\dis([a])^2\cong \mathbb{Z}_2^6$. Therefore $\dis(Q)\cong \dis^\alpha\rtimes \mathbb{Z}_p$ and in particular $p$ divides the size of $\aut{\dis^\alpha}=GL_{3+k}(2)$ for some $0\leq k\leq 3$. Hence $p\in \{7,31\}$. An exhaustive GAP computer search on the groups of the form $\mathbb{Z}_2^{3+k}\rtimes \mathbb{Z}_p$ and their automorphisms for $0\leq k\leq 3$ and $p\in \{7,31\}$ shows that there are no such quandles.
%\end{proof}

We are particularly interested to understand minimal congruences of quandles, in particular for solvable ones (e.g. finite latin quandles, see Corollary 1.3 of \cite{CP}).

\begin{lemma}\label{minimal cong}
Let $Q$ be a finite faithful quandle, and $\alpha\in \ma_Q$. 
\begin{itemize}
    \item[(i)] $\alpha=\mathcal{O}_{\dis_\alpha}$ and $\dis_\alpha$ is a minimal normal subgroup of $\lmlt(Q)$.
    \item[(ii)] If $\alpha$ is abelian, then $\dis_\alpha$ is an elementary abelian $p$-group for some prime $p$ and the blocks of $\alpha$ have size a power of $p$.
\item[(iii)]  If $Q$ is solvable, then $\alpha$ is abelian. In particular, $\nu_Q$ is abelian.
\end{itemize}

\end{lemma}
%\comment{Does it follow that the blocks of $\alpha$ are simple? Check on GAP}
\begin{proof}

(i) Let $\beta=\mathcal{O}_{\dis_\alpha}\leq \alpha$. So, $\beta=0_Q$ or $\alpha=\beta$. If $\beta=0_Q$ then $\dis_\alpha=1$. Thus, $\alpha\leq \lambda_Q=0_Q$, contradiction.

Assume that $N\leq \dis_\alpha$. Then $\delta=\mathcal{O}_N\leq \alpha$. Hence $\delta=0_Q$ and accordingly $N=1$ or $\delta=\alpha$ and so $\dis_\alpha\leq N$ (see \cite{CP}, Section 3.2). 

(ii) According to \cite[Corollary 5.4]{CP} the subgroup $\dis_\alpha$ is abelian. Let $H$ be a characteristic subgroup of $\dis_\alpha$. Then $H$ is normal in $\lmlt(Q)$ and therefore trivial since $\dis_\alpha$ is a minimal normal subgroup of $\lmlt(Q)$. Hence $\dis_\alpha$ is necessarily an elementary abelian $p$-group for some prime $p$ (the $p$-Sylow subgroup and the subgroups of the form $p^n\dis_\alpha$ for $n\in \mathbb{N}$ are characteristic). So, $|[x]|=[\dis_\alpha:(\dis_\alpha)_x]$ is a power of $p$.

(iii) The group $\dis(Q)$ is solvable, since $Q$ is solvable. Let $\alpha$ be a minimal congruence. Then $\dis_\alpha$ is a solvable minimal normal group. Therefore, $[\dis_\alpha,\dis_\alpha]<\dis_\alpha$, and thus $[\dis_\alpha,\dis_\alpha]=1$. Hence $\dis_\alpha$ is abelian and so is $\alpha$. According to \cite[Proposition 3.2(2)]{CP} we have 
$$\dis_{\nu_Q}=\dis_{\vee \{\alpha\in \ma_Q\}}=\langle\dis_\alpha,\,\alpha \in \ma_Q\rangle.$$
Moreover, $[\dis_\alpha,\dis_\beta]\leq \dis_\alpha\cap\dis_\beta=1$ and the subgroup $\dis_\alpha$ is abelian for every $\alpha,\beta\in\ma_Q$. Hence $\dis_{\nu_Q}$ is abelian and so $\nu_Q$ is also abelian. %\comment{maybe I have just an embedding!}
\end{proof}

%\begin{corollary}\label{nu is abelian}
%Let $Q$ be a finite solvable faithful quandle. Then $\nu_Q$ is abelian.
%\end{corollary}

%\comment{maybe this is true for Malt'sev algebras in general}

%\begin{proof}
%The minimal congruences are abelian. Then, according to \cite[Proposition 3.2(2)]{CP} we have 
%$$\dis_{\nu_Q}=\prod_{\alpha\in \ma_Q}\dis_\alpha.$$
%The subgroups $\dis_\alpha$ are abelian and so it is $\dis_{\nu_Q}$. %\comment{maybe I have just an embedding!}
%\end{proof}
		


%\begin{lemma}\comment{add ref} Finite solvable superconnected quandles have the Lagrange property.
%\end{lemma}




\begin{lemma}\label{dis_gamma}
Let $Q$ be a non-nilpotent quandle. Then: 
\begin{itemize}
  \item[(i)] If $Q$ is faithful and $\gamma_Q\in \ma_Q$, then $\dis_{\gamma_Q}=\gamma_2(\dis(Q))$. 
    \item[(i)] If $\ma_Q=\{\gamma_Q\}$, then $Z(\dis(Q))=1$.
  
\end{itemize}
\end{lemma}

\begin{proof}
Let $G=\dis(Q)$. 


(i) According to \cite[Proposition 3.3(1)]{CP} we have $\gamma_2(G)=[G,\gamma_1(G)]\leq \dis_{\gamma_Q}\neq 1$ and Lemma \ref{minimal cong}(i) implies that the subgroup $\dis_{\gamma_Q}$ is a minimal normal subgroup of $\lmlt(Q)$. So, $\gamma_2(G)=1$ or $\gamma_2(G)=\dis_{\gamma_Q}$. In the first case, $G$ is nilpotent and so also $Q$ is nilpotent by \cite[Theorem 1.2(2)]{CP}, contradiction. Then $\dis_{\gamma_Q}=\gamma_2(G)$.

(ii) Assume that $Z(G)\neq 1$. Then $\beta=\mathcal{O}_{Z(G)}\neq 0_Q$. According to \cite[Lemma 5.12]{CP} $\beta$ is central and $\gamma_Q\leq \beta$ since $\gamma_Q$ is the only minimal congruence of $Q$. Thus, $\gamma_2(Q)=[\gamma_Q,1_Q]=0_Q$ and so $Q$ is nilpotent, which is a contradiction.
\end{proof}

%\subsection{Central congruences}
%We are going to use a well-known fact about connected quandles: if $Q$ is connected then $\dis(Q)=[\lmlt(Q),\lmlt(Q)]$ (\cite[Proposition 2.3]{hsv}).

We now introduce a corollary of Lemma 1.4 of \cite{LSS}.
\begin{corollary}
    \label{embedding as quandle}
	Let $Q$ be a connected quandle, $\alpha\in \con(Q)$ such that $Q/\alpha=Sg([a_1],\ldots, [a_n])$. Then:
	\begin{itemize}
		\item[(i)] the mapping
		\begin{displaymath}\psi: \dis^\alpha\longrightarrow \prod_{i =1}^n \aut{[a_i]},\quad h\mapsto (h|_{[a_1]},\ldots, h|_{[a_n]}),
		\end{displaymath}
		is an embedding of groups and $\dis^\alpha|_{[a_i]_\alpha}\cong \dis^\alpha|_{[a_j]_\alpha}$ for every $1\leq i,j\leq n$.
		
		\item[(ii)] If $Q$ is a finite latin quandle and $\alpha$ is abelian, then $\dis_\alpha|_{[a]}\cong \dis([a])$ and $\dis_\alpha$ embeds in $\dis([a])^n$. In particular $|\dis_\alpha|$ divides $|[a]|^n$. 
		%In particular $\dis(Q)_{a_1}$ embeds into $\prod_{i =2}^n \aut{[a_i]}$.
	\end{itemize}

\end{corollary}

 \begin{proof}
	(i) Apply Lemma 1.4(i) of \cite{LSS} to $\dis^\alpha$.
	
	(ii) The blocks of $\alpha$ are finite latin affine quandles. The group $\dis_\alpha$ is abelian and so $\dis_\alpha|_{[x]}$ is regular on $[x]$. Then Lemma 1.4(ii) of \cite{LSS} applies to $\dis_\alpha$. So $\dis_\alpha$ embeds in $\dis([a])^n$ that has size $|[a]|^n$.  
\end{proof}

In the following, we are displaying some sufficient conditions for a congruence of a quandle to be central. 
\begin{lemma}\label{central 2}
    Let $Q$ be a connected faithful quandle and $\alpha\in \con(Q)$. If $\dis_\alpha$ is cyclic, then $\alpha$ is central.
\end{lemma}

\begin{proof}
According to \cite[Proposition 2.3]{hsv} we have $\dis(Q)=[\lmlt(Q),\lmlt(Q)]$. Assume that $H=\dis_\alpha$ is cyclic. The subgroup $H$ is normal in $\lmlt(Q)$, so we have a group action $\rho:\lmlt(Q)\longrightarrow \aut{H}$ where $\rho_x=\widehat{x}|_H$ for every $x\in \lmlt(Q)$. The group $\aut{H}$ is abelian, so $\dis(Q)=[\lmlt(Q),\lmlt(Q)]\leq \ker{\rho}$, i.e. $H$ is central in $\dis(Q)$. According to \cite[Corollary 5.4]{CP}, we see that $\alpha$ is a central congruence.
\end{proof}
Note that Lemma \ref{central 2} only uses that the automorphism group of $\dis_\alpha$ is abelian, so if this condition holds, then $\alpha$ is central.
%
%\begin{corollary}
 %   Let $Q$ be a finite principal latin quandle and $\alpha$ be a congruence of $Q$. If the blocks of $\alpha$ have prime order then $\alpha$ is central.
%\end{corollary}

%\begin{proof}
 %   The blocks of $\alpha$ have size equal to the size of $\dis_\alpha$. Then if $|[x]|=|\dis_\alpha|$ is prime then $\dis_\alpha$ is cyclic.
%\end{proof}

%
%\begin{lemma}\label{min/max cong} \comment{NEEDED?}
%Let $Q$ be a quandle, $\alpha\in \Ma_Q$ and $\beta\in \ma_Q$. Then either $Q$ embeds into $Q/\alpha\times Q/\beta$ or $\beta\leq \alpha$.
%\end{lemma}
%
%\begin{proof}
%It follows since $\alpha\wedge\beta\leq \beta$ and $\beta$ is minimal.
%\end{proof}


\begin{lemma}\label{central}
Let $Q$ be a finite connected faithful quandle and let $\alpha\in \con(Q)$ such that $\alpha\wedge \gamma_Q=0_Q$. Then: 
\begin{itemize}
\item[(i)] $\alpha$ is central; 
\item[(ii)] $\dis_\alpha=\dis^\alpha$; 
\item[(iii)] $Q/\alpha$ is faithful.
\end{itemize}

\end{lemma}

\begin{proof}
(i) Let $A$ be an algebra, and $\alpha$ be a congruence. Then $[1_A,\alpha]\leq \alpha\wedge \gamma_A=0_A$ \cite[Proposition 3.4]{comm}. Therefore, $\alpha$ is central.

(ii) According to \cite[Proposition 3.3(i)]{CP} we have
$$[\dis(Q),\dis^\alpha]\leq \gamma_1(\dis(Q))\cap \dis_\alpha=\dis^{\gamma_Q}\cap \dis_\alpha\leq \dis^{\gamma_Q}\cap \dis^\alpha=\dis^{\gamma_Q\wedge \alpha}=\dis^{0_Q}=1.$$
Then $\dis^\alpha\leq Z(\dis(Q))$. According to \cite[Lemma 4.4(i)]{Nilpotent}, $\dis^\alpha=\dis_\alpha\times \left(\dis^\alpha\right)_x$ and since $\dis^\alpha$ is a central subgroup of the transitive group $\dis(Q)$ then $\left(\dis^\alpha\right)_x=1$. Consequently $\dis_\alpha=\dis^\alpha$. 

(iii) Using \cite[Corollary 4.5(i)]{Nilpotent}, we have $\alpha=\c{\dis_\alpha}=\c{\dis^\alpha}$. Then we can apply \cite[Proposition 3.7]{CP} and we find that $Q/\alpha$ is faithful.
\end{proof}


\section{latin quandles of size \texorpdfstring{$16p$}{16p}} \label{sec:latin16p}

%\subsection{$\LSS$}

\subsection{Locally strictly simple quandles}
A quandle is {\it strictly simple} if it does not have proper subquandles. Finite strictly simple quandles coincide with simple latin quandles \cite{Principal}. Latin quandles are solvable (see \cite{CP}, Corollary 1.3), thus finite strictly simple quandles are abelian and hence affine. Recall that finite simple abelian quandles are given by affine quandles of the form $Q=\Aff(\Z_p^n,f)$ where $\Z_p^n$ does not have $f$-invariant subgroups (see \cite{Principal} and \cite{Medial} for more details on abelian quandles). In particular, $\aut{Q}\cong \Z_p^n\rtimes \Z_{p^n-1}$, $\aut{[x]}_x= C_{GL_n(p)}(f)\cong \Z_{p^n-1}$ and $\aut{Q}/\dis(Q)$ is abelian \cite[Proposition 3.8]{Principal}. %Moreover, $Q$ is generated by any pair of elements and accordingly $\aut{Q}_{x}\cap\aut{Q}_{y}=1$ for every $x\neq y\in Q$. 

The following result has already been implicitly stated in the proof of \cite[Theorem 3.14]{LSS} for which we include a cleaner proof divided into two parts. 
\begin{lemma}\label{aut of simple}
    Let $Q=\aff(\Z_p^n, f)$ be a simple quandle, and $g\in \aut{Q}$. If $|g|$ and $p$ are relatively prime, then $g\in \aut{Q}_x$ for some $x\in Q$. 
\end{lemma}


\begin{proof}
   %The automorphism $f$ acts irreducibly, $\aut{[x]}=\Z_p^n\rtimes \Z_{p^n-1}$ and $\aut{[x]}_x= C_{GL_n(p)}(f)\cong \Z_{p^n-1}$ \cite[Proposition 3.8]{Principal}. 
   If $g\in \aut{[x]}$ has order relatively prime with $p$ then $g(y)=t+h(y)$ for $1\neq h\in C_{GL_n(p)}(f)$ and some $t\in \Z_p^n$ for every $y\in Q$. Assume that $0\neq y\in Fix(h)$. Then $f(y)=fh(y)=h(f(y))$, so $Fix(h)$ is an $f$-invariant subgroup. % and so it is a subquandle of $Q$ containing $0$ (see \cite[Proposition 3.14]{Principal}). 
   Therefore, $Fix(h)=0$ and so $1-h$ is invertible. Hence, there exists $z\in Q$ such that $(1-h)(z)=t$, i.e. $g(z)=z$.
\end{proof}

\begin{lemma}\label{N cyclic}
    Let $Q$ be a finite quandle, $\alpha\in \con(Q)$ and $N\leq \dis^\alpha$. If $[x]_\alpha$ is a simple affine quandle of order relatively prime with $|N|$ then $N$ embeds into $\aut{[x]}_x$ for some $x\in Q$. In particular $N$ is cyclic.
\end{lemma}
  
 

\begin{proof}
 Let $g\in N$. Then $g|_{[x]}\in\aut{[x]}$ has order relatively prime with $|[x]|$ and so for every $[x]\in Q/\alpha$ there exists $c\in [x]$ such that $g(c)=c$ (see Lemma \ref{aut of simple}). Let $X$ be any set of representatives of the blocks of $\alpha$. Then $Q$ is generated by $X\cup [x]_\alpha$ for every $x\in Q$ \cite[Lemma 6.1]{GB}.
% 
 %Moreover $N$ embeds into a power of $\aut{[x]}_{c}$ that is abelian. So $N$ is abelian, and then $h h'(c)=h'h(c)=h'(c)$. Thus $h\in \aut{[x]}_{h'(c)}\cap \aut{[x]}_c=\aut{[x]}_{c,h'(c)}$. Since the quandle $[x]$ is generated by any pair of elements, then $\aut{[x]}_{a,b}=1$ whenever $a\neq b$. Thus, $h'(c)=c$. If $N|_{[x]}=1$ for every $[x]\in Q/\alpha$ then $N=1$. Otherwise, $N|_{[x]}\leq \aut{[x]}_c$ for some $x\in Q$. 
 Let $\phi:g\mapsto g|_{[x]}$ map $N$ to $\aut{[x]}$. If $\phi(g)=1$ then $g$ fixes all the elements of block $[x]$ and an element for every other block of $\alpha$. Then $g$ fixes a set of generators of $Q$ and so $g=1$. Hence $N$ embeds into $\aut{[x]}$.

Since $\aut{[x]}/\dis([x])$ is abelian, then $[\aut{[x]},\aut{[x]}]\leq \dis([x])$ and so is also $[\phi(N),\phi(N)]=\phi([N,N])\leq \dis([x])$. The order of $N$ and of $\dis([x])$ are relatively prime (indeed $|\dis([x])|=|[x]|$), so $[\phi(N),\phi(N)]=1$ and thus $N$ is abelian. So, if $g,g'\in N$ and $g(x)=x$ then $g'(x)=g'(g(x))=g(g'(x))$. So $g$ fixes $S=Sg(x,g'(x))$. If $g'(x)\neq x$ then $S=[x]$ and $g=1$ accordingly. Therefore, $g'(x)=x$, so we have that $N$ embeds in $\aut{[x]}_x$ which is cyclic. So $N$ is also cyclic. 
\end{proof}

We now consider the class of {\it locally strictly simple} (LSS) quandles, that is, the class of quandles such that all proper subquandles are strictly simple. In \cite{LSS} we proved that the non-affine subdirectly irreducible LSS quandles have a unique proper congruence with strictly simple factor and strictly simple blocks (see \cite{LSS}). Let $\sigma_Q$ be the equivalence relation defined by setting
$$x\, \sigma_Q\, y \text{ if and only if }\dis(Q)_x=\dis(Q)_y.$$ 
We denote by $\LSS(p^m,q^n,\alpha)$ the class of subdirectly irreducible non-affine LSS quandles such that the unique factor has size $q^n$, the blocks of the unique congruence have size $p^m$ and $\sigma_Q=\alpha$.

We can use \cite[Proposition 5.7]{Principal} (also mentioned in \cite[Proposition. 3.11]{LSS}) to describe the structure of the displacement group of quandles with the congruence lattice given by the three-element chain that includes all LSS latin quandles. 

\begin{lemma}\label{dis for 3-chains}
    Let $Q$ be a finite connected faithful solvable quandle with congruence lattice given by a chain of length $3$. If $Q$ is not nilpotent then one of the following holds: 
    \begin{itemize}
        \item[(i)] $\gamma_1(\dis(Q))=\gamma_2(\dis(Q))$ and $\dis(Q)=\Z_p^k\rtimes_\rho \Z_q^m$ where $p,q$ are primes and $k\in \mathbb{N}$.  
        \item[(ii)] $\gamma_1(\dis(Q))\neq\gamma_2(\dis(Q))$ and $\dis(Q)=\Z_p^k\rtimes_\rho K$ where $p$ is a prime, $k\in \mathbb{N}$ and $K$ is a special $q$-group. 
    \end{itemize}
In both cases, $Z(\dis(Q))=1$ and $\rho$ is faithful.
\end{lemma}

\begin{proof}
Let $G=\dis(Q)$. Since $Q$ is solvable but not nilpotent, then the unique non-trivial congruence of $Q$ is $\gamma_Q$, it is abelian and the blocks have size a power of a prime (see Lemma \ref{minimal cong}). According to Lemma \ref{dis_gamma}(ii) we have $\dis_{\gamma_Q}=\gamma_2(G)$ and it is an elementary abelian $p$-group according to Lemma \ref{minimal cong}(ii). The quandle $Q/\gamma_Q$ is simple and abelian, therefore $\dis(Q/\gamma_Q)=G/\gamma_1(G)$ is elementary abelian. So $Q/\gamma_Q\cong \aff(\Z_q^n,f)$ where $f$ has no invariant subgroups. The quandle $Q$ is not nilpotent, so $p\neq q$ (connected quandles of size a power of a prime are nilpotent \cite[Theorem 1.4]{CP}).

Let $\gamma_1(G)\leq H\neq G$ where $H$ is a characteristic subgroup of $G$. Then $H/\gamma_1(G)$ is an $f_{\gamma_1(G)}$-invariant subgroup of $G/\gamma_1(G)$, so $H=\gamma_1(G)$ and $\gamma_1(G)$ is a maximal characteristic subgroup of $G$. Thus, the group $G$ satisfies the hypothesis of \cite[Proposition 5.7]{Principal} and the statement follows. 

According to Lemma \ref{dis_gamma}(i), $\dis(Q)$ has trivial center, and therefore $\rho$ is faithful according to Lemma \ref{no center}. 
\end{proof}



For the scope of this paper, we need to investigate the classes $\LSS(p,16,\sigma_Q)$ and $\LSS(16,p,\sigma_Q)$ as they are classes of quandles of size $16p$.
%According to \cite[Theorem 3.14]{LSS}, if $Q\in \LSS(16,p,\sigma_Q)$ then $\dis(Q)=\Z_2^{4+k}\rtimes_\rho \Z_p$ and accordingly $p\in \{3,5,7,17,31,127\}$ \comment{are they latin?}. 

\begin{proposition}\label{no LSS}
    The class $\LSS(p,16,\alpha)$ is empty for every $\alpha$.
\end{proposition}


\begin{proof}
    %According to \cite[Proposition 3.18]{LSS} $\sigma_Q=0_Q$. 
  Let $Q\in \LSS(p,16,\alpha)$ and $G=\dis(Q)$.  According to the results of Section 3 of \cite{LSS} and following the argument of the proof of Proposition 3.21 of \cite{LSS} it follows that $G=\Z_p^2\rtimes_\rho K$ where $K$ is an extraspecial $2$-group of size $32$, $\rho$ is faithful, $\gamma_1(G)= \Z_p^2\rtimes_\rho Z(K)$ and $Z(\gamma_1(G))=1$. There are only $2$ extraspecial $2$-groups of size $32$, i.e. $K_i=$ \texttt{SmallGroup(32,i)} for $i=49,50$. If $K=K_{49}$, then it has a subgroup isomorphic to $\Z_2^3$. Thus, $\rho$ is not faithful according to Corollary \ref{ro faithful on Z_p}. If $K=K_{50}$, then $\gamma_1(G)$ has non-trivial center, according to Proposition \ref{no rep}. So, there are no quandles in the class $\LSS(p,16,\sigma_Q)$.
    %The quandle $Q_{\gamma_Q}$ as defined in ... is a connected quandle of size $32$ with a unique proper congruence. The connected quandles of size $32$ are classified, so there is a unique quandle matching such properties, namely SmallQuandle(32,4) and $K=\dis(Q_\alpha)\cong $ SmallGroup(32,50). According to Proposition \ref{no rep}, $\gamma_1(G)$ has non-trivial center, contradiction.
\end{proof}

%\comment{what can we say about $\LSS(16,p,\sigma_Q)$?}


\subsection{Factors of latin quandles of size \texorpdfstring{$16p$}{16p}}

The size of factors of latin quandles of size $16p$ is limited according to the known results on latin quandles of size dividing $16p$. In this subsection, we identify the size of the factors of latin quandles of size $16p$ and we also recall the known results about such quandles.




\begin{lemma}\label{meet of max is min}
Let $Q$ be a latin quandle of size $16p$. Then:

\begin{itemize}

\item[(i)] $|Q/\alpha|\in \{1,4,16,p,4p,16p\}$ for every $\alpha	\in \con(Q)$.	

\item[(ii)] $Q$ is congruence slim.

%\item[(iii)] $Q$ has at most one factor of size $p$. \comment{Needed?}

%\item[(iii)] If $|\Ma_Q|\geq 2$ then $\alpha\wedge\beta=\mu_Q=\gamma_Q \in \ma_Q$ for every $\alpha,\beta\in \Ma_Q$ and $\alpha\neq \beta$.  
 
%\item[(iv)] If $|\Ma_Q|\geq 2$ then $\mu_Q=\gamma_Q$. 
 
\end{itemize}



\end{lemma}


\begin{proof}
(i) The size of $Q/\alpha$ divides the size of $Q$, so $|Q/\alpha|=2^k p^s$ for $0\leq k\leq 4$ and $0\leq s\leq 1$. The factors or congruence blocks of $Q$ cannot have size $2$ or $2p$ because there are no connected quandles for such orders for $p>5$ (see \cite{McCarron}) and there are no latin quandles of size $6$ and $10$ (see the database of connected quandles \cite{RIG} in GAP). 

(ii) Let $\alpha\in \ma_Q$. Then according to (i), we have $|Q/\alpha|\in \{4,16,p,4p\}$. In all cases, the height of $Q/\alpha$ is less than or equal to $3$ and therefore the height of $Q$ is less than or equal to $4$. 
%
%(ii) Let $$0_Q=\alpha_0<\alpha_1<\ldots<\alpha_{n-1}<\alpha_n=1_Q$$
%where $\alpha_{n-1}\in \Ma_Q$. Then $|Q/\alpha_{n-1}|$ is a power of a prime, so $|Q/\alpha_{n-1}|\in \{4,16,p\}$ and accordingly $Q'=[x]_{\alpha_n}$ have size $4p,p,16$ respectively. We have a chain of congruences $$0_{Q'}=\alpha_0|_{Q'}<\alpha_1|_{Q'}<\ldots<\alpha_{n-2}|_{Q'}=\alpha_{n-1}|_{Q'}=1_{Q'}.$$
%In all cases then $n-1=3$.
%
%(ii) \comment{Better!} Assume that there exists a chain of length $5$, namely
%$$0_Q=\alpha_1<\alpha_2<\alpha_3<\alpha_4<\alpha_5=1_Q.$$
%
%Then $|Q/\alpha_{i+1}|$ divides $|Q/\alpha_i|$ for every $1\leq i\leq 5$ and $|Q/\alpha_i|\in \{4,16,p,4p\}$ for $2\leq i\leq 4$. A case by case discussion shows that necessarily $|Q|>16p$.%$\alpha_i=\alpha_{i+1}$ for some $1\leq i\leq 5.$ 
%
%(ii) Take $\alpha \in \ma_Q$. Then $Q/\alpha$ has size $4,16,p,4p$ and $\con(Q/\alpha)$ has height less than or equal to $3$. 
%
%(iii) If $|Q/\alpha|=|Q/\beta|=p$ then, by Lemma \ref{DP}, $|Q/(\alpha\wedge \beta)|=p^2$ divides $16p$, which is a contradiction since $p$ is odd.
%
%(iii) According to Lemma \ref{LSS 16p}, the size of $Q/\alpha$ and of $Q/\beta$ is not $16$. Since $Q/(\alpha\wedge\beta)$ is the direct product of $Q/\alpha$ and $Q/\beta$ (see Lemma \ref{DP} again) has size either $16$, or $4p$. Hence the blocks of $\alpha\wedge \beta$ have size $p$ or $4$ and therefore it is a minimal congruence and $Q/(\alpha\wedge \beta)$ is affine. Then
%$$\gamma_Q\leq \mu_Q\leq \alpha\wedge \beta$$
%and since $Q$ is not affine and $\alpha\wedge \beta\in \ma_Q$ then the equality holds.
\end{proof}
According to Lemma \ref{meet of max is min}(ii) we can apply the results of Section \ref{slim} to latin quandles of size $16p$, and so we know how their congruence lattices look like.


Let us list the latin quandles of size $\{4,16,p,4p\}$ where $p$ is an odd prime. 
\begin{itemize}
    \item[(i)] We denote by $\Q_4$ the unique latin quandle of order $4$. In particular $\Q_4$ is affine over $\Z_2^2$ and it is simple.% and latin.%, namely $\aff(\Z_2^2,f)$ where $f$ is the matrix
%$$f=\begin{bmatrix}
%0 & 1\\
%1 & 1
%\end{bmatrix}.$$


\item[(ii)] Latin quandles of size $16$ are affine over elementary abelian groups $\Z_2^4$ \cite{RIG}.

\item[(iii)] Latin quandles of size $p$ are affine over cyclic groups of prime size \cite{EGS}. 

%\item[(iv)] There are no connected quandle of order $2p$ for $p>5$ \cite{McCarron}. Moreover there are no latin quandles of size $6$ and $10$ \cite{RIG}.




 \item[(iv)] Latin quandles of size $4p$ with $p=3,5$ are affine \cite{RIG}. 
 Latin quandles of size $4p$ for $p>5$ are affine or are locally strictly simple, and they are completely classified in \cite{LSS}. The affine quandles of size $4p$ are isomorphic to $\mathcal{Q}_4\times \aff(\Z_p,f)$ and are generated by every pair $(x,0),(0,y)$ for $x,y\neq 0$. We denote the non-affine latin quandles of size $4p$ by $Q(p,j)$ for $j=1,2$. Such quandles are in the class $\LSS(p,4,0_Q)$ and exist only when $p=1\pmod{3}$. Let $k\in \Z_p$ such that $k^3=1$ and $k\neq 1$. The coset representation of $Q(p,j)$ is $\mathcal{Q}(\mathcal{G}_k,Fix(f(j)),f(j))$ where $\mathcal{G}_k\cong \Z_p^2\rtimes_\rho \mathcal{Q}_8$ is the group defined in Subsection \ref{grup for sr}, where
 $$\mathcal{Q}_8=\langle x,y,z\,|\, x^2,\, y^2,\, [x,y]z^{-1},\, z^2,\,[z,y],\,[z,x]\rangle,$$
 the action $\rho$ is defined on the generators of $\mathcal{Q}_8$ by setting 
\begin{displaymath}
\rho_x=\begin{bmatrix}
0 & -1\\
1 & 0
\end{bmatrix},\quad
\rho_{y}=\begin{bmatrix}
k^2 & k\\
k & -k^2
\end{bmatrix},\quad \rho_z=\begin{bmatrix}
-1 & 0\\
0 & -1
\end{bmatrix},
\end{displaymath}
 and 
\begin{align}\label{fj}
      f(j)=\begin{cases}
          x\mapsto y,\\
                  y\mapsto xy,\\
                  z\mapsto z\\   f|_{\gamma_2(\mathcal{G}_k)}=\begin{bmatrix}
                      -k(1+k^j) & -(1+k^j)\\
                      0 & -k^2(1+k^j)
                  \end{bmatrix}
                  \end{cases}
\end{align}
 for $j=1,2$.

 
%\item[(vi)] latin quandles of size $8p$ are affine \cite[Theorem 3.20]{LSS}. 
\end{itemize}



In the paper, we are also using the classification of non-faithful connected quandles of size $32$ that can be found in \cite{RIG}. We are collecting the data we will use later in Table \ref{32}, in particular, we list the connected non-faithful quandles of size $32$ such that the abelianization of the displacement group is the elementary abelian $2$-group of size $16$.

%\[
%\begin{array}{|c|c|c|c|c|}
%\hline
%Q \text{ (Size 16)} & \text{Dis(Q)} & |\con(Q)| & \text{Simple} & \text{SI} \\
%\hline
%\text{SmallQuandle(16,1)} & \text{SmallGroup(16,2)} & 3 & \text{false} & \text{true} \\
%\text{SmallQuandle(16,2)} & \text{SmallGroup(16,14)} & 7 & \text{false} & \text{false} \\
%\text{SmallQuandle(16,3)} & \text{SmallGroup(16,14)} & 2 & \text{true} & \text{true} \\
%\text{SmallQuandle(16,4)} & \text{SmallGroup(16,14)} & 3 & \text{false} & \text{true} \\
%\text{SmallQuandle(16,5)} & \text{SmallGroup(16,2)} & 3 & \text{false} & \text{true} \\
%\text{SmallQuandle(16,6)} & \text{SmallGroup(16,2)} & 3 & \text{false} & \text{true} \\
%\text{SmallQuandle(16,7)} & \text{SmallGroup(16,2)} & 3 & \text{false} & \text{true} \\
%\text{SmallQuandle(16,8)} & \text{SmallGroup(16,14)} & 2 & \text{true} & \text{true} \\
%\text{SmallQuandle(16,9)} & \text{SmallGroup(16,14)} & 2 & \text{true} & \text{true} \\
%\hline
%\end{array}
%\]


\begin{table}[ht!]
    \centering
    \begin{tabular}{|c|c|c|c|}
\hline
$Q$   & $\dis(Q)$ &  $|\con(Q/\gamma_Q)|$  \\
\hline
%\text{SmallQuandle(32,1)} & \text{SmallGroup(32,2)} & $\Z_4^2$ & 5 & \text{SmallQuandle(16,1)}&   \text{false} \\
\text{{\tt SmallQuandle(32,2)}} & \text{{\tt SmallGroup(32,47)}}  &  7   \\
\text{{\tt SmallQuandle(32,3)}} & \text{{\tt SmallGroup(32,49)}}   &   7  \\
\text{{\tt SmallQuandle(32,4)}} & \text{{\tt SmallGroup(32,50)}}   &   2  \\
\text{{\tt SmallQuandle(32,5)}} & \text{{\tt SmallGroup(32,47)}}   &   3  \\
\text{{\tt SmallQuandle(32,6)}} & \text{{\tt SmallGroup(32,49)}}    &   3  \\
%\text{SmallQuandle(32,7)} & \text{SmallGroup(32,2)} & $\Z_4^2$& 5 & \text{SmallQuandle(32,1)} & \text{false} \\
%\text{SmallQuandle(32,8)} & \text{SmallGroup(32,2)} &$\Z_4^2$ & 5 & \text{SmallQuandle(32,1)} & \text{false} \\
%\text{SmallQuandle(32,9)} & \text{SmallGroup(32,2)} &$\Z_4^2$ & 5 & \text{SmallQuandle(32,1)} & \text{false} \\
\hline 
    \end{tabular}
    \caption{Non-affine connected quandles of size $32$ such that the largest affine factor is affine over the group $\Z_2^4$.}
    \label{32}
\end{table}



%First let us show that subdirectly reducible latin quandles of size $16p$ are directly decomposable.


\subsection{Subdirectly reducible case}

%\comment{ If $p\neq 1\pmod{3}$ and $p\notin \{3,5,7,17,31,127\}$ then $Q$ is directly decomposable and abelian.}
Latin quandles of size $4,16,p$ and $4p$ are classified, so we can focus on directly indecomposable latin quandles of size $16p$, where $p$ is an odd prime.  




\begin{lemma}\label{16p nilpotent are abelian}
Let $Q$ be a latin quandle of size $16p$. The following are equivalent: 
\begin{itemize}
    \item[(i)] $Q$ is nilpotent.
  \item[(ii)] $Q$ is abelian.
\end{itemize}
Moreover, if $Q$ is nilpotent, then $Q$ is directly decomposable. 
\end{lemma}

%\comment{equivalent also to $Z(\dis(Q))\neq 1$.}
\begin{proof}
Abelian quandles are nilpotent. Assume that $Q$ is nilpotent. Finite latin nilpotent quandles are direct products of prime power order quandles \cite[Theorem 1.4]{CP}. Latin quandles of size $16$ and $p$ are abelian, therefore so is their direct product. Hence $Q$ is abelian and directly decomposable.
\end{proof}


%Let $p$ be a prime such that $p=1\pmod{3}$. Let us denote by $Q(p,i)$ for $i=1,2$ the subdirectly irreducible quandles of size $4p$. 
%Let $E(p,i,f)=Q(p,i)\times \aff(\Z_2^4,f)$ for $i=1,2$ and $\aff(\Z_2^4,f)$ is an S.I. latin quandle.  


\begin{proposition}\label{Sub red}
Let $Q$ be a subdirectly reducible latin quandle of size $16p$. If $Q$ is directly indecomposable then 
%
%\begin{itemize}
    %
    %\item[(i)] 
    $\con(Q)$ is the following:
\begin{align}\label{cong of not SI}
    \xymatrixrowsep{0.15in}
		\xymatrixcolsep{0.15in}
		\xymatrix{ 
			 & 1_Q  \ar@{-}[d]^4& &\\
			 & \nu_Q \ar@{-}[dr]^p\ar@{-}[dl]_4 & &\\
			\gamma_Q \ar@{-}[dr]_p  & & \zeta_Q\ar@{-}[dl]^4 \\
			 & 0_Q&  &
		}
\end{align}   
%    \item[(ii)] $Q$ subdirectly embeds into $E(p,i,f)$ for some $i=1,2$ and some $f\in \aut{\Z_2^4}$;

%\item[(ii)] $|\dis(Q)|=32p^2$, $|\gamma_1(\dis(Q))|=2p^2$ and $Z(\dis(Q))\cong \Z_2^2$. 
%
%\end{itemize}



%
%Then $$\dis(Q)\cong \Z_p^2\rtimes_{\rho} K$$ where $K=${\tt SmallGroup}(32,47) and $Q$ subdirectly embeds into $E(p,i)$ for some $i=1,2$. In particular, $p=1\pmod{3}$.

% Therefore, for a latin quandle of size $16p$, being directly indecomposable is equivalent to being subdirectly irreducible.
 
%There is no directly indecomposable, subdirectly reducible latin quandle of size $16p$.
%
%Let $Q$ be a directly indecomposable, subdirectly reducible latin quandle of size $16p$. Then:
%\begin{itemize}
%\item[(i)] $Q\cong \Q_4\times \widetilde{Q}$ where $\widetilde{Q}$ is a non-affine S.I, quandle of size $4p$.
%
%\item[(i)] The congruence lattice of $Q$ has the form 
%	\begin{center}
%$\xymatrixrowsep{0.15in}
%		\xymatrixcolsep{0.15in}
%		\xymatrix{ 
%			 & 1_Q  \ar@{-}[d]^4& &\\
%			 & \nu_Q \ar@{-}[dr]^p\ar@{-}[dl]_4 & &\\
%			\gamma_Q \ar@{-}[dr]_p  & & \zeta_Q\ar@{-}[dl]^4 \\
%			 & 0_Q&  &
%		}$
%$\xymatrixrowsep{0.15in}
%		\xymatrixcolsep{0.15in}
%		\xymatrix{ 
%			& & 1_Q  \ar@{-}[d]& &\\
%			& & \delta  \ar@{-}[d]\ar@{-}[dr]\ar@{-}[drr] \ar@{-}[dl] \ar@{-}[dll] & &\\
%			\gamma_Q \ar@{-}[drr] & \alpha_1\ar@{-}[dr]  &\alpha_2\ar@{-}[d]  & \ldots\ar@{-}[dl]  & \alpha_n \ar@{-}[dll] \\
%			& & 0_Q&  &
%		}$
%	\end{center}
%	
%\item[(ii)] $Q/\gamma_Q$ is a S.I. latin quandle of size $16$.
%\item[(iii)]  $Q/\zeta_Q$ is a S.I. latin quandles of size $4p$ and $p=1\pmod{3}$.



%\item[(ii)] $p=7$ and the congruence lattice of $Q$ has the form 
%\begin{center}
%		$\xymatrixrowsep{0.15in}
%		\xymatrixcolsep{0.15in}
%		\xymatrix{ 
%			& & 1_Q  \ar@{-}[d]& &\\
%			& & \gamma_Q  \ar@{-}[d]\ar@{-}[dr]\ar@{-}[drr] \ar@{-}[dl] \ar@{-}[dll] & &\\
%	\alpha_1 \ar@{-}[drr] & \alpha_2\ar@{-}[dr] &\ldots\ar@{-}[d] &  \alpha_{n-1}\ar@{-}[dl] & \alpha_n \ar@{-}[dll]\\
%			\alpha_1 \ar@{-}[drr] & \alpha_1\ar@{-}[dr] &\alpha_2\ar@{-}[d] & \ldots\ar@{-}[dl] & \alpha_n \ar@{-}[dll]\\
%			& & 0_Q&  &
%		}$
%	\end{center}
	
%
%\noindent and $2\leq n\leq 5$, $Q/\alpha_i\in \LSS(4,7,\gamma_Q)$ for every $\alpha_i\in \ma_Q$. \comment{We can classify computationally this case, by looking at subquandles of $Q_1\times Q_2$ with $|Q_1|=|Q_2|=28$.} \comment{The quandles in $\LSS(4,7,\sigma_Q)$ are not latin!}

%\item[(iv)]  $Q$ is $2$-step solvable.
%\end{itemize}



\end{proposition}
%\begin{itemize}

%\item[(1)]  \comment{$Q/\gamma_Q$ has size $16$, so $Q$ embeds into the product of classified quandles! }
 %   \item[(2)] \comment{$Q/\gamma_Q$ is an S.I. latin quandle of size $16$. Then either it is SmallQuandle(16,4) or it is affine over $\Z_4^2$. Can the second option happen? According to what follows it does not.}
%
 %   \item[(3)] \comment{
%$\dis_{\zeta_Q}=\dis^{\zeta_Q}=Z(\dis(Q))$ has size $4$ and $|\dis(Q)|=|\dis_{\zeta_Q}||\dis(Q/\zeta_Q)|=4*8p^2=16|\gamma_1(\dis(Q))|$ and so $|\gamma_1(\dis(Q))|=2p^2$. Moreover $|\dis(Q)|=4|\dis^{\nu_Q}|=32p^2$ and so $|\dis^{\nu}|=8p^2$. Since it contains $\langle \gamma_1(\dis(Q)),Z(\dis(Q))\rangle=\gamma_1(\dis(Q))\times Z(\dis(Q))\leq \dis^{\nu_Q}$ then equality holds by size. }
%
%\item[(4)] \comment{ $\gamma_1(\dis(Q))$ have a normal $p$-Sylow $P$ that is characteristic (has index $2$). Thus I can consider the quandle $Q'$ of size $32$ again over the group $\dis(Q)/P$. $Q'$ is connected and not faithful ($Q'$ is nilpotent but has subquandle of size $2$). Moreover $Q'/\gamma_{Q'}$ is an S.I. affine quandle. There is only one of such quandles SmallQuandle(32,6) with order of $|L_{[x]}|=6$. Thus, $Q/\gamma_Q\cong Q'/\gamma_{Q'}$ is SmallQuandle(16,4). The order of $L_x$ is $6$ since $Q$ embeds into the product of quandles with order $6$ and $3$.  SmallQuandle(16,4) is not simply connected, and SmallQuandle(32,6)/$\gamma_Q$=SmallQuandle(16,4).} 
%
%\item[(5)] \comment{We have that $\gamma_1(\dis(Q)))=\Z_p^2\rtimes \Z_2$ and $\dis(Q)/\gamma_1(\dis(Q))=K=$SmallGroup(32,49). I think the sequence $\Z_p^2\longrightarrow\dis(Q)\longrightarrow K$ splits. So can investigate the representations of such group over $\Z_p^2$. The subgroup $Z=Z(\gamma_1(\dis(Q))$ is normal in $\lmlt(Q)$. Indeed it is characteristic a characteristic subgroup of $\dis(Q)$ that is normal in $\lmlt(Q)$. If so $\dis_\alpha\leq Z$ or $Z=1$.  Moreover, $Z(\dis(Q))$ has size $4$ and $\ker{\rho}\cap Z(K)\leq Z(\dis(Q))$. Let $x=vh\in Z(\dis(Q))$ then $h\in Z(K)$ and so $Z(\dis(Q))\leq \gamma_1(\dis(Q))$, contadiction.}

%\comment{Another way to get a contradiction: $Z(\dis(Q))=\setof{(a_1,1),(a_2,1),(b_1,z),(b_2,z)}{z\in Z(K)}\cong \Z_2^2$. Hence $(a_i,1)(a_i,1)=(2a_i,1)=(1,1)$, then $a_1=a_2=1$, contradiction!} 
%\end{itemize}

% \comment{si può fare forse a meno di invocare la classificazione dei quandles of size $32$, se si riesce a dire che $K$ è extraspeciale in un altro modo}





\begin{proof}
%(i) 
Let $Q$ be a subdirectly reducible latin quandle of size $16p$. If $Q$ is directly indecomposable then the congruence lattice of $Q$ is as follows according to Proposition \ref{lattices}(ii):
\begin{center}
		$\xymatrixrowsep{0.15in}
		\xymatrixcolsep{0.15in}
		\xymatrix{ 
			& & 1_Q  \ar@{-}[d]& &\\
			& & \nu_Q  \ar@{-}[d]\ar@{-}[dr]\ar@{-}[drr] \ar@{-}[dl] \ar@{-}[dll] & &\\
	\alpha_1 \ar@{-}[drr] & \alpha_2\ar@{-}[dr] &\ldots\ar@{-}[d] &  \alpha_{n-1}\ar@{-}[dl] & \alpha_n \ar@{-}[dll]\\
%			\alpha_1 \ar@{-}[drr] & \alpha_1\ar@{-}[dr] &\alpha_2\ar@{-}[d] & \ldots\ar@{-}[dl] & \alpha_n \ar@{-}[dll]\\
			& & 0_Q&  &
		}$
	\end{center}
where $n\geq 2$. Moreover, $|Q/\nu_Q|\in \{4,16,p\}$ since finite simple latin quandles have size a power of a prime (as they are of the form $\aff(\Z_p^n,f)$). Let us proceed case by case according to the size of $Q/\nu_Q$.
\begin{itemize}
    \item If $|Q/\nu_Q|=16$, then $\nu_Q$ has blocks of prime size and so $\nu_Q$ is minimal. Thus, $Q$ is not subdirectly reducible.


\item If $|Q/\nu_Q|=p$, then $Con(Q/\alpha_i)$ is a chain of length $3$ and so $Q/\alpha_i\in \LSS(4,p,\sigma_Q)$ for $i=1,\ldots, n$. Therefore $p=7$ (see \cite[Proposition 5.3]{LSS}). Such quandles are not latin, as shown in \cite[Section 5]{LSS}.



\item If $|Q/\nu_Q|=4$, then $[x]_{\nu_Q}$ has size $4p$. Latin quandles are solvable and faithful, so according to Corollary \ref{minimal cong}(iii) the congruence $\nu_Q$ is abelian and $[x]_{\nu_Q}$ is an abelian quandle of size $4p$. In particular $[x]_{\nu_Q}\cong \Q_4\times \aff(\Z_p,f)$ and $[x]_{\nu_Q}$ has $2$ proper congruences. Hence $n=2$ according to Lemma \ref{below abelian}(ii). 

\end{itemize}

We can assume that $\ma_Q=\{\alpha_1,\alpha_2\}$ with $|Q/\alpha_1|=16$ and $|Q/\alpha_2|=4p$. Then $Q/\alpha_1$ is abelian (latin quandles of size $16$ are affine \cite{RIG}), so we have $\gamma_Q\leq \alpha_1$. Then we can conclude that $\gamma_Q=\alpha_1$, since $Q$ is not abelian.  

The size of $Q/\alpha_2$ is $4p$. Since $\alpha_2\wedge \gamma_Q=0_Q$, according to Lemma \ref{central}, $\alpha_2\leq \zeta_Q$ and moreover $Q/\alpha_2$ is not abelian given that $\gamma_Q$ is not contained in $\alpha_2$. Hence $Q/\alpha_2\in \LSS(p,4,\sigma_Q)$. If $\alpha_2<\zeta_Q$ then $\zeta_Q=\nu_Q$ and therefore $\gamma_Q\leq \zeta_Q$. Thus, $Q$ is nilpotent and is therefore directly decomposable according to Lemma \ref{Sub red}, a contradiction. 

Therefore, the congruence lattice of $Q$ is as in \eqref{cong of not SI}.
\end{proof}
%\begin{align}\label{cong of not SI}
 %   \xymatrixrowsep{0.15in}
%		\xymatrixcolsep{0.15in}
%		\xymatrix{ 
%			 & 1_Q  \ar@{-}[d]^4& &\\
%			 & \nu_Q \ar@{-}[dr]^p\ar@{-}[dl]_4 & &\\
%			\gamma_Q \ar@{-}[dr]_p  & & \zeta_Q\ar@{-}[dl]^4 \\
%			 & 0_Q&  &
%		}
%\end{align}    

%(ii) The quandle $Q$ embeds into $Q/\gamma_1(Q)\times Q/\zeta_Q$ where $Q/\gamma_Q=Q(p,i)$ for $i=1,2$ and $Q/\zeta$ is an S.I. affine quandle of size $16$. So $Q/\zeta_Q\cong \aff(\Z_2^4,f)$.

%
%(iii) Let $G=\dis(Q)$. According to Lemma \ref{central}, $\dis_{\zeta_Q}=\dis^{\zeta_Q}=Z(G)$ and it has size $4$ (a central group of a transitive group is semiregular). Moreover, $Z(G)$ embeds into a power of $\dis([x]_{\zeta_Q})$ and so $Z(G)\cong \Z_2^2$. %Moreover $\zeta_Q=\mathcal{O}_{Z(G)}$ and so $\dis_{\zeta_Q}\leq Z(G)\leq \dis^{\zeta_Q}$. Therefore $|Z(G)|=4$. %$Z(G)=\dis_{\zeta_Q}\cong \Z_2^2$.

%Since $Q/\zeta_Q\in \LSS(p,4,\sigma_Q)$ then $|\dis(Q/\zeta_Q)|=8p^2$ (see \cite[Proposition 5.4]{LSS}) and therefore $$|G|=|\dis(Q/\zeta_Q)||\dis^{\zeta_Q}|= 8p^2\cdot 4=|G/\gamma_1(G)||\gamma_1(G)|=16|\gamma_1(G)|.$$ So we have that $|G|=32p^2$ and $\gamma_1(G)$ has size $2p^2$.% and accordingly $\gamma_1(G)\cong P\rtimes \Z_2$ where $P$ is the $p$-Sylow of $\gamma_1(G)$. 
%
%\comment{Add that $Z(G)\cong \Z_2^2$. Forse da qui in poi non serve.}
%
%According to Lemma \ref{central 2}, $\dis_{\gamma_Q}$ has size a power of $p$, it is not cyclic and so $P=\dis_{\gamma_Q}\cong \Z_p^2$. We can consider the quandle $Q_P$ as defined in \cite[...]{Involutory}. Such quandle is a connected quandle of size $32$, such that $Q_P/\gamma_{Q_P}$ is affine over $\Z_2^4$, it has 3 congruences and $\dis(Q_P)=G/P=K$. Connected quandles of size $32$ are classified and according to Table \ref{32} we have that $K$ is {\tt SmallGroup}(32,i) for $i=47,49$ and $Q/\gamma_Q=\text{{\tt SmallQuandle(16,4)}}$. 
%
%If $K=\text{{\tt SmallGroup(32,49)}}$ then $\rho$ is not faithful by Corollary \ref{ro faithful on Z_p} since $K$ has a subgroup isomorphic to $\Z_2^3$. The subgroup $\gamma_1(K)$ is a minimal central subgroup of $K$. Since $Z(K)\cap\ker{\rho}\leq Z(G)$, we have $\gamma_1(K)\leq \ker{\rho}$ and so $$\gamma_1(K)\leq Z(K)\cap \ker{\rho}\cap \gamma_1(G)\leq Z(G)\cap \gamma_1(G)=1,$$ contradiction. Thus $K=\text{{\tt SmallGroup(32,47)}}$.
%
%Therefore $Q$ embeds subdirectly into $Q/\gamma_Q\times Q/\zeta_Q\cong \text{{\tt SmallQuandle(16,4)}}\times Q(p,i)=E(p,i)$ for $i=1,2$.
%
%
%and $[a]_{\gamma_Q}$ is isomorphic to $\Q_4\times \Q_4$ (that is the unique subdirectly reducible latin quandle of size $16$). Thus, congruence lattice is as in (iii) since $[a]_{\nu_Q}$ has $5$ proper congruences (see Lemma \ref{below abelian}.
%\end{proof}





%\begin{proposition}\label{Sub red}
 %Subdirectly reducible latin quandle of size $16p$ are directly decomposable.% Therefore, for a latin quandle of size $16p$, being directly indecomposable is equivalent to being subdirectly irreducible.
 
%There is no directly indecomposable, subdirectly reducible latin quandle of size $16p$.
%
%Let $Q$ be a directly indecomposable, subdirectly reducible latin quandle of size $16p$. Then:
%\begin{itemize}
%\item[(i)] $Q\cong \Q_4\times \widetilde{Q}$ where $\widetilde{Q}$ is a non-affine S.I, quandle of size $4p$.
%
%\item[(i)] The congruence lattice of $Q$ has the form 
%	\begin{center}
%$\xymatrixrowsep{0.15in}
%		\xymatrixcolsep{0.15in}
%		\xymatrix{ 
%			 & 1_Q  \ar@{-}[d]^4& &\\
%			 & \nu_Q \ar@{-}[dr]^p\ar@{-}[dl]_4 & &\\
%			\gamma_Q \ar@{-}[dr]_p  & & \zeta_Q\ar@{-}[dl]^4 \\
%			 & 0_Q&  &
%		}$
%$\xymatrixrowsep{0.15in}
%		\xymatrixcolsep{0.15in}
%		\xymatrix{ 
%			& & 1_Q  \ar@{-}[d]& &\\
%			& & \delta  \ar@{-}[d]\ar@{-}[dr]\ar@{-}[drr] \ar@{-}[dl] \ar@{-}[dll] & &\\
%			\gamma_Q \ar@{-}[drr] & \alpha_1\ar@{-}[dr]  &\alpha_2\ar@{-}[d]  & \ldots\ar@{-}[dl]  & \alpha_n \ar@{-}[dll] \\
%			& & 0_Q&  &
%		}$
%	\end{center}
%	
%\item[(ii)] $Q/\gamma_Q$ is a S.I. latin quandle of size $16$.
%\item[(iii)]  $Q/\zeta_Q$ is a S.I. latin quandles of size $4p$ and $p=1\pmod{3}$.



%\item[(ii)] $p=7$ and the congruence lattice of $Q$ has the form 
%\begin{center}
%		$\xymatrixrowsep{0.15in}
%		\xymatrixcolsep{0.15in}
%		\xymatrix{ 
%			& & 1_Q  \ar@{-}[d]& &\\
%			& & \gamma_Q  \ar@{-}[d]\ar@{-}[dr]\ar@{-}[drr] \ar@{-}[dl] \ar@{-}[dll] & &\\
%	\alpha_1 \ar@{-}[drr] & \alpha_2\ar@{-}[dr] &\ldots\ar@{-}[d] &  \alpha_{n-1}\ar@{-}[dl] & \alpha_n \ar@{-}[dll]\\
%			\alpha_1 \ar@{-}[drr] & \alpha_1\ar@{-}[dr] &\alpha_2\ar@{-}[d] & \ldots\ar@{-}[dl] & \alpha_n \ar@{-}[dll]\\
%			& & 0_Q&  &
%		}$
%	\end{center}
	
%
%\noindent and $2\leq n\leq 5$, $Q/\alpha_i\in \LSS(4,7,\gamma_Q)$ for every $\alpha_i\in \ma_Q$. \comment{We can classify computationally this case, by looking at subquandles of $Q_1\times Q_2$ with $|Q_1|=|Q_2|=28$.} \comment{The quandles in $\LSS(4,7,\sigma_Q)$ are not latin!}

%\item[(iv)]  $Q$ is $2$-step solvable.
%\end{itemize}



%\end{proposition}
%\begin{itemize}

%\item[(1)]  \comment{$Q/\gamma_Q$ has size $16$, so $Q$ embeds into the product of classified quandles! }
 %   \item[(2)] \comment{$Q/\gamma_Q$ is an S.I. latin quandle of size $16$. Then either it is SmallQuandle(16,4) or it is affine over $\Z_4^2$. Can the second option happen? According to what follows it does not.}
%
 %   \item[(3)] \comment{
%$\dis_{\zeta_Q}=\dis^{\zeta_Q}=Z(\dis(Q))$ has size $4$ and $|\dis(Q)|=|\dis_{\zeta_Q}||\dis(Q/\zeta_Q)|=4*8p^2=16|\gamma_1(\dis(Q))|$ and so $|\gamma_1(\dis(Q))|=2p^2$. Moreover $|\dis(Q)|=4|\dis^{\nu_Q}|=32p^2$ and so $|\dis^{\nu}|=8p^2$. Since it contains $\langle \gamma_1(\dis(Q)),Z(\dis(Q))\rangle=\gamma_1(\dis(Q))\times Z(\dis(Q))\leq \dis^{\nu_Q}$ then equality holds by size. }
%
%\item[(4)] \comment{ $\gamma_1(\dis(Q))$ have a normal $p$-Sylow $P$ that is characteristic (has index $2$). Thus I can consider the quandle $Q'$ of size $32$ again over the group $\dis(Q)/P$. $Q'$ is connected and not faithful ($Q'$ is nilpotent but has subquandle of size $2$). Moreover $Q'/\gamma_{Q'}$ is an S.I. affine quandle. There is only one of such quandles SmallQuandle(32,6) with order of $|L_{[x]}|=6$. Thus, $Q/\gamma_Q\cong Q'/\gamma_{Q'}$ is SmallQuandle(16,4). The order of $L_x$ is $6$ since $Q$ embeds into the product of quandles with order $6$ and $3$.  SmallQuandle(16,4) is not simply connected, and SmallQuandle(32,6)/$\gamma_Q$=SmallQuandle(16,4).} 
%
%\item[(5)] \comment{We have that $\gamma_1(\dis(Q)))=\Z_p^2\rtimes \Z_2$ and $\dis(Q)/\gamma_1(\dis(Q))=K=$SmallGroup(32,49). I think the sequence $\Z_p^2\longrightarrow\dis(Q)\longrightarrow K$ splits. So can investigate the representations of such group over $\Z_p^2$. The subgroup $Z=Z(\gamma_1(\dis(Q))$ is normal in $\lmlt(Q)$. Indeed it is characteristic a characteristic subgroup of $\dis(Q)$ that is normal in $\lmlt(Q)$. If so $\dis_\alpha\leq Z$ or $Z=1$.  Moreover, $Z(\dis(Q))$ has size $4$ and $\ker{\rho}\cap Z(K)\leq Z(\dis(Q))$. Let $x=vh\in Z(\dis(Q))$ then $h\in Z(K)$ and so $Z(\dis(Q))\leq \gamma_1(\dis(Q))$, contadiction.}

%\comment{Another way to get a contradiction: $Z(\dis(Q))=\setof{(a_1,1),(a_2,1),(b_1,z),(b_2,z)}{z\in Z(K)}\cong \Z_2^2$. Hence $(a_i,1)(a_i,1)=(2a_i,1)=(1,1)$, then $a_1=a_2=1$, contradiction!} 
%\end{itemize}

 %\comment{si può fare forse a meno di invocare la classificazione dei quandles of size $32$, se si riesce a dire che $K$ è extraspeciale in un altro modo}



%\comment{How can I handle the case $i=47$?}

%\comment{finora la proof mostra che possiamo supporre che i subdirectly reducible di ordine $16p$ si embeddano nel prodotto diretto di un SI di ordine $4p$ e SmallQuandle(16,4). GAP non riesce a calcolare i sottoquandle di questo prodotto diretto neanche per $p=7$.}
%\begin{proof}
%Let $Q$ be a subdirectly reducible latin quandle of size $16p$. If $Q$ is directly decomposable then the congruence lattice of $Q$ is as follows according to Proposition \ref{lattices}(ii):
%\begin{center}
	%	$\xymatrixrowsep{0.15in}
%		\xymatrixcolsep{0.15in}
%		\xymatrix{ 
%			& & 1_Q  \ar@{-}[d]& &\\
%			& & \nu_Q  \ar@{-}[d]\ar@{-}%[dr]\ar@{-}[drr] \ar@{-}[dl] %\ar@{-}[dll] & &\\
	%\alpha_1 \ar@{-}[drr] & %\alpha_2\ar@{-}[dr] &\ldots\ar@{-}[d] &  %\alpha_{n-1}\ar@{-}[dl] & \alpha_n %\ar@{-}[dll]\\
%			\alpha_1 \ar@{-}[drr] & \alpha_1\ar@{-}[dr] &\alpha_2\ar@{-}[d] & \ldots\ar@{-}[dl] & \alpha_n \ar@{-}[dll]\\
%			& & 0_Q&  &
%		}$
%	\end{center}
%for $n\geq 2$. Moreover, $|Q/\nu_Q|\in \{4,16,p\}$ since finite simple latin quandles have size a power of a prime. 

%If $|Q/\nu_Q|=16$ then $\nu_Q$ has blocks of prime size and so $\nu_Q$ is minimal. Thus $Q$ is not subdirectly reducible.


%If $|Q/\nu_Q|=p$, then $Q/\alpha_i\in \LSS(4,p,\sigma_Q)$ for $i=1,\ldots, n$ and therefore $p=7$ (\cite{LSS}). Such quandles are not latin \cite{LSS}.


%If $|Q/\nu_Q|=4$, $[x]_{\nu_Q}$ has size $4p$. latin quandles are solvable and faithful, so according to Corollary \ref{minimal cong}(iii) the congruence $\nu_Q$ is abelian and so $[x]_{\nu_Q}$ is an abelian quandle of size $4p$. Then $[x]_{\nu_Q}\cong \Q_4\times \aff(\Z_p,f)$ and $n=2$, since $[x]_{\nu_Q}$ has $2$ congruences according to Lemma \ref{below abelian}(ii). Given $\alpha_1$ such that $|Q/\alpha_1|=16$  then $\gamma_Q\leq \alpha_1$. Since $Q$ is not abelian then $\gamma_Q=\alpha_1$. Moreover $Q/\alpha_2\in \LSS(p,4,\sigma_Q)$. According to Lemma \ref{central}, $\alpha_2\leq \zeta_Q$ and so equality holds, otherwise $Q$ is nilpotent.

%Then, the congruence lattice of $Q$ is as follows:
%\begin{center}
    
%$\xymatrixrowsep{0.15in}
%		\xymatrixcolsep{0.15in}
%		\xymatrix{ 
%			 & 1_Q  \ar@{-}[d]^4& &\\
%			 & \nu_Q \ar@{-}[dr]^p\ar@{-}%[dl]_4 & &\\
		%	\gamma_Q \ar@{-}[dr]_p  & & %\zeta_Q\ar@{-}[dl]^4 \\
%			 & 0_Q&  &
%		}$
%\end{center}




%Let $G=\dis(Q)$. According to Lemma \ref{central}, $\dis_{\zeta_Q}=\dis^{\zeta_Q}$ and it has size $4$ (a central group of a transitive group is semiregular). Since $Q/\zeta_Q\in \LSS(p,4,\sigma_Q)$ then $|\dis(Q/\zeta_Q)|=8p^2$ (see \cite{LSS}) and therefore $$|G|=|\dis(Q/\zeta_Q)||\dis^{\zeta_Q}|=4\cdot 8p^2=|G/\gamma_1(G)||\gamma_1(G)|=16|\gamma_1(G)|.$$ So we have that $\gamma_1(G)$ has size $2p^2$ and accordingly $\gamma_1(G)\cong P\rtimes \Z_2$ where $P$ is the $p$-Sylow of $\gamma_1(G)$. According to Lemma \ref{central 2}, $\dis_{\gamma_Q}$ has size a power of $p$, it is not cyclic and so $P=\dis_{\gamma_Q}\cong \Z_p^2$. We can consider the quandle $Q_P$ as defined in ... Such quandle is a connected quandle of size $32$, such that $Q_P/\gamma_{Q_P}$ is affine over $\Z_2^4$, it has 3 congruences and $\dis(Q_P)=G/P=K$. Connected quandles of size $32$ are classified and according to Table \ref{32} we have that $K$ is SmallGroup(32,i) for $i=47,49$. 



%If $K=$SmallGroup(32,49) then $\rho$ is not faithful by Corollary \ref{ro faithful on Z_p} since $K$ has a subgroup isomorphic to $\Z_2^3$. The subgroup $\gamma_1(K)$ is a minimal central subgroup of $K$. Therefore $\gamma_1(K)\leq \ker{\rho}$ and so $\gamma_1(K)\leq \gamma_1(G)\cap Z(G)=1$, contradiction.
%
%
%and $[a]_{\gamma_Q}$ is isomorphic to $\Q_4\times \Q_4$ (that is the unique subdirectly reducible latin quandle of size $16$). Thus, congruence lattice is as in (iii) since $[a]_{\nu_Q}$ has $5$ proper congruences (see Lemma \ref{below abelian}.
%\end{proof}
%\begin{lemma}
%Let $Q\in \LSS(p,16,\sigma_Q)$, then:
%\begin{itemize}
 %   \item[(i)]  $\sigma_Q=0_Q$. 
    



 %   \item[(iii)] According to \cite[Proposition 3.21]{LSS} if the unique proper congruence of $Q$ is strict. In particular, if $k=|L_x|$ then $|L_x|_{[x]}|$ divides $k$ and since $[x]=\aff(\Z_p,f)$ then $|f|$ divides $k$. In particular, since the simple quandles of size $16$ are classified, $k\in \{5,15\}$, and so $|f|\in\{3,5,15\}$ accordingly. Since $p=1\pmod {|f|}$ we have that $p=1\pmod{3}$ (resp. $\pmod{5}$). 

  %  \item[(ii)] $\dis(Q)=\Z_p^2\rtimes_\rho K$, where $K=$SmallGroup(32,50) and the action $\rho$ is faithful. In particular, $32$ divides $(p-1)^2(p+1)$. \comment{If $2^s$ divides $p-1$ and $p+1$ then $s=1$. Hence if $v_2(p+1)=1$ then $p=1\pmod{4}$. If $v_2(p-1)=1$ then $p=-1\pmod{8}$.}

    
   % \item[(iii)] $Q$ is latin, $Q/\gamma_Q=$SmallQuandle(16,3), $|L_x|=5$ and $p=1\pmod{5}$. \comment{Since $p=1\pmod{5}$ then there exists a subgroup of $\dis_{\gamma_Q}$ that is invariant under $L_x$, i.e. $Fix(\widehat{L_x})$}


%\end{itemize}
    
%\end{lemma}
%\comment{According to Lemma 7.1 \cite{LSS} the $\rho_z$ is diagonalizable for eigenvalues $1$ and $-1$ given $z\in Z(K)$. If one eigenvalues of $\rho_z$ is $1$ the eigenspace is in the center, contradiction. Then we can assume that $\rho_z=-Id$. Let us restrict to the subgroup of $K$ isomorphic to the quaternion group. Let $x$ or order $4$ in $K$. Then $\rho_x^2=-1$. If $\rho_x$ is not diagonalizable then $\rho_x=\begin{bmatrix}
 %   0 & -1\\
  %  1&0
%\end{bmatrix}$. If $y\in K$ s.t. $[x,y]=z$ then $\rho_y=\begin{bmatrix}
 %   a& b\\ b & -a
%\end{bmatrix}$ with $a^2+b^2=1.$ If $[x,t]=0$ then $\rho_t=\begin{bmatrix}
 %   a	& -b\\
%b	& a
%\end{bmatrix}$ with $ab=0$. If $\rho_x$ is diagonalizable then $\rho_x=\begin{bmatrix}
 %   x &0 \\ 0 &-x
%\end{bmatrix}$ and $\rho_y=\begin{bmatrix}
 %   0 & -c^{-1}\\
  %  c & 0
%\end{bmatrix}$ and $\rho_t=\begin{bmatrix}
 %   a & 0 \\ 0 & b
%\end{bmatrix}$. Since I have a subgroup isomorphic to $\mathcal{Q}_8$ the representation is irreducible and we can assume that $\rho_x$ is not diagonalizable }

%\begin{proof}
 %   (i)  It follows by \cite[Proposition 3.18]{LSS}.

    
  %  (ii) According to \cite[Theorem 3.14]{LSS}, $\dis(Q)=\Z_p^2\rtimes_\rho K$ where $K$ is an extraspecial $2$-group of size $32$. The quandle $Q_{\gamma_Q}$ as defined in ... is a connected quandle of size $32$ with a unique proper congruence. The connected quandles of size $32$ are classified, so there is a unique quandle matching such properties, namely SmallQuandle(32,4) and $K=\dis(Q_\alpha)\cong $ SmallGroup(32,50).
    
   % (iii)  According to  \cite[Lemma 3.17, Proposition 3.16]{LSS}, $Q$ is latin. The quandle $Q/\gamma_Q$ is a factor of $Q_{\gamma_Q}$ and so necessarily if SmallQuandle(16,3). In particular, $|L_{[x]}|=5$ and $p=1\pmod{5}$. According to Proposition \cite[Proposition 3.21]{LSS} we have $|L_x|=|L_{[x]}|=5$.\end{proof}


   \begin{proposition}\label{group K}
       Let $Q$ be a subdirectly reducible directly indecomposable latin quandle of size $16p$ and $G=\dis(Q)$. Then: 
       \begin{enumerate}
                  \item[(i)] $Z(G)=\dis_{\zeta_Q}=\dis^{\zeta_Q}\cong \Z_2^2$.
           \item[(ii)]  $G/\gamma_2(G)\cong$ \texttt{SmallGroup(32,47)} $\cong \mathcal{Q}_8\times \Z_2^2$. %\item $Q/\gamma_Q\cong${\tt SmallQuandle(16,4)}. \comment{needed?}
           \item[(iii)] $p=1\pmod{3}$ and $G\cong \mathcal{G}_k \times Z(G)$.
              \end{enumerate}

   \end{proposition}

\begin{proof}

(i) Let $G=\dis(Q)$. The blocks of $\zeta_Q$ have size $4$, so they are isomorphic to $\mathcal(Q)_4\cong \aff(\Z_2^2,f)$ and $\gamma_Q\wedge\zeta_Q=0_Q$. According to Lemma \ref{central}, $\dis_{\zeta_Q}=\dis^{\zeta_Q}=Z(G)$ (see Lemma \ref{center for latin}) and has size $4$ (a central group of a transitive group is semiregular). Moreover, $Z(G)$ embeds into a power of $\dis([x]_{\zeta_Q})$ and so $Z(G)\cong \Z_2^2$. %Moreover $\zeta_Q=\mathcal{O}_{Z(G)}$ and so $\dis_{\zeta_Q}\leq Z(G)\leq \dis^{\zeta_Q}$. Therefore $|Z(G)|=4$. %$Z(G)=\dis_{\zeta_Q}\cong \Z_2^2$.

(ii) Since $Q/\zeta_Q\in \LSS(p,4,\sigma_Q)$ then $|\dis(Q/\zeta_Q)|=8p^2$ (see \cite[Proposition 5.4]{LSS}) and therefore $$|G|=|\dis(Q/\zeta_Q)||\dis^{\zeta_Q}|= 8p^2\cdot 4=|G/\gamma_1(G)||\gamma_1(G)|=16|\gamma_1(G)|.$$ So we have that $|G|=32p^2$ and $\gamma_1(G)$ has size $2p^2$. According to Lemma \ref{dis_gamma} then $\dis_{\gamma_Q}=\gamma_2(G)\cong \Z_p^2$ as it is a non-cyclic elementary abelian $p$-group ($\gamma_Q$ is a minimal abelian non-central congruence).

Let $Q=\mathcal{Q}(G,G_x,f)$ where $f=\widehat{L_x}$ for some $x\in Q$. According to \cite[Lemma 1.10]{Involutory} the principal quandle $Q'=\mathcal{Q}(G/\gamma_2(G),f_{\gamma_2(G)})$ is a connected cover of $Q/\gamma_Q$ of size $32$ and $K=G/\gamma_2(G)\cong \dis(Q')$. In particular $\lambda_{Q'}=\gamma_{Q'}$ is a minimal congruence with blocks of size $2$, $Q'/\lambda_{Q'}\cong Q/\gamma_Q\cong \aff(\Z_2^4,f_{\gamma_1(G)})$ and $\con(Q'/\gamma_{Q'})$ is a chain of length $3$. The quandles of size $32$ are classified in \cite{RIG} and by looking of the size of $\con(Q/\gamma_Q)$ we have $K\cong {\tt SmallGroup(32,i)}$ for $i=47,49$. Assume that $K=G/\gamma_2(G)\cong{\tt SmallGroup(32,49)}$. Let $x=yk\in Z(G)$ for $y\in \Z_p^2$ and $k\in K$. Then $k\in Z(K)$ and so $k\in \gamma_1(K)\leq \gamma_1(G)$. Thus $x\in Z(G)\cap \gamma_1(G)=1$, contradiction. Hence $K\cong {\tt SmallGroup(32,47)}$.

%Let us discuss the two cases:
%\begin{itemize}
%    \item If $K=G/\gamma_2(G)\cong${\tt SmallGroup(32,49)}. Let $x=yk\in Z(G)$ for $y\in \Z_p^2$ and $k\in K$. Then $k\in Z(K)$ and so $k\in \gamma_1(K)\leq \gamma_1(G)$. Thus $x\in Z(G)\cap \gamma_1(G)=1$, contradiction.

 %   \item Assume that $\gamma_{Q'}\neq \alpha\in \ma_{Q'}$. Then $\alpha\wedge \gamma_{Q'}=0_Q$ and $Q'$ embeds into $Q/\gamma_{Q'}\times Q/\alpha$. Since $Q'$ is not abelian then necessarily $Q/\alpha$ is not abelian. Hence $|Q/\alpha|=8$ and so $\con(Q/\alpha)$ is a chain of length $3$ and the unique proper congruence of $Q'$ is abelian. Thus the congruence lattice of $Q'$ is as 

 %   \begin{align*}
 %   \xymatrixrowsep{0.15in}
%		\xymatrixcolsep{0.15in}
%		\xymatrix{ 
%			& & 1_{Q'}  \ar@{-}[d]& &\\
%			&  &\nu_{Q'} \ar@{-}[dr] \ar@{-}[dl]&  &\\
 %            & \gamma_{Q'}\ar@{-}[dr] & &  \alpha\ar@{-}[dl] & \\
%			\alpha_1 \ar@{-}[drr] & \alpha_1\ar@{-}[dr] &\alpha_2\ar@{-}[d] & \ldots\ar@{-}[dl] & \alpha_n \ar@{-}[dll]\\
%			& & 0_E&  &
%		}
 %       \end{align*}	
%        By looking again at the classification of connected quandles of size $32$ given in Table \ref{32} we have that $Q'\cong${\tt SmallQuandle}(32,5) and $K\cong${\tt SmallGroup}(32,47). %Therefore $Q/\gamma_Q\cong Q'/\gamma_{Q'}\cong${\tt SmallQuandle}(16,4).
%\end{itemize}

(iii) Note that $G\cong \gamma_2(G)\rtimes_\rho K$. Let $xk\in Z(G)$ with $x\in \Z_p^2$ and $k\in K$. Then $k\in Z(K)\cap \ker{\rho}$ and $x\in Fix(\rho)=\setof{y\in \Z_p^2}{\rho_k(y)=y \text{ for every }k\in K}$. Then $|xk|=|x||k|=2$, therefore $x=1$. Then $Z(G)=\setof{(1,k)}{k\in Z(K)\cap \ker{\rho}}$. Moreover, $K\cong \mathcal{Q}_8\times \Z_2^2$ and $Z(K)=Z(\mathcal{Q}_8)\times \Z_2^2$. Since $\gamma_1(G)\cap Z(G)=1$ we have that $$Z(G)\cap ((\Z_p^2\rtimes \mathcal{Q}_8)\times 1) \leq Z(G)\cap \left(1\times Z(\mathcal{Q}_8)\times 1\right)\leq Z(G)\cap \gamma_1(K)\times 1\leq Z(G)\cap \gamma_1(G)=1.$$ So $G=\left(\Z_p^2\rtimes_\rho \mathcal{Q}_8\times 1\right)\times Z(G)$. Note that $Q/\zeta_Q$ has size $4p$ and is subdirectly reducible, so $p=1\pmod{3}$. Hence $G/Z(G)=G/\dis^\zeta_Q\cong \dis(Q/\zeta_Q)\cong \mathcal{G}_k$.
\end{proof}

%\begin{theorem}\label{p not 1 3 affine}
%Let $p$ be a prime such that $p\neq 1\pmod 3$ and $Q$ be a subdirecly reducible latin quandles of size $16p$. Then $Q$ is directly decomposable and affine.
%\end{theorem}

%\begin{proof}
 %   According to Proposition \ref{group K} the quandle $Q$ is directly decomposable. Latin quandles of size $4,p$ and $16$ are affine and latin quandles of size $4p$ are affine provided $p\neq 1\pmod{3}$. Hence $Q$ is the direct product of affine quandles and so it is also affine.%
%\end{proof}

\begin{theorem}\label{sr quandles}
Let $p=1\pmod{3}$ and $1\neq k\in \Z_p$ such that $k^3=1\pmod{3}$. The representatives of isomorphism classes of subdirectly reducible directly indecomposable latin quandles of size $16p$ are of the form $\mathfrak{Q}(p,j)=\mathcal{Q}(\mathcal{G}_k\times \Z_2^2,Fix(f(j,a)),f(j,a))$ where
   \begin{align}\label{aut of G_k x Klein_0}
        f(j,a)=\begin{cases}
          x\mapsto y a,\\
                  y\mapsto xy ,\\
                  z\mapsto z\\   f|_{\gamma_2(G)}=\begin{bmatrix}
                      -k(1+k^j) & -(1+k^j)\\
                      0 & -k^2(1+k^j)
                  \end{bmatrix},
                  \\
                  f|_{Z(G)}=\begin{bmatrix}
                      0 & 1\\
                      1 & 1
                  \end{bmatrix},
        \end{cases}
    \end{align}
   for $j=1,2$ and $a=(1,0)$.
\end{theorem}

\begin{proof}
Let $Q$ be a subdirectly reducible latin quandles of size $16p$ and $G=\mathcal{G}_k\times \Z_2^2$. According to Proposition \ref{group K} then $\dis(Q)\cong G$ and so $Q\cong \mathcal{Q}(G,Fix(f),f)$ where $f\in \aut{G}$. Since the factor $Q/\zeta_Q\cong Q(p,i)\cong \mathcal{Q}(G/Z(G),Fix(f_{Z(G)}), f_{Z(G)})$ for some $i=1,2$, then $f_{Z(G)}\in \mathcal{F}=\setof{g\in \aut{\mathcal{G}_k}}{|g|=3,\, |Fix(g)|=2p}$, see \cite{LSS}.
Moreover, the isomorphism classes of these quandles are in one-to-one correspondence with the conjugacy classes of automorphisms of $G$. According to Lemma \ref{iso classes of sr} a set of representatives of conjugacy classes of such automorphisms are as in \eqref{aut of G_k x Klein_0} with $a\in \{(0,0),(1,0)\}$. Note that if $a=(0,0)$ the associated quandle decomposes as the direct product of $Q(p,j)$ and $\mathcal{Q}_4$.
\end{proof}

%\begin{proposition}\label{embedding of Q}
 %   Let $Q$ be a subdirectly reducible quandle of size $16p$ and $j:Q\longrightarrow E=E(p,i,f)$ be a subdirect embedding for $i=1,2$. If $Q$ is not directly decomposable then:
  %  \begin{itemize}
   %     \item[(i)] $H=\langle L_{j(x)} L_{j(y)}^{-1}\colon\, x,y\in Q\rangle$ is a normal subgroup of $\lmlt(E)$ and $H\cong \dis(Q)$; 
    %   \item[(ii)]  $[x]_{\mathcal{O}_H}\cong j(Q)$ for every $x\in E$; 
       
     % \item[(iii)] the interval between $0_E$ and $\mathcal{O}_H$ in $\con(E)$ is as follows:
%\begin{align*}
 %   \xymatrixrowsep{0.15in}
%		\xymatrixcolsep{0.15in}
%		\xymatrix{ 
%			& & \mathcal{O}_H  \ar@{-}[d]& &\\
%			&  &\nu_E \ar@{-}[dr] \ar@{-}[dl]&  &\\
 %            & \gamma_{E}\ar@{-}[dr] & &  \mathcal{O}_{Z(H)}\ar@{-}[dl] & \\
%			\alpha_1 \ar@{-}[drr] & \alpha_1\ar@{-}[dr] &\alpha_2\ar@{-}[d] & \ldots\ar@{-}[dl] & \alpha_n \ar@{-}[dll]\\
%			& & 0_E&  &
%		}
 %       \end{align*}		    
        
        
%\item[(iv)]  $\ma(E)=\{\gamma_E,\mathcal{O}_{Z(H)}\}$, $\mathcal{O}_{Z(H)}\leq \zeta_E$ and $[x]_{\zeta_E}\cong \aff(\Z_2^4,f)$. 
 %       \end{itemize}


%\end{proposition}

%    \comment{just need that $Q/\gamma_Q$ is SI of size $16$}

%\begin{proof}

%Let $G=\dis(E)$. 

%(i), (ii) According to Lemma \ref{DD LQGs} we have that $G=\dis(Q(p,i))\times \dis(\aff(\Z_2^4,f))\cong (\Z_p^2\rtimes_\rho \mathcal{Q}_8)\times \Z_2^4$, which has size $128p^2$, and the derived subgroup of $G$ has size $2p^2$. The maps $$H\longrightarrow\dis(Q),\quad h\mapsto h|_{j(Q)}$$
%is an onto morphism of groups and so $H$ has size greater or equal to $32p^2$. So $\gamma_1(H)$ maps onto $\gamma_1(\dis(Q))$ that has size $2p^2$. Since $\gamma_1(H)\leq \gamma_1(G)$ and $2p^2=|\gamma_1(H)|=|\gamma_1(G)|$ it follows that $\gamma_1(G)=\gamma_1(H)\leq H$ and so $H$ is normal in $G$. Moreover, $H$ is stable under conjugation by $L_x$ for $x\in j(Q)$, so $H$ is normal also in $\lmlt(E)$. Thus, $\alpha=\mathcal{O}_H$ is a congruence of $E$ and $[x]_{\alpha}=x^H=j(Q)$. 

%According to \cite[Theorem 2.3]{GB} we have that $\dis(E)_x\leq \gamma_1(G)\leq H$. Therefore
%\begin{align}
    %|H|=|Q||\dis(E)_x|=16p \cdot\frac{|\dis(E)|}{|E|}=16p\cdot\frac{128p^2}{4p\cdot 16} =32p^2=|\dis(Q)|.
%\end{align}
%Hence, $H\cong \dis(Q)$.

% (iii)   The subgroups $Z(H)$ and $\gamma_1(G)$ are normal in $\lmlt(E)$ and they are contained in $H$. So, the orbits of such subgroups provide the congruences $\mathcal{O}_{\gamma_1(G)}=\gamma_E$ and $\mathcal{O}_{Z(H)}$ contained in $\mathcal{O}_H$. The subgroups generated by $\gamma_1(G)$ and $Z(H)$ has size $8p^2$ and its orbits provide a congruence with blocks of size $4p$. The blocks of $\mathcal{O}_H$ have $3$ nontrivial congruences, since they are isomorphic to $Q$ so, according to Lemma \ref{below abelian}(ii), there are no other congruences below $\mathcal{O}_H$. Hence, the interval between $0_E$ and $\mathcal{O}_H$ is as in the diagram. 

 %   Note that $|G_x|={|G|}:{|E|}={128 p^2}:{64p}=2p$ and it is contained in $\gamma_1(G)$ \cite[Theorem 2.3]{GB}. So $G_x\leq H$ and $[H:G_x]=16p$. Thus, $|H|=32p^2=|\dis(Q)|$. Therefore, the map $h\mapsto h|_{j(Q)}$ is an isomorphism between $H$ and $\dis(Q)$.
%\end{proof}


%\comment{If $Q$ has size $16p$ for $p\geq 7$ and it has a congruence with factor of size $p$ then it is abelian.}


%\begin{lemma}\label{minimal cong of DP}
%Let $p$ be a prime with $p=1\pmod{3}$, $Q$ be a subdirectly reducible and directly indecomposable latin quandle of size $16p$, and $E=E(p,i,f)$ for some $i=1,2$. If $j:Q\longrightarrow E$ is a subdirect embedding and $H=\langle L_{j(x)} L_{j(y)}^{-1},\, x,y\in Q\rangle$ then $\ma(E)=\{\gamma_E,\mathcal{O}_{Z(H)}\}$, $\mathcal{O}_{Z(H)}\leq \zeta_E$ and $[x]_{\zeta_E}\cong \aff(\Z_2^4,f)$.
%\end{lemma}
%\comment{together with \ref{embedding of Q}?}
%\begin{proof}

%(iv) Recall that $G=\left(\Z_p^2\rtimes_\rho \mathcal{Q}_8\right)\times \Z_2^4$ and $Z(G)=1\times \Z_2^4$. Therefore $\zeta_Q=\mathcal{O}_{Z(G)}$ is isomorphic to $\aff(\Z_2^4,f)$. 

%Moreover $\gamma_E=\mathcal{O}_{\gamma_1(G)}=\Z_p^2\rtimes_\rho \Z_2   \times 1\leq H$. So $|[x]_{\gamma_E}|=p$ and then $\gamma_E$ is minimal. Moreover the blocks of $\mathcal{O}_{Z(H)}$ has size $4$ and so also $\mathcal{O}_{Z(H)}$ is minimal and it is properly contained in $ \zeta_E$ since $\mathcal{O}_{Z(H)}\wedge \gamma_E=0_E$. Let $\alpha\in \ma(Q)$. Then $\alpha\wedge \gamma_E=0_Q$ since also $\gamma_E$ is minimal. Thus, $\alpha\leq \zeta_E$. The subquandle $[x]_{\zeta_E}$ is $\aff(\Z_2^4,f)$ has a unique proper congruence. Thus $\ma(E)=\{\gamma_E,\mathcal{O}_{Z(H)}\}$. 
%\end{proof}

%\comment{Maybe I need $p>5$.}
%\begin{lemma}
 %   Let $p$ be a prime with $p=1\pmod{3}$ and $Q=E(p,i)$ for $i=1,2$. Then $Q/\zeta_Q$ is a S.I. latin qaundle of size $4p$ with congruence lattice isomorphic to a three elements chain. If $\mu$ is the unique maximal congruence over $\zeta_Q$ then $[x]_\mu$ is an affine latin quandle of size $16p$.
%\end{lemma}


%\begin{proof}
 %   The blocks of the congruence $\mu$ has size $16p$ and it has a congruence with factor of size $p$. According to the previous results then $[x]_\mu$ is affine. In particular $\dis_\mu$ is abelian and so the Sylow subgroups of $\dis_\alpha$ provide trivially intersecting congruences.
%\end{proof}


%\begin{theorem}
 %   Subdirectly reducible quandles of size $16p$ are directly decomposable.
%\end{theorem}

%\begin{proof}
%    Assume that $Q$ is a subdirectly reducible quandle of size $16p$, $j:Q\longrightarrow E=E(p,i,f)$ is a subdirect embedding for $i=1,2$, $G=\dis(Q)$ and $H=\langle L_{j(x)} L_{j(y)}^{-1},\, x,y\in Q\rangle$. Assume also that $Q$ is not directly decomposable. According to Proposition \ref{embedding of Q} the minimal congruences of $E$ are $\mathcal{O}_{Z(H)}$ and $\gamma_E$ and they are contained in $\mathcal{O}_H$. %Moreover, $E\cong E/\zeta_E\times E/\alpha$ for some $\alpha$ such that $E/\alpha=\text{{\tt SmallQuandl(16,4)}}$ such that $\alpha\circ \zeta_{E}=1_E$. 


%Note that $G=\dis(Q(p,i)\times \Z_2^4$ and so $Z(G)=1\times \Z_2^4$. Hence $\zeta_E=\mathcal{O}_{Z(G)}$ has blocks of size $16$. Accordingly we have that $E/\zeta_E=Q(p,i)$ and its congruence lattice is the chain of length $3$. Let us denote by $\mu_{\zeta}$ the unique maximal congruence containing $\zeta_E$. Note that $Q/\mu_\zeta$ has size $4$, and so it is abelian. Thus $\gamma_E,\mathcal{O}_{Z(H)}\leq \mu_\zeta$.  So necessarily $\gamma_E\vee \mathcal{O}_{Z(H)}\leq \mu_\zeta$. Moreover $\mathcal{O}_H\neq \mu_\zeta$, otherwise $\mathcal{O}_H$ would contain more than $3$ proper congruences. So we have the following subset of the congruence lattice of $E$:
      %\begin{center}
		%$\xymatrixrowsep{0.15in}
		%\xymatrixcolsep{0.15in}
		%\xymatrix{ 
	%		& & 1_E  \ar@{-}[dr]^4\ar@{-}[dl]& &\\
	%		& \mathcal{O}_H\ar@{-}[dr] & &\mu_\zeta \ar@{-}[dl] \ar@{-}[dr]^p &\\
         %   &  &\nu_E\ar@{-}[dr] \ar@{-}[dl]&   & \zeta_{E} \ar@{-}[dl]^4 \\
          %   & \gamma_{E}\ar@{-}[dr] & &  \mathcal{O}_{Z(H)}\ar@{-}[dl]^4 & \\
%			\alpha_1 \ar@{-}[drr] & \alpha_1\ar@{-}[dr] &\alpha_2\ar@{-}[d] & \ldots\ar@{-}[dl] & \alpha_n \ar@{-}[dll]\\
		%	& & 0_E&  &
		%}$
%	\end{center}

%The factor $Q'=Q/\mathcal{O}_{Z(H)}$ has size $16p$ and it is not abelian since $\gamma_E$ is not contained in $\mathcal{O}_{Z(H)}$. According to the results given in the previous sections, the congruence lattice of $Q'$ is not one of those allowed for non-abelian latin quandles of size $16p$, since both $|\ma(Q')|$ and $|\Ma(Q')|$ are greater than $2$. Therefore $Q$ is directly decomposable.     %
%    
 %   . The congruence lattice of $E(p,i)$ has at least two maximal congruences, therefore the congruence lattice of $E(p,i)$ looks like this:
  %  \begin{center}
%		$\xymatrixrowsep{0.15in}
%		\xymatrixcolsep{0.15in}
%		\xymatrix{ 
%			& & 1_Q  \ar@{-}[dr]^4\ar@{-}[dl]^4& &\\
%			& \delta\ar@{-}[dr] & &\mu \ar@{-}[dl] \ar@{-}[dr]^p &\\
 %           &  &\mu\cap \delta \ar@{-}[dr] \ar@{-}[dl]&  & \zeta_{E(p,i)} \ar@{-}[dl]^4\\
  %           & \gamma_{E(p,i)}\ar@{-}[dr] & &  \mathcal{O}_{Z(H)}\ar@{-}[dl]^4 & \\
%	%		\alpha_1 \ar@{-}[drr] & \alpha_1\ar@{-}[dr] &\alpha_2\ar@{-}[d] & \ldots\ar@{-}[dl] & \alpha_n \ar@{-}[dll]\\
	%		& & 0_Q&  &
	%	}$
	%\end{center}
%
%Since $E(p,i)\cong E(p,i)/\zeta_Q\times E(p,i)/\alpha$ for some $\alpha$ such that $E(p,i)/\alpha=\text{{\tt SmallQuandl(16,4)}}$ such that $\alpha\circ \zeta_{E(p,i)}=1_Q$. Thus, $\gamma_{E(p,i)}\leq \alpha$ and there exists a unique maximal congruence $\mu_\alpha$ containing $\alpha$. If $\mu\wedge \delta\neq \delta\wedge \mu_\alpha$ then we have 
 %   \begin{center}
%		$\xymatrixrowsep{0.15in}
%		\xymatrixcolsep{0.15in}
%		\xymatrix{ 
%			& & 1_Q \ar@{-}[dll]^4 \ar@{-}[dr]^4\ar@{-}[dl]^4& &\\
%			\mu_\alpha\ar@{-}[d] & \delta\ar@{-}[dr] \ar@{-}[dl]& &\mu \%ar@{-}[dl] \ar@{-}[dr]^p &\\
  %         \mu_\alpha\cap \delta\ar@{-}[dr] &  &\mu\cap \delta \ar@{-}[dr] \ar@{-}[dl]&  & \zeta_{E(p,i)} \ar@{-}[dl]^4\\
   %          & \gamma_{E(p,i)}\ar@{-}[dr] & &  \mathcal{O}_{Z(H)}\ar@{-}[dl]^4 & \\
%			\alpha_1 \ar@{-}[drr] & \alpha_1\ar@{-}[dr] &\alpha_2\ar@{-}[d] & \ldots\ar@{-}[dl] & \alpha_n \ar@{-}[dll]\\
%			& & 0_Q&  &
%		}$
%	\end{center} 
%In this case the blocks of $\delta$ are directly decomposable. So we have that $\mu_\alpha\wedge \delta =\mu\wedge\delta$, i.e. 
 %     \begin{center}
%		$\xymatrixrowsep{0.15in}
%		\xymatrixcolsep{0.15in}
%		\xymatrix{ 
%			& & 1_Q  \ar@{-}[dr]^4\ar@{-}[d]^4\ar@{-}[dl]^4& &\\
%			& \mu_\alpha\ar@{-}[dl]\ar@{-}[dr] & \delta \ar@{-}[d]&\mu \ar@{-}[dl] \ar@{-}[dr]^p &\\
 %         \alpha\ar@{-}[dr]   & &\mu\cap \delta \ar@{-}[dr] \ar@{-}[dl]&  & \zeta_{E(p,i)} \ar@{-}[dl]^4\\
  %           & \gamma_{E(p,i)}\ar@{-}[dr] & &  \mathcal{O}_{Z(H)}\ar@{-}[dl]^4 & \\
%			\alpha_1 \ar@{-}[drr] & \alpha_1\ar@{-}[dr] &\alpha_2\ar@{-}[d] & \ldots\ar@{-}[dl] & \alpha_n \ar@{-}[dll]\\
%			& & 0_Q&  &
%		}$
%	\end{center}
 % 
%\comment{Using that $Q/(\alpha\wedge \beta)\cong Q/\alpha\times Q/\beta$ for maximal congruences of finite latin quandles.}
%
%\comment{the maximal congruences have factor of size $4$, otherwise we get a contradiction. Moreover, if $\mu$ is maximal then $\mu\cap\delta=\mu_\alpha\cap \delta.$ Therefore $Q/(\mu_\alpha\cap \delta)=\mathcal{Q}_4^2$ and there are $5$ maximal congruences. The congruence lattice should look like this:}
%\begin{center}
%		$\xymatrixrowsep{0.15in}
%		\xymatrixcolsep{0.15in}
%		\xymatrix{ 
%			&  & & 1  \ar@{-}[dll]^4\ar@{-}[drr]^4\ar@{-}[dr]^4\ar@{-}[d]^4\ar@{-}[dl]^4& &\\
%			& \mu_\alpha\ar@{-}[dl]\ar@{-}[drr] & \mu_1 \ar@{-}[dr]& \mathcal{O}_H \ar@{-}[d]&\mu_2\ar@{-}[dl] &\mu_{\zeta} \ar@{-}[dll] \ar@{-}[dr]^p &\\
 %         \alpha\ar@{-}[dr]  & &  & \mu_i\cap \mathcal{O}_H \ar@{-}[drr] \ar@{-}[dll]& &  & \zeta_{E(p,i)} \ar@{-}[dl]^4\\
  %           & \gamma_{E(p,i)}\ar@{-}[drr] & & & &  \mathcal{O}_{Z(H)}\ar@{-}[dll]^4 & \\
%			\alpha_1 \ar@{-}[drr] & \alpha_1\ar@{-}[dr] &\alpha_2\ar@{-}[d] & \ldots\ar@{-}[dl] & \alpha_n \ar@{-}[dll]\\
%		&	& & 0&  & &
%		}$
%	\end{center}
%\comment{The blocks of $\mu_\alpha$ have size $16p$ and need to be affine. Indeed they are subdirectly reducible and has more than one maximal congruences. Does this implies that the congruence $\mu_\alpha$ is abelian? Is the join of two abelian congruences also abelian?} 
%
%\comment{The factor $Q/\mathcal{O}_{Z(H)}$ has size $16p$, is subdirectly reducible and has more than one maximal congruence. Since $\gamma_{E(p,i)}$ is not contained in $\mathcal{O}_{Z(H)}$ then $Q/\mathcal{O}_{Z(H)}$ is not abelian. Thus, the congruence lattice of $Q/\mathcal{O}_{Z(H)}$ is not one of the lattices described above.}
%\end{proof}

 
\subsection{Subdirectly irreducible case}

In this section, we will study subdirectly irreducible latin quandles of size $16p$. We start considering quandles for which the congruence lattice is not a chain.

\begin{proposition}\label{SI 16p case 4p}
Let $Q$ be a latin subdirectly irreducible quandle of size $16p$. If $\con(Q)$ is not a chain we have:
\begin{itemize}
    \item[(i)] $\ma_Q=\{\gamma_Q\}$ and $Z(\dis(Q))=1$.
    \item[(ii)] If $|Q/\gamma_Q|=4p$, then $p\in \{3,5,7\}$ and 
    
%    $$
 %   |\dis(Q)| \text{ divides } \begin{cases} 2^{4+k}p \text{ for } 0\leq k\leq 2 \text{ if } p\neq 3,\\
  %  2^{4+k}3^2 \text{ for } 0\leq k\leq 2 \text{ if } p= 3.
   % \end{cases}$$
     $$
    |\dis(Q)| \text{ divides } \begin{cases} 2^{6}p \text{ if } p\neq 3,\\
    2^{6}\cdot 3^2  \text{ if } p= 3.
    \end{cases}$$
    
    \item[(iii)] If $p>7$ then $Q/\gamma_Q\cong \mathcal{Q}_4^2$. 
\end{itemize}
\end{proposition}
 

%\comment{Chi è $\dis(Q)$ quando $Q/\gamma_Q$ ha ordine $4p$? In case (ii) we have a congruence with blocks of size $4p$ that is abelian. Can I use that?}




\begin{proof}
Let $G=\dis(Q)$ and $\alpha=\gamma_Q$. The quandle $Q$ is $2$-step solvable, so according to \cite[Corollary 3.14]{Involutory} the prime divisors of $|G|$ are $2$ and $p$. 


(i) Since $\con(Q)$ is not a chain, we can assume that $\con(Q)$ is as in Proposition \ref{lattices}(i), and then $|\Ma_Q|\geq 2$. By Lemma \ref{mu and gamma} we see that $\mu_Q=\alpha$ is the unique minimal congruence of $Q$. The quandle $Q$ is not nilpotent, so according to Lemma \ref{dis_gamma}(ii) $Z(G)=1$.


(ii) Assume that $Q/\alpha$ has size $4p$. Then the blocks of $\gamma_Q$ have size $4$, and so the blocks are isomorphic to $\mathcal{Q}_4$ (the unique latin quandle of size $4$). The quandle $Q/\alpha\cong \mathcal{Q}_4\times \aff(\Z_p,k)$ is generated by $(a,0),(0,b)$ for every $0\neq a\in \Z_2^2$ and $0\neq b\neq \Z_p$. According to Corollary \ref{embedding as quandle} we have that $\dis_{\alpha}$ embeds into $\Z_2^{2+k}$ for $0\leq k\leq 2$ and that the group $\dis^{\alpha}=\gamma_1(G)$ embeds into $(\aut{[x]_{\alpha}})^2\cong (\Z_2^2\rtimes \Z_3)^2$ which $3$-Sylow subgroups are isomorphic to $\Z_3^3$. Furthermore, $G/\gamma_1(G)=\dis(Q/\alpha)\cong \Z_2^2\times \Z_p$. We have to discuss two cases.





%
%According to Lemma \ref{dis_gamma} $\gamma_2(G)=\dis_{\gamma_Q}$ and it embeds into $(\Z_2^2)^2$ since \comment{$Q/\gamma_Q\cong \aff(\Z_2^2,f)\times \aff(\Z_p,k)$ is generated by $(a,0),(0,b)$ for every $0\neq a\in \Z_2^2$ and $0\neq b\neq \Z_p$ - move?} \cite{}. According to... $E=Q_{\gamma_2(G)}$ is a connected cover of $Q/\gamma_Q$ that is affine. So $rad(|E|)=rad(|Q/\gamma_Q|)=2p$ according to... . Therefore 
%$$|G|=|G/\gamma_2(G)||\gamma_2(G)|=|E||\gamma_2(G)|,$$
%and so $rad(|G|)=2p$.
%
%Since $Q/\gamma_Q\cong \aff(\Z_2^2,f)\times \aff(\Z_p,k)$ is generated by $(a,0),(0,b)$ for every $0\neq a\in \Z_2^2$ and $0\neq b\neq \Z_p$ the group $\gamma_1(G)$ embeds into $(\aut{[x]_{\gamma_Q}})^2\cong (\Z_2^2\rtimes \Z_3)^2$. 

\begin{itemize}
    \item Let $p>3$. Then $\gamma_1(G)$ embeds into $(\Z_2^2)^2$ and $|G|=|G/\gamma_1(G)||\gamma_1(G)|=|Q/\alpha||\gamma_1(G)|=4p\cdot 2^t $ for some $2\leq t\leq 4$. Let $x\in G$ be an element of order $p$ and let $H=\langle \gamma_1(G),x\rangle$. Note that $H$ is normal in $G$ and therefore contains the $p$-Sylow subgroups of $G$. Moreover, $H$ is a characteristic subgroup, so in particular $H$ is normal in $\lmlt(Q)$. If $H$ is abelian, then its Sylow subgroups provide trivially intersecting congruences, according to Lemma \ref{Sylow}. So, $H\cong \gamma_1(G)\rtimes_\rho \Z_p$, and accordingly $p$ divides $GL_{2+k}(2)$ where $0\leq k\leq 2$. Thus, $p\in \{3,5,7\}$.

\item Let $p=3$. According to Lemma \ref{N cyclic} the $3$-Sylow of $\gamma_1(G)$ is cyclic and so isomorphic to $\Z_3^s$ for $s=0,1$, so $|\gamma_1(G)|$ divides $2^{2+k}3^s$ for $0\leq k\leq 2$ and $0\leq s\leq 1$. Thus $|G|$ divides $2^{6}3^2$.
\end{itemize}


(iii) We have that $\mu_Q=\alpha$ and it is the unique minimal congruence of $Q$ (see item (i)). The factor $Q/\alpha$ is not simple. Quandles of size $4$ and $p$ are simple, so according to Lemma \ref{meet of max is min}(i) $Q/\alpha$ has size $4p$ or $16$. According to item (ii), then $Q/\alpha$ have size $16$ and is subdirectly reducible. So $Q/\alpha$ is the unique subdirectly reducible latin quandle of size $16$, that is, $\mathcal{Q}_4^2$ (latin quandles of size $16$ are classified in \cite{RIG}).
%
%In the first case, $[x]_{\alpha_i}$ has size $4p$ for every  $i$ and it not abelian by Lemma \ref{Sylow}. Then necessarily $[x]_{\alpha_i}\in \LSS(p,4,\sigma_Q)$ and so $p=1\pmod{3}$.
%In the first case, we can assume that $|Q/\alpha_1|=4$. Then $[x]_{\alpha_1}$ has size $4p$ and it not abelian by Lemma \ref{Sylow}. Then necessarily $[x]_{\alpha_1}\in \LSS(4,p,\sigma_Q)$ and so according to \cite{LSS} $p=7$ and $[x]_{\alpha_1}$ is not latin. Thus, $Q/\gamma_Q$ has size $16$, it is subdirectly reducible, and so $Q/\gamma_Q\cong \mathcal{Q}_4\times \mathcal{Q}_4$ (latin quandles of size $16$ are classified). Accordingly, $Q/\gamma_Q$ has $5$ congruences.
\end{proof}

%\begin{proposition}\label{SI 16p}
%Let $Q$ be a subdirectly irreducible latin quandle of size $16p$ for $p>7$. If $\con(Q)$ is not a chain then %we have:
%\begin{itemize}
    %\item[(i)] 
%    $Q/\gamma_Q\cong \Q_4\times \Q_4$. % and $\con(Q)$ is	
%	\begin{center}
%		$\xymatrixrowsep{0.15in}
%		\xymatrixcolsep{0.15in}
%		\xymatrix{ 
%			& & 1_Q  \ar@{-}[d]\ar@{-}[dr]\ar@{-}[drr] \ar@{-}[dl] \ar@{-}[dll] & &\\
%			\alpha_1 \ar@{-}[drr] & \alpha_2\ar@{-}[dr] &\alpha_3\ar@{-}[d] & \alpha_4\ar@{-}[dl] & \alpha_5 \ar@{-}[dll]\\
%			& & \gamma_Q\ar@{-}[d] &  &\\
%			& & 0_Q & &
%		}$
%		\end{center}
%\item[(ii)] $Q/\gamma_Q\cong \Q_4\times \Q_4$;
%\item[(iii)] $[x]_{\alpha_i}\in \LSS(p,4,\sigma_Q)$ and $p=1\pmod{3}$. \comment{needed?}
%\item[(iv)] $Q$ is $2$-step solvable and $\zeta_Q=0_Q$  \comment{needed? Using $Z(\dis(Q))=1$}
%\item[(ii)] $Z(\dis(Q))=1$. 
%\end{itemize}




%one of the following holds:

%\begin{itemize}

%\item[(i)] the congruence lattice of $Q$ is	
%	\begin{center}
		%$\xymatrixrowsep{0.15in}
		%\xymatrixcolsep{0.15in}
	%	\xymatrix{ 
	%		& & 1_Q  \ar@{-}[d]\ar@{-}[dr]\ar@{-}[drr] \ar@{-}[dl] \ar@{-}[dll] & &\\
	%		\alpha_1 \ar@{-}[drr] & \alpha_2\ar@{-}[dr] &\alpha_3\ar@{-}[d] & \alpha_4\ar@{-}[dl] & \alpha_5 \ar@{-}[dll]\\
	%		& & \gamma_Q\ar@{-}[d] &  &\\
	%		& & 0_Q & &
	%	}$
	%	\end{center}
%Moreover, $Q/\gamma_Q\cong \Q_4\times \Q_4$, $[x]_{\alpha_i}\in \LSS(p,4,\sigma_Q)$ and $p=1\pmod{3}$.
%
%\item[(ii)] $p=7$ and the congruence lattice of $Q$ is	
%	\begin{center}		
%		$
		%\xymatrixrowsep{0.15in}
		%\xymatrixcolsep{0.15in}
	%	\xymatrix{ 
	%		& 1_Q \ar@{-}[dr]\ar@{-}[dl] & \\
%			\alpha_1 \ar@{-}[dr] & &\alpha_2\ar@{-}[dl] \\
%			& \gamma_Q \ar@{-}[d] & \\
%			& 0_Q & 
%		}
%		$
%	\end{center}
%and $Q/\gamma_Q\cong \Q_4\times \aff(\mathbb{Z}_7,f)$. 
%\end{itemize}

%\end{proposition}
%\begin{itemize}




%\item[(1)] \comment{$\mathcal{Q}_4^2$ is $2$-generated? No! $2$ generated subquandles are all iso to $\mathcal{Q}_4$.}


%\end{itemize}


%\begin{proof}
%(i) 
%By Lemma \ref{SI 16p case 4p}(i) we have that $\mu_Q=\gamma_Q$ is the unique minimal congruence of $Q$. The factor $Q/\gamma_Q$ is not simple. Quandles of size $4$ and $p$ are simple, so according to Lemma \ref{meet of max is min}(i) $Q/\gamma_Q$ has either size $4p$ or size $16$. According to Lemma \ref{SI 16p case 4p}(ii) then $Q/\gamma_Q$ have size $16$ and it is subdirectly reducible. So $Q/\gamma_Q\cong \mathcal{Q}_4\times \mathcal{Q}_4$ (latin quandles of size $16$ are classified). %Accordingly, $Q/\gamma_Q$ has $5$ congruences.
%
%
%In the first case, $[x]_{\alpha_i}$ has size $4p$ for every  $i$ and it not abelian by Lemma \ref{Sylow}. Then necessarily $[x]_{\alpha_i}\in \LSS(p,4,\sigma_Q)$ and so $p=1\pmod{3}$.
%
%In the first case, we can assume that $|Q/\alpha_1|=4$. Then $[x]_{\alpha_1}$ has size $4p$ and it not abelian by Lemma \ref{Sylow}. Then necessarily $[x]_{\alpha_1}\in \LSS(4,p,\sigma_Q)$ and so according to \cite{LSS} $p=7$ and $[x]_{\alpha_1}$ is not latin. Thus, $Q/\gamma_Q$ has size $16$, it is subdirectly reducible, and so $Q/\gamma_Q\cong \mathcal{Q}_4\times \mathcal{Q}_4$ (latin quandles of size $16$ are classified). Accordingly, $Q/\gamma_Q$ has $5$ congruences.
%
%The congruence $\gamma_Q$ is abelian since it is minimal and so $Q$ is $2$-step solvable. 
%
%(ii) The congruence $\gamma_Q$ is the unique minimal congruence of $Q$ that is not nilpotent. Then according to Lemma \ref{dis_gamma} $Z(\dis(Q))=1$.
%
%
%Let $\alpha,\beta$ be maximal congruences of $Q$ and let $\alpha\neq \beta$. According to Lemma \ref{LSS 16p} and Lemma \ref{meet of max is min}(ii) either both factors have size $4$ or one has size $p$ and one has size $4$.
%
%Let $\alpha,\beta,\delta$ be maximal congruences. If $\alpha\wedge	\beta\neq \alpha\wedge\delta$ then $\alpha\wedge\beta\wedge\delta=0_Q$ since $\alpha\wedge\delta$ is minimal according to Lemma \ref{meet of max is min}. Then $Q$ embeds into $Q/(\alpha\wedge\beta)\times Q/\delta$ which is affine. Hence the meet of any pair of maximal congruence is the same and the factor is affine, by virtue of \comment{cite}. Let $\nu$ be such congruence, then $\gamma_Q\leq \nu$ and since $\gamma_Q\neq 0_Q$ and $\nu$ is minimal, the equality holds. Moreover $Q/\alpha\times Q/\beta\cong Q/\gamma_Q$ for every pair of maximal congruences $\alpha,\beta$, therefore either $Q/\gamma_Q$ is the direct product of $\Q_4$ and a quandle of size $p$ or all the simple factors of $Q$ are isomorphic to $\Q_4$ and $Q/\gamma_Q\cong \Q_4\times \Q_4$.
%\end{proof}



%
%
%\begin{corollary}\label{SI factor}
%Let $Q$ be a non-affine latin quandle of size $16p$ and $\alpha\in \ma_Q$. If $\alpha\neq \gamma_Q$ then $Q/\alpha$ is S.I.
%%If $Q/\alpha$ is not SI then $\alpha=\gamma_Q$.
%\end{corollary}
%\comment{This needs just height $3$ and solvable (maybe)}
%
%\begin{proof}
%Assume that $Q/\alpha$ is not S.I. Since the height of $\con(Q/\alpha)$ is $2$, then there exists $\beta_1,\beta_2\in \Ma_Q$ such that $\alpha\leq \beta_1\wedge\beta_2=\gamma_Q$. Hence, the equality holds.
%\end{proof}
%%
%\subsection{Subdirectly irreducible case}
%
%
%\begin{proposition}
%Let $Q$ be a non-affine latin subdirectly irreducible quandle of size $16p$. Then one of the following holds:
%
%\item[(i)] $|\Ma_Q|=1$ and the congruence lattice of $Q$ is 
%
%	\begin{center}
%		$ \xymatrixrowsep{0.15in}
%		\xymatrixcolsep{0.15in}
%		\xymatrix{ 		
%			 1_Q  \ar@{-}[d]\ar@{-}\\
%			\gamma_Q\ar@{-}[d]\\
%			 0_Q 
%		}$
%		\end{center}
%		
%and if $|Q/\gamma_Q|=p$ then $p\in \{3,5,7,17,31,127\}$, otherwise $Q\in \LSS(p,16,0_Q)$. 
%
%\item[(ii)] $|\Ma_Q|=1$ and $\con(Q)$ has the form
%
%	
%	\begin{center}
%		$ \xymatrixrowsep{0.15in}
%		\xymatrixcolsep{0.15in}
%%		\xymatrix{ 		
%%			 1_Q  \ar@{-}[d]\ar@{-}\\
%%			\gamma_Q\ar@{-}[d]\\
%%			 0_Q 
%%		}
%%		\qquad\qquad
%		\xymatrixrowsep{0.15in}
%		\xymatrixcolsep{0.15in}
%		\xymatrix{ 		
%			 1_Q  \ar@{-}[d]\ar@{-}\\
%			\alpha \ar@{-}[d]\\
%			 \beta\ar@{-}[d]\\\
%			 0_Q 
%		}
%		$
%	\end{center}
%	
%
%\item[(iii)] The congruence lattice of $Q$ is	
%	\begin{center}
%		$\xymatrixrowsep{0.15in}
%		\xymatrixcolsep{0.15in}
%		\xymatrix{ 
%			& & 1_Q  \ar@{-}[d]\ar@{-}[dr]\ar@{-}[drr] \ar@{-}[dl] \ar@{-}[dll] & &\\
%			\alpha_1 \ar@{-}[drr] & \alpha_2\ar@{-}[dr] &\alpha_3\ar@{-}[d] & \alpha_4\ar@{-}[dl] & \alpha_5 \ar@{-}[dll]\\
%			& & \gamma_Q\ar@{-}[d] &  &\\
%			& & 0_Q & &
%		}$
%		\end{center}
%and $Q/\gamma_Q\cong \Q_4\times \Q_4$ and $\zeta_Q=0_Q$.
%
%\item[(iv)] The congruence lattice of $Q$ is	
%	\begin{center}		
%		$
%		\xymatrixrowsep{0.15in}
%		\xymatrixcolsep{0.15in}
%		\xymatrix{ 
%			& 1_Q \ar@{-}[dr]\ar@{-}[dl] & \\
%			\alpha_1 \ar@{-}[dr] & &\alpha_2\ar@{-}[dl] \\
%			& \gamma_Q \ar@{-}[d] & \\
%			& 0_Q & 
%		}
%		$
%	\end{center}
%and $Q/\gamma_Q\cong \Q_4\times \aff(\mathbb{Z}_p,f)$ and $\zeta_Q=0_Q$.
%\end{proposition}
%
%\begin{proof}
%(i) - (ii) Assume that $|\Ma_Q|\geq 1$.Let $\alpha$ be the unique maximal congruence and let $\beta$ be the unique minimal congruence of $Q$. If $\alpha=\beta$ we have the first lattice, since $Q$ is solvable but not abelian. In particular, both $Q/\gamma_!$ and the blocks of $\gamma_Q$ have prime power order. So either $|Q/\alpha|=16$ and $|[a]|=p$ or the other way around. In the first case, we can apply Lemma \ref{LSS 16p} and so $Q\in \LSS(p,16,0_Q)$.
%
%
%In the second case, let $G=\dis(Q)$: we have that 
%$$\gamma_2(G)\leq \dis_{\gamma_Q}\leq \gamma_1(G)$$
%The group $G$ is not nilpotent and $\dis_\alpha$ is a minimal subgroup of $\N(Q)$, hence $\dis_\alpha=\gamma_2(G)$. Moreover $G/\gamma_1(G)$ is cyclic and $\gamma_1(G)/\gamma_2(G)$ is generated by the elements $\setof{[x,y]\gamma_2(G)}{x,y\in X}$ and $X$ generates $G/\gamma_1(G)$ \cite[Proposition 9.2.5]{Sims}, thus $\gamma_1(G)=\dis_\alpha$. According to Lemma \ref{minimal cong}, $\dis_\alpha$ is an elementary abelian $2$-group and so the short exact sequence
%$$\dis_\alpha\longrightarrow G	\longrightarrow \Z_p$$
%splits. By \cite[Lemma 1.4(ii)]{LSS} $\dis_\alpha$ embeds into $\Z_2^{4+k}$ for $k\leq 4$ and so $p$ divides $|GL_2(4+k)|$. Therefore $p\in \{3,5,7,17,31,127\}$. 
%\comment{Is $Q\in\LSS(16,p,\sigma_Q)$? Can I say something more?} \comment{Yes! $Q$ is $2$-step solvable}
%
%If $\alpha\neq \beta$, we have the congruence lattice displayed in (ii), since the height of $\con(Q)$ is at most $3$ by Lemma \ref{meet of max is min}(i).
%
%
%(ii) - (iii) Assume that $|\Ma_Q|\geq 2$. By Lemma \ref{meet of max is min}(iv) we have that $\mu_Q=\gamma_Q$ is the unique minimal congruence of $Q$. If $Q/\alpha$ has size $16$ then we are is case (ii), otherwise $Q/\gamma_Q$ has size $4p$ and we have case (iii).
%
%In both cases, if $\zeta_Q\neq 0_Q$ then $\gamma_Q\leq \zeta_Q$ and so $Q$ is nilpotent and so abelian by Lemma \ref{16p nilpotent are abelian}, contradiction.
%%
%%
%%Let $\alpha,\beta$ be maximal congruences of $Q$ and let $\alpha\neq \beta$. According to Lemma \ref{LSS 16p} and Lemma \ref{meet of max is min}(ii) either both factors have size $4$ or one has size $p$ and one has size $4$.
%%
%%Let $\alpha,\beta,\delta$ be maximal congruences. If $\alpha\wedge	\beta\neq \alpha\wedge\delta$ then $\alpha\wedge\beta\wedge\delta=0_Q$ since $\alpha\wedge\delta$ is minimal according to Lemma \ref{meet of max is min}. Then $Q$ embeds into $Q/(\alpha\wedge\beta)\times Q/\delta$ which is affine. Hence the meet of any pair of maximal congruence is the same and the factor is affine, by virtue of \comment{cite}. Let $\nu$ be such congruence, then $\gamma_Q\leq \nu$ and since $\gamma_Q\neq 0_Q$ and $\nu$ is minimal, the equality holds. Moreover $Q/\alpha\times Q/\beta\cong Q/\gamma_Q$ for every pair of maximal congruences $\alpha,\beta$, therefore either $Q/\gamma_Q$ is the direct product of $\Q_4$ and a quandle of size $p$ or all the simple factors of $Q$ are isomorphic to $\Q_4$ and $Q/\gamma_Q\cong \Q_4\times \Q_4$.
%\end{proof}

%
%
%\begin{corollary}
%	Let $Q$ be a subdirectly irreducible quandle of size $16p$ and let $n>1$. Then $\con(Q)$ is one of the following:
%	
%	\medskip
%	
%	\begin{center}
%		$\xymatrixrowsep{0.15in}
%		\xymatrixcolsep{0.15in}
%		\xymatrix{ 
%			& & 1_Q  \ar@{-}[d]\ar@{-}[dr]\ar@{-}[drr] \ar@{-}[dl] \ar@{-}[dll] & &\\
%			\alpha_1 \ar@{-}[drr] & \alpha_2\ar@{-}[dr] &\alpha_3\ar@{-}[d] & \alpha_4\ar@{-}[dl] & \alpha_5 \ar@{-}[dll]\\
%			& & \gamma_Q\ar@{-}[d] &  &\\
%			& & 0_Q & &
%		}
%		\qquad \xymatrixrowsep{0.15in}
%		\xymatrixcolsep{0.15in}
%		\xymatrix{ 
%			& 1_Q \ar@{-}[dr]\ar@{-}[dl] & \\
%			\alpha_1 \ar@{-}[dr] & &\alpha_2\ar@{-}[dl] \\
%			& \gamma_Q \ar@{-}[d] & \\
%			& 0_Q & 
%		}
%		$
%	\end{center}
%\end{corollary}


%\subsection{Subdirectly reducible case}
%%
%%\begin{lemma}
%%Let $Q$ be a subdirectly reducible non-affine quandle of size $16p$. Then one of the following holds:
%%\begin{itemize}
%%\item[(i)] $Q\cong \Q_4\times \widetilde{Q}$ where $\widetilde{Q}$ is a non-affine S.I, quandle of size $4p$.
%%\item[(ii)] $|Q/\alpha|\in \{16,4p\}$ for every $\alpha\in \ma_Q$.
%%\end{itemize}
%%
%%\end{lemma}
%%
%%\begin{proof}
%%Let $\alpha\in \ma_Q$. If $|Q/\alpha|\in \{4,p\}$ then $\alpha$ is both a minimal and a maximal congruence. Then $\alpha\wedge \beta=0_Q$ and $\alpha\vee \beta=1_Q$ for every $\beta\in \ma_Q$, i.e. $Q$ is a direct product. Therefore, if $Q$ is not a direct product $|Q/\alpha|\in \{16,4p\}$. 
%%\end{proof}
%%
%
%
%\begin{lemma}\label{sub red max and min cong}
%Let $Q$ be a subdirectly reducible non-affine quandle of size $16p$. Then one of the following holds:
%\begin{itemize}
%\item[(i)] $Q\cong \Q_4\times \widetilde{Q}$ where $\widetilde{Q}$ is a non-affine S.I, quandle of size $4p$.
%\item[(ii)] $|\Ma_Q|=1$ and $|Q/\alpha|\in \{16,4p\}$ for every $\alpha\in \ma_Q$.
%\end{itemize}
%
%\end{lemma}
%\comment{NEEDED?}
%\begin{proof}
%Let $\alpha\in \ma_Q$. If $|Q/\alpha|\in \{4,p\}$ then $\alpha$ is both a minimal and a maximal congruence. Then $\alpha\wedge \beta=0_Q$ and $\alpha\vee \beta=1_Q$ for every $\beta\in \ma_Q$, i.e. $Q$ is a direct product. Therefore, if $Q$ is not a direct product $|Q/\alpha|\in \{16,4p\}$. 
%
%Assume that $|\Ma_Q|\geq 2$, and let $\alpha_1,\alpha_2\in \Ma_Q$ with factor respectively of size $4$ and $p$ (they can not be both of size $p$ and for case $16$ see Lemma \ref{LSS 16p}). According to Lemma \ref{gamma and max} we have that $\alpha_1\wedge\alpha_2=\gamma_Q\in \ma_Q$. Let $\gamma_Q\neq \beta\in \ma_Q$. Then $Q/\beta$ is a non-affine latin quandle of size $4p$ and so either $\beta\leq \alpha_i$ or $Q$ embeds in $Q/\alpha_i\times Q/\beta$ and accordingly $|Q|$ divides $|Q/\alpha_i\times Q/\beta|=4p|Q/\alpha_i|$. If $i=1$ then $Q$ is a direct product, otherwise $|Q|$ divides $4p^2$, contradiction.
%\end{proof}
%
%
%\begin{proposition}\label{Sub red}
%Let $Q$ be a subdirectly reducible non-affine quandle of size $16p$. Then one of the following holds:
%\begin{itemize}
%\item[(i)] $Q\cong \Q_4\times \widetilde{Q}$ where $\widetilde{Q}$ is a non-affine S.I, quandle of size $4p$.
%
%\item[(ii)] The congruence lattice of $Q$ has the form 
%	\begin{center}
%$\xymatrixrowsep{0.15in}
%		\xymatrixcolsep{0.15in}
%		\xymatrix{ 
%			 & 1_Q  \ar@{-}[d]& &\\
%			 & \delta \ar@{-}[dr]\ar@{-}[dl] & &\\
%			\gamma_Q \ar@{-}[dr]  & & \alpha\ar@{-}[dl] \\
%			 & 0_Q&  &
%		}$
%%$\xymatrixrowsep{0.15in}
%%		\xymatrixcolsep{0.15in}
%%		\xymatrix{ 
%%			& & 1_Q  \ar@{-}[d]& &\\
%%			& & \delta  \ar@{-}[d]\ar@{-}[dr]\ar@{-}[drr] \ar@{-}[dl] \ar@{-}[dll] & &\\
%%			\gamma_Q \ar@{-}[drr] & \alpha_1\ar@{-}[dr]  &\alpha_2\ar@{-}[d]  & \ldots\ar@{-}[dl]  & \alpha_n \ar@{-}[dll] \\
%%			& & 0_Q&  &
%%		}$
%	\end{center}
%and $Q/\gamma_Q$ is a S.I. latin quandle of size $16$ and $Q/\alpha$ is a S.I. latin quandles of size $4p$. 
%\item[(iii)] $p=7$ and the congruence lattice of $Q$ has the form 
%\begin{center}
%		$\xymatrixrowsep{0.15in}
%		\xymatrixcolsep{0.15in}
%		\xymatrix{ 
%			& & 1_Q  \ar@{-}[d]& &\\
%			& & \gamma_Q  \ar@{-}[d]\ar@{-}[dr]\ar@{-}[drr] \ar@{-}[dl] \ar@{-}[dll] & &\\
%	\alpha_1 \ar@{-}[drr] & \alpha_2\ar@{-}[dr] &\alpha_3\ar@{-}[d] &  \alpha_4\ar@{-}[dl] & \alpha_5 \ar@{-}[dll]\\
%%			\alpha_1 \ar@{-}[drr] & \alpha_1\ar@{-}[dr] &\alpha_2\ar@{-}[d] & \ldots\ar@{-}[dl] & \alpha_n \ar@{-}[dll]\\
%			& & 0_Q&  &
%		}$
%	\end{center}
%and $Q/\alpha_i$ is a S.I. latin quandles of size $28$ for every $\alpha_i\in \ma_Q$.
%
%
%\end{itemize}
%\end{proposition}
%
%\comment{Can I say something more? Yes, $Q$ is solvable of length $2$.}
%
%\begin{proof}
%According to Lemma \ref{sub red max and min cong} if $Q$ is not a direct product, then $|\Ma_Q|=1$ and the factors with respect to minimal congruences have size $16$ or $4p$.
%
%
%Assume that there exists $\alpha\in \ma_Q$ such that $|Q/\alpha|=16$. Since $Q$ is not affine and latin quandles of size $16$ are affine, then $\alpha$ is the unique congruence with such properties (otherwise if $\beta$ has the same properties, $Q$ embeds into $Q/\alpha\times Q/\beta$ and so $Q$ is affine) and $\alpha=\gamma_Q$. Moreover,
%
%$$\dis_\delta=\dis_{\gamma_Q}\times \dis_{\alpha_i}\cong \Z_p^m\times \Z_2^n\cong \dis_{\alpha_i}\times \dis_{\alpha_j}\cong \Z_2^r$$
%
%for every $\alpha_i\alpha_j\in \ma_Q$ and so, since $p>2$ we have that $Q$ has $2$ minimal congrueces, i.e. we are in case (ii).
%
%If all the factors with respect the minimal congruences have size $4p$, they are S.I. and so the minimal congruences have blocks of size $4$. Hence 
%$$\dis_{\gamma_Q}\cong \dis_{\alpha_i}\times \dis_{\alpha_j}\cong \Z_2^n$$
%and so $\gamma_Q$ have blocks of size $16$. Namely $|Q/\gamma_Q|$ and so according to \cite[]{LSS} $p=7$. The minimal congruences of $Q$ induce congruences of $[a]_{\gamma_Q}$ and therefore, according to GAP, the unique latin quandle of size $16$ with more than $1$ congruence is $\Q_4\times \Q_4$ and it has $5$ nontrivial congruences.
%\end{proof}

%
%
%Is $\con(Q)$ is one of the following?
%
%\begin{center}
%		$\xymatrixrowsep{0.15in}
%		\xymatrixcolsep{0.15in}
%		\xymatrix{ 
%		&	& & 1_Q \ar@{-}[dlll] \ar@{-}[d]\ar@{-}[dr]\ar@{-}[drr] \ar@{-}[dl] \ar@{-}[dll] & &\\
%	\beta\ar@{-}[drrr] &		\delta_1 \ar@{-}[drr] & \delta_2\ar@{-}[dr] &\delta_3\ar@{-}[d] & \delta_4\ar@{-}[dl] & \delta_5 \ar@{-}[dll]\ar@{-}[d]\\
%	&		& & \alpha_1=\gamma_Q\ar@{-}[d] & & \alpha_2\ar@{-}[dll] &\\
%		&	& & 0_Q & &
%		}
%		\qquad \xymatrixrowsep{0.15in}
%		\xymatrixcolsep{0.15in}
%		\xymatrix{ 
%			& 1_Q \ar@{-}[dr]\ar@{-}[dl] & \\
%			\beta \ar@{-}[dr] & &\delta \ar@{-}[dl] \ar@{-}[dr]\\
%			& \alpha_1=\gamma_Q \ar@{-}[dr] & & \alpha_2 \ar@{-}[dl] \\
%			&  & 0_Q & 
%		}
%		$
%	\end{center}





According to Proposition \ref{SI 16p case 4p} if $p\leq 7$ the size of the displacement group is bounded, so we can study such quandles computationally in the next section. Let us proceed with the case $p>7$.% First, let us identify the congruence lattice.

%\comment{Let us proceed by identifying the displacement group of subdirectly irreducible latin quandles of size $16p$ such that the congruence lattice is not a chain, provided $p>7$.}

%\comment{In SmallGroup(32,49) ci sono sottogruppi $\Z_2^3$. Ristretta a questi sottogruppi, $\rho$ non può essere fedele (sono matrici diagonali di ordine $2$).}


%\begin{proposition}\label{chain}
 %   Let $Q$ be a subdirectly irreducible latin quandle of size $16p$ for $p>7$. Then $\con(Q)$ is a chain.
%\end{proposition}

%\begin{proposition}\label{Dis for SI}
  %  Let $Q$ be a subdirectly irreducible latin quandle of size $16p$ for $p>7$ such that $\con(Q)$ is not a chain. Then
   % $$\dis(Q)=\begin{cases}
    %    \Z_p^3\rtimes_{\rho} \Z_2^4\\
     %   \Z_p^3\rtimes_\rho K
    %\end{cases}$$
    %where $\rho$ is an irreducible faithful action and $K$=SmallGroup(32,49). In the second case $Z(\gamma_1(\dis(Q)))=1$.
%$$\dis(Q)=\Z_p^2\rtimes_{\rho} K$$
%where $\rho$ is a faithful action, $\gamma_1(G)=\Z_p^2\rtimes_\rho Z(K)$, $K=SmallGroup(32,49)$ and $Z(\dis(Q))=Z(\gamma_1(\dis(Q)))=1$.
%\end{proposition}

%\comment{si può fare forse a meno di invocare la classificazione dei quandles of size $32$, se si riesce a dire che $K$ è extraspeciale in un altro modo}

%\comment{Usare notazione $K_{i}?$}

%\comment{the center of SmallGroup(32,47) is $\Z_2^3$. Hence the action on $\Z_2^2$ is not faithful. Then $t=3$. Usando \ref{center = -1} $t=2$, contradiction!}

%\begin{proof}
%Let $G=\dis(Q)$ and $\alpha=\gamma_Q$. Assume that $\con(Q)$ is not a chain. According to Lemma \ref{SI 16p case 4p}, $\alpha$ is the unique minimal congruence of $Q$, the factor $Q/\alpha$ is $\mathcal{Q}_4^2$ and then $G/\gamma_1(G)\cong \Z_2^4$. Moreover, $Q$ is not nilpotent, so $\gamma_2(G)=\dis_\alpha$ and $Z(G)=1$, according to Lemma \ref{dis_gamma}. The blocks of $\alpha$ have size $p$, the quandle $Q/\alpha$ is generated by $3$ elements, so the group $\dis_\alpha$ embeds into $\Z_p^3$ according to Corollary \ref{embedding as quandle}(ii). Moreover $|\dis_\alpha|\neq p$ since $\alpha$ is not central (see Lemma \ref{central 2}).



%We have to distinguish between two cases.

%The congruence $\alpha$ is the unique minimal congruence of $Q$ that is not nilpotent. Then according to Lemma \ref{dis_gamma} $Z(G)=1$. %The action $\rho$ is faithful since $Z(G)=1$ (see Lemma \ref{no center}).

%\begin{itemize}
 %   \item Assume that $s=0$. Then $\gamma_1(G)=\gamma_2(G)$ and so $G=\Z_p^t\rtimes_\rho \Z_2^4$ for $t\leq 3$. By Lemma \ref{ro faithful on Z_p}, $\rho$ is not faithful since $t\leq 3\leq 4$, leading to a contradiction. 
 

%\item Assume that $s=1$. The quandle $Q'$ is a connected quandle of size $32$ with factor $Q'/\gamma_{Q'}$ isomorphic $\mathcal{Q}_4^2$, having $7$ congruences. Then, looking at Table \ref{32}, we have that $Q'$ is {\tt SmallQuandle(32,j)} for $j=2,3$ and accordingly $K=G/\gamma_2(G)=\dis(Q')$ is {\tt SmallGroup(32,i)} for $i=47,49$. So $G=\Z_p^t\rtimes_\rho K$ for $t=2,3$. Both groups have a subgroup isomorphic to $\Z_2^3$, so we have $t=3$ by Corollary \ref{ro faithful on Z_p}. According to Corollary \ref{center of dis^alpha} then $Z=Z(\gamma_1(G))=Z(\dis^{\gamma_Q})=1$ since it is not a group of size a power of a prime. On the other hand we have that $Z\neq 1$ according to Corollary \ref{center = -1}, a contradiction. %we have $Z=Z(\gamma_1(G))\neq 1$. Since $\beta=\mathcal{O}_Z\leq \alpha$, then we have $\beta=\alpha$ and so $\dis_\alpha\leq Z$. Consequently, $\gamma_1(G)$ is abelian and has size $2p$. By Lemma \ref{Sylow} the Sylow subgroups of $\gamma_1(G)$ provide trivially intersecting congruences, leading to a contradiction.  \qedhere
%
%So the elements of $\gamma_1(K)$ commute with the elements of $K$ and the elements of $\dis_\alpha$. Hence $Z(K)\leq Z(G)=1$, contradiction. \qedhere
%\end{itemize}

%
%
%According to Proposition \ref{no reps} there is no such groups for $t=2$. Thus, $t=3$. 
%
%Finally, note that $\rho$ is irreducible according to Lemma \ref{irreducible2} and Lemma \ref{irreducible}. According to \ref{no irred Z_2^n} there is no irreducible representation of $\Z_2^4$ of dimension $3$.
%\end{proof}

%As a direct consequence of Corollary \ref{ro faithful on Z_p} we have the following theorem. 


%\begin{proof}
 %   
%\end{proof}


%\begin{proposition}\label{on order of LeftMul of LSS}
%Let $Q$ be a non-affine latin subdirectly irreducible quandle of size $16p$ as in Proposition \ref{SI 16p}. Then $|L_x|=3$ for every $x\in Q$.
%\end{proposition}

%\comment{Check!}

%\begin{proof}

%Let $Q=\mathcal{Q}(G,H,f)$ where $G=\dis(Q)$, $f=\widehat{L}_a$ and $H=Fix(f)$. According to Proposition \ref{Dis for SI} $G\cong \gamma_2(G)\rtimes K\cong\Z_p^3\rtimes K$ where $K$ is an extraspecial $2$-group. The block $[x]_{\gamma_Q}$ is isomorphic to $\aff(\gamma_2(G)/\gamma_2(G)_x,f|_{\gamma_2(G)})\cong \aff(\mathbb{Z}_p,t)$ for some $t\neq 1$. 


%Let $x,y\in Q$ such that $[x]_{\gamma_Q}\neq [y]_{\gamma_Q}$ and $S=Sg(x,y)$. The quandle $Q/\gamma_Q$ is not $2$-generated (see \cite[Lemma 4.9]{Principal}), so $|S|\in \{4,4p\}$. If $|S|=4p$ then $S\in \LSS(p,4,\sigma_Q)$ and it is not latin. So $|S|=4$ and so $|y^{L_x}|=3$. If $S=\mathcal{Q}_4\times \aff(\Z_p,t)$ then $|y^{L_x}|=\lcm\{3,|t|\}$.  \comment{Can $S$ be abelian of size $4p$??} If $|S|=4$ then $|y^{L_x}|=3.$

%The restriction of $f$ to $\gamma_2(G)$ is diagonalizable with eigenvalues $1$ and $t\neq 1$. We denote by $F$ the matrix with respect to a basis of eigenvectors. %Let $x\in K=G/\gamma_2(G)$ be a non central element of $K$. 
%since $x$ does not normalize $\left(\dis_\alpha\right)_a$ (see \ref{faithful action and sigma =0}). 
%The center of $K$ is fixed by $f_{\gamma_2(G)}$ and its induced mapping over $K/Z(K)\cong G/\gamma_1(G)$ is $f_{\gamma_1(G)}$. The order of $f_{\gamma_2(G)}$ and the order of $f_{\gamma_1(G)}$ coincide (since $Q_K$ is a covering of $Q/\alpha$, see \cite{covering_paper}) and so $f_{\gamma_2(G)}^3(x)=x $ for every $x\in K$ and we have %where $d\in\gamma_2(G)$. Hence
% \begin{equation}\label{commutator condition}
% \rho_{f_{\gamma_2(G)}^3(x)}=F^3\rho_{x}F^{-3}= \rho_{x}.
% \begin{bmatrix} 
%1 & 0\\
%0 & t^k
%\end{bmatrix} 
%\begin{bmatrix}
%a & c\\
%b & d
%\end{bmatrix}
%\begin{bmatrix} 
%1 & 0\\
%0 & t^{-k}\end{bmatrix}=\begin{bmatrix}
%a	& ct^{-k}\\
%bt^k	 &  d
%\end{bmatrix}=
% \rho_{x}=\begin{bmatrix} 
%a & c\\
%b & d\end{bmatrix}\, ,
% \end{equation}
% The coefficient $b$ is not $0$ since $N(\dis(Q)_a)=\dis(Q)_a$ for every $a\in Q$, and $x\notin \dis(Q)_a\cap K=Z(K)$. 
%Thus, $\rho(K)$ centralizes the diagonal matrix $F^k$ with entries $1$ and $t^3$. The eigenspaces of $F^3$ are $K$-invariant subspaces. Since $\rho$ is irreducible then $t^3=1$ and accordingly, $|z^{L_x}|=3$ whenever $z\in [x]_{\gamma_Q}$. 
%
%Equation \eqref{commutator condition} is equivalent to $c(t^{-k}-1)=b(t^k-1)=0$. If $b=c=0$ for every $x\in K$, then the image of $\rho$ is abelian and so it is $K$, contradiction. Thus we have $t^k=1$. %If $b=c=0$, the matrix $\rho_x$ is diagonal with respect to this basis, then $F\rho_x F^{-1}=\rho_{f(x)}=\rho_x$. The action $\rho$ is faithful and so $f(x)=x$ and so $x\in K\cap Fix(f)=Z(K)$, contradiction. Hence $t^k=1$.
%\begin{displaymath}
%F^k\rho_x-\rho_x F^k=\begin{bmatrix}
%1 & 0\\
%0 & t^k
%\end{bmatrix} \begin{bmatrix}
%a & c\\
%b & d
%\end{bmatrix}- \begin{bmatrix}
%a & c\\
%b & d
%\end{bmatrix}\begin{bmatrix}
%1 & 0\\
%0 & l^k
%\end{bmatrix}=\begin{bmatrix}
%0 & c(t^k-1)\\
%-b(t^k-1) & 0
%\end{bmatrix}.
%\end{displaymath}
%Since $x$ does not normalize $\left(\dis_\alpha\right)_a$, then $b\neq 0$, and so $t^k=1$. 

%Therefore $|L_x|=\lcm\setof{|y^{L_x}|}{y\in Q}=3$ for every $x\in Q$.
%\end{proof}
In the following proof, we use the fact that if $Q=\aff(A,f)$ is a finite connected quandle, then $\aut{Q}\cong A\rtimes C_{\aut{A}}(f)$, so in particular if $A=\Z_2^2\rtimes \Z_p$ for an odd prime $p$, then $\aut{A}=\aut{\Z_2^2}\times \aut{\Z_p}$, consequently $C_{\aut{A}}(f_1,f_2)=C_{\aut{\Z_2^2}}(f_1)\times C_{\aut{\Z_p}}(f_2)$.




\begin{lemma}\label{lemma 4 chain}
Let $Q$ be a latin quandle of size $16p$ such that $\con(Q)$ is a chain of length $4$. Then $\ma_Q=\{ \gamma_Q\}$ and $|Q/\gamma_Q|=16$.
%\begin{itemize}
    %\item[(i)] 
   % $\gamma_Q$ is the unique minimal congruence of $Q$ and $|Q/\gamma_Q|=16$.
    %$\con(Q)$ is: 

%\begin{center}
%		$\xymatrixrowsep{0.15in}
%		\xymatrixcolsep{0.15in}
%		\xymatrix{ 
%			 1_Q  \ar@{-}[d]\\
%			\alpha\ar@{-}[d] \\
%			 \gamma_Q\ar@{-}[d] \\
%			0_Q 
%		}$
%		\end{center}

%    \item[(ii)] $|Q/\gamma_Q|=16$.
    
    
%   \item[(ii)] The blocks of $\alpha$ are non-abelian quandles of size $4p$ and $|L_x|_{\gamma_Q}|=3$.  

%\item[(ii)] $Z(\dis(Q))=1$. 
%\end{itemize}
\end{lemma}

%    \comment{Prove or justify that $\aut{\mathcal{Q}_4\times \aff(\Z_p,f)}\cong \Z_2^2\times \Z_p \rtimes (\Z_3\times \Z_{p-1})$??}


\begin{proof}

%\begin{center}
%		$\xymatrixrowsep{0.15in}
%		\xymatrixcolsep{0.15in}
%		\xymatrix{ 
%			 1_Q  \ar@{-}[d]\\
%			\alpha\ar@{-}[d] \\
%			 \beta\ar@{-}[d] \\
%			0_Q 
%		}$
%		\end{center}

    %(i) 
    
    Let $\Ma_Q=\{\alpha\}$ and $\ma_Q=\{\beta\}$. Note that $|Q/\beta|\in \{16,4p\}$ and $\con(Q/\beta)$ is a chain of length $3$. Assume that $|Q/\beta|=4p$. In this case, we can also assume that $Q/\beta$ is not abelian (otherwise $Q/\beta$ would be directly decomposable), that is, $\gamma_Q=\alpha$. Therefore $p=1\pmod{3}$, since non-abelian latin quandles of size $16p$ exist only for such primes. We have to distinguish between three cases.
    \begin{itemize}

    %\item Assume that $|Q/\alpha|=16$. Then $\alpha\in \ma_Q$, contradiction.
    
    \item Assume that $|Q/\alpha|=p$. Then $Q/\beta\in \LSS(4,p,\sigma_Q)$ and so $p=7$ according to \cite[Proposition 5.3]{LSS}. Such quandles are not latin. 
    
        \item Assume that $|Q/\alpha|=4$ and $[x]_\alpha$ is in $\LSS(4,p,\sigma_Q)$. Then, $p=7$ again, and the blocks of $\alpha$ are not latin.  
                
        
        
%        \item Assume that $|Q/\alpha|=4$ and $[x]_\alpha$ is in $\LSS(p,4,\sigma_Q)$. Then $Q/\beta$ has size $16$ and so it is abelian. Therefore $\gamma_Q\leq \beta$ and since $Q$ is not abelian the equality holds.


                \item Assume that $|Q/\alpha|=4$ and $[x]_\alpha\cong \mathcal{Q}_4\times \aff(\Z_p,k)$. %Then $|Q/\beta|=4p$ and it is not abelian (otherwise $Q/\beta$ would be subdirectly reducible and so $|\Ma_Q|\geq 2$). Hence, $p=1\pmod{3}$ according to the classification of connected quandles of size $4p$ obtained in \cite{LSS}. 
                We can assume that $p\neq 3$, since $p=1\pmod{3}$. Let $N$ be a $p$-Sylow subgroup of $\dis(Q)$. Since $|\dis(Q/\alpha)|=4$, hence $N$ is contained in $\dis^\alpha$ and since $|\dis(Q/\beta)|=8p^2$ we see that $p^2$ divides $|N|$. On the other hand, $\dis^\alpha$ embeds into $\aut{[x]_\alpha}^2\cong [(\Z_2^2\times \Z_p)\rtimes (\Z_3\times \Z_{p-1})]^2$ since $Q/\alpha$ is $2$-generated (again by using Corollary \ref{embedding as quandle}). The group $\aut{[x]_\alpha}^2$ has a unique normal $p$-Sylow subgroup of size $p^2$ and therefore $|N|=p^2$ and $N$ is normal in $\dis^\alpha$. Then $N$ is characteristic in $\dis^\alpha$ and is therefore normal in $\lmlt(Q)$ accordingly. Hence $\delta=\mathcal{O}_N$ is a congruence of $Q$ with blocks a power of $p$. Hence $\delta=\beta$ and so $|Q/\beta|=16$, contradiction.  
        \end{itemize}
    Therefore, $|Q/\beta|=16$ and so $Q/\beta$ is abelian (as latin quandles of size $16$ are affine). Then $\gamma_Q\leq \beta$ and since $Q$ is not abelian, then $\beta=\gamma_Q$. 
    %
%    (ii) The blocks of $\alpha$ have size $4p$ and they are not abelian according to Lemma \ref{Sylow}. So $[x]_\alpha\in \LSS(p,4,\sigma_Q)$ and $L_x|_{[x]_\alpha}$ is $3$. So $L_x|_{[x]_{\gamma_Q}}$ is also $3$.% and in particular $p=1\pmod{3}$.
%
%
%
%(ii) The congruence $\gamma_Q$ is the unique minimal congruence of $Q$ that is not nilpotent. Then according to Lemma \ref{dis_gamma} we have $Z(\dis(Q))=1$.
%
%    (iii) If $Z(\dis(Q))\neq 1$ then $\zeta_Q\neq 0_Q$ and so $\gamma_Q\leq \zeta_Q$ since $\gamma_Q$ is minimal. So $Q$ is nilpotent and directly reducible, contradiciton.
\end{proof}

%We can characterize the displacement group of latin quandles of size $16p$ with congruence lattice given by a chain of length $4$.


We can use the coset representation of connected quandles in order to conclude that the congruence lattice of latin quandles of order $16p$ is a chain of length $3$ but in a few cases for small primes. 



\begin{proposition}\label{chain}
    Let $Q$ be a subdirectly irreducible latin quandle of size $16p$. 
    \begin{itemize}
        \item[(i)] If $\con(Q)$ is a chain then it has length $3$.
        \item[(ii)] If $\con(Q)$ is not a chain then $p\leq 7$.
    \end{itemize}
\end{proposition}

%\begin{proposition}\label{Dis for SI}
  %  Let $Q$ be a subdirectly irreducible latin quandle of size $16p$ for $p>7$ such that $\con(Q)$ is not a chain. Then
   % $$\dis(Q)=\begin{cases}
    %    \Z_p^3\rtimes_{\rho} \Z_2^4\\
     %   \Z_p^3\rtimes_\rho K
    %\end{cases}$$
    %where $\rho$ is an irreducible faithful action and $K$=SmallGroup(32,49). In the second case $Z(\gamma_1(\dis(Q)))=1$.
%$$\dis(Q)=\Z_p^2\rtimes_{\rho} K$$
%where $\rho$ is a faithful action, $\gamma_1(G)=\Z_p^2\rtimes_\rho Z(K)$, $K=SmallGroup(32,49)$ and $Z(\dis(Q))=Z(\gamma_1(\dis(Q)))=1$.
%\end{proposition}

%\comment{si può fare forse a meno di invocare la classificazione dei quandles of size $32$, se si riesce a dire che $K$ è extraspeciale in un altro modo}

%\comment{Usare notazione $K_{i}?$}

%\comment{the center of SmallGroup(32,47) is $\Z_2^3$. Hence the action on $\Z_2^2$ is not faithful. Then $t=3$. Usando \ref{center = -1} $t=2$, contradiction!}

\begin{proof}
Let $G=\dis(Q)$ and $\alpha=\gamma_Q$. Assume that $\con(Q)$ is not a chain and $p>7$ or that $\con(Q)$ is a chain of length $4$. According to Lemma \ref{SI 16p case 4p} and Lemma \ref{lemma 4 chain}, $\alpha$ is the unique minimal congruence of $Q$, the factor $Q/\alpha$ is $\mathcal{Q}_4^2$ and then $G/\gamma_1(G)\cong \Z_2^4$. 


Moreover, in both cases, $Q$ is not nilpotent, so $\gamma_2(G)=\dis_\alpha$ and $Z(G)=1$, according to Lemma \ref{dis_gamma}. The blocks of $\alpha$ have size $p$, the quandle $Q/\alpha$ is generated by $3$ elements, so the group $\dis_\alpha$ embeds into $\Z_p^3$ according to Corollary \ref{embedding as quandle}(ii). Moreover $|\dis_\alpha|\neq p$ since $\alpha$ is not central (see Lemma \ref{central 2}).


Let $Q=\mathcal{Q}(G,G_x,f)$ where $f=\widehat{L_x}$ for some $x\in Q$. According to \cite[Lemma 1.10]{Involutory} the principal quandle $Q'=\mathcal{Q}(G/\gamma_2(G),f_{\gamma_2(G)})$ is a connected cover of $Q/\gamma_Q$, $K=G/\gamma_2(G)\cong \dis(Q')$ and it has size a power of $2$. In particular the group $K=G/\gamma_2(G)$ is a $2$-step nilpotent group $2$-group and so $G\cong \gamma_2(G)\rtimes_\rho K$. Moreover $K/\gamma_1(K)\cong G/\gamma_1(G)$ is an elementary $2$-group and so by Lemma \ref{derived si elemetary} $\gamma_1(K)$ is an elementary abelian $2$-group. Moreover $\gamma_1(G)=\gamma_2(G)\rtimes_\rho \gamma_1(K)$. The $2$-Sylow of $\dis^{\gamma_Q}=\gamma_1(G)$ is cyclic according to Lemma \ref{N cyclic}, so the $2$-Sylow of $\gamma_1(G)$ is isomorphic to $\Z_2^s$ for $s=0,1$. 

We have to distinguish between two cases.

%The congruence $\alpha$ is the unique minimal congruence of $Q$ that is not nilpotent. Then according to Lemma \ref{dis_gamma} $Z(G)=1$. %The action $\rho$ is faithful since $Z(G)=1$ (see Lemma \ref{no center}).

\begin{itemize}
    \item Assume that $s=0$. Then $\gamma_1(G)=\gamma_2(G)$ and so $G=\Z_p^t\rtimes_\rho \Z_2^4$ for $t\leq 3$. By Lemma \ref{ro faithful on Z_p}, $\rho$ is not faithful since $t\leq 3\leq 4$, leading to a contradiction. 
 

\item Assume that $s=1$. The quandle $Q'$ is a connected quandle of size $32$ with factor $Q'/\gamma_{Q'}$ isomorphic $\mathcal{Q}_4^2$, having $7$ congruences. Then, looking at Table \ref{32}, we have that $Q'$ is {\tt SmallQuandle(32,j)} for $j=2,3$ and accordingly $K=G/\gamma_2(G)=\dis(Q')$ is {\tt SmallGroup(32,i)} for $i=47,49$. So $G=\Z_p^t\rtimes_\rho K$ for $t=2,3$. Both groups have a subgroup isomorphic to $\Z_2^3$, so we have $t=3$ by Corollary \ref{ro faithful on Z_p}. According to Corollary \ref{center of dis^alpha} then $Z=Z(\gamma_1(G))=Z(\dis^{\alpha})=1$ since $\dis^\alpha$ is not a group of size a power of a prime. On the other Corollary \ref{center = -1} implies that $t$ is even, a contradiction. \qedhere%we have $Z=Z(\gamma_1(G))\neq 1$. Since $\beta=\mathcal{O}_Z\leq \alpha$, then we have $\beta=\alpha$ and so $\dis_\alpha\leq Z$. Consequently, $\gamma_1(G)$ is abelian and has size $2p$. By Lemma \ref{Sylow} the Sylow subgroups of $\gamma_1(G)$ provide trivially intersecting congruences, leading to a contradiction.  \qedhere
%
%So the elements of $\gamma_1(K)$ commute with the elements of $K$ and the elements of $\dis_\alpha$. Hence $Z(K)\leq Z(G)=1$, contradiction. \qedhere
\end{itemize}

%
%
%According to Proposition \ref{no reps} there is no such groups for $t=2$. Thus, $t=3$. 
%
%Finally, note that $\rho$ is irreducible according to Lemma \ref{irreducible2} and Lemma \ref{irreducible}. According to \ref{no irred Z_2^n} there is no irreducible representation of $\Z_2^4$ of dimension $3$.
\end{proof}

%\begin{proposition}
%Let $Q$ be a subdirectly irreducible latin quandle of order $16p$ for $p>7$. Then $\con(Q)$ is a chain of length $3$.    
%\end{proposition}


%\begin{proposition}\label{group 4-chain}
 %   Let $Q$ be a latin quandle of size $16p$ such that $\con(Q)$ is a chain of length $4$. Then $$\dis(Q)\cong \Z_p^t\rtimes_\rho \Z_4^2$$ where $\rho$ is a faithful action, $t\in \{2,3\}$ and $Z(\dis(Q))=1$.
%\end{proposition}
%\comment{According to \ref{case q=2 then n even}, if $t=3$ then $e(p,2)=1$, i.e. $p=1\pmod{4}$.}

%\comment{forse dato che $Z(G)=1$ viene che $t$ è multiplo di $e$ che può essere $1$ o $2$. Se $t=3$ quindi $e=1$, cioè $p=1\pmod{4}$? Se $t=2$ dal fatto che $\rho_x^2$ è diagonale e deve commutare con matrici diagonali non scalari allora viene che $\rho_x$ è diagonalizzabile}




%\item[(2)] \comment{CHECK: Sia $G=\dis(Q)$. Se $\gamma_1(G)\neq \gamma_2(G)$ allora viene che $\gamma_1(G)/\gamma_2(G)$ is cyclic of order $2$. Moreover, if $K=G/\gamma_2(G)$ we have $K$ is $2$-step nilpotent and $[K,K]$ is cyclic of order $2$ and so also $G/\gamma_1(G)=K/\gamma_1(K)$ is $2$-elementary abelian. Moreover $K$ has $2$ non-trivial normal subgroups. The quandle $Q_{\gamma_2(G)}$ is connected of size $32$ and it has $2$ non-trivial congruences and so necessarily $Q_{\gamma_Q(G)}$ is SmallQuandle(32,6) and $Q/\gamma_Q$ is SmallQuandle(16,4). Moreover $K=$SmallGroup(32,49).}

%\item[(3)] \comment{What if $\gamma_1(G)=\gamma_2(G)$? In this case $\gamma_1(G)$ is an elementary $p$-abelian group and $G=\Z_p^k\rtimes_\rho G/\gamma_1(G)$. What can I say about $G/\gamma_1(G)$?} 



%\comment{si può fare forse a meno di invocare la classificazione dei quandles of size $32$, se si riesce a dire che $K$ è extraspeciale in un altro modo}


 %   \comment{I latini di ordine $16$ sono affini sopra $\Z_2^4$! Il caso $\dis(Q/\gamma_Q)=\Z_4^2$ non può avvenire!} 

%\begin{proof}
%We can assume that $\con(Q)$ is a chain according to Proposition \ref{chain}. Assume that the length of $\con(Q)$ is $4$ and let $\Ma_Q=\{\alpha\}$, $\ma_Q=\{\beta\}$ and $G=\dis(Q)$. According to Lemma \ref{lemma 4 chain} we have that $|Q/\beta|=16$ and $\beta=\gamma_Q$ is the unique minimal congruence. The quandle $Q$ is not nilpotent, so we have $\dis_{\beta}=\gamma_2(G)$ and $Z(G)=1$ according to Lemma \ref{dis_gamma}. Latin quandles of size $16$ are affine over $\Z_2^4$, so $G/\gamma_1(G)\cong \Z_2^4$. The quandle $Q/\beta$ has a congruence whose factor and blocks are of size $4$, so they are both generated by $2$ elements. According to Lemma 6.1 of \cite{GB} then $Q/\beta$ is generated by three elements. Thus, $\dis_\beta$ embeds into $\Z_p^3$ according to Corollary \ref{embedding as quandle}, and it is not cyclic since $\beta$ is not central (see Lemma \ref{central 2}).

%The congruence $\gamma_Q$ is the unique minimal congruence of $Q$ that is not nilpotent. Then according to Lemma \ref{dis_gamma} we have $Z(G)=1$.

%Following the same argument as in Proposition \ref{chain},
%we can define the principal quandle $Q'$ as in Theorem \ref{chain} that is a covering of $Q/\beta$. According to \cite[Proposition 3.10]{Involutory} $|Q'|=|Q/\beta||\gamma_1(G)/\gamma_2(G)|$ is a power of $2$ and the $2$-Sylow subgroup of $\gamma_1(G)$ is cyclic (see Lemma \ref{N cyclic}). Hence, the group $K=G/\gamma_2(G)$ is a $2$-step nilpotent $2$-group of size $2^{4+s}$ where $s=0,1$ and accordingly $G=\gamma_2(G)\rtimes_\rho K$. Following the same argument as in Proposition \ref{chain} in both cases we get a contradiction.
%
%
%The map $G\longrightarrow K/\gamma_1(K)$ is a group homomorphism onto an abelian group, so $\gamma_1(G)\leq \gamma_2(G)\rtimes \gamma_1(K)$. The other inclusion is clear, so $\gamma_1(G)=\gamma_2(G)\rtimes_\rho \gamma_1(K)$. Hence, $K$ is a $2$-group, it is $2$-step nilpotent. %and $K/\gamma_1(K)$ is elementary $2$-abelian. According to Lemma \ref{derived si elemetary} $\gamma_1(K)$ is elementary $2$-abelian.
%
%
%\begin{itemize}
 %   \item Assume that $s=0$. Then $\gamma_2(G)=\gamma_1(G)$ and $G=\Z_p^t\rtimes_{\rho} \Z_2^4$ where $t\in \{2,3\}$, $\rho$ is a faithful action (since $Z(G)=1$ we can apply Lemma \ref{no center}). %and $A$ is an abelian group of size $16$. %According to the classification of connected quandles of size $16$, then $A\in \{\Z_2^4, \Z_4^2\}$. 
%According to Lemma \ref{ro faithful on Z_p} then $\rho$ is not faithful. %Therefore $A=\Z_4^2$. 
%
 %
%
%\item Assume that $s=1$. The quandle $Q'/\gamma_{Q'}$ has $3$ congruences. Looking at Table \ref{32} we find that $K=G/\gamma_2(G)=\dis(Q_{\gamma_2(G)})$ is {\tt SmallGroup(32,i)} for $=47,49$, and consequently $G=\Z_p^t\rtimes_{\rho} K$ where $Z(G)=1$ and so $\rho$ is faithful according to Lemma \ref{no center}. Following the same argument as in Proposition \ref{chain} 
%Both groups have a subgroup isomorphic to $\Z_2^3$, so we have $t=3$ by Corollary \ref{ro faithful on Z_p}. Hence, by Corollary \ref{center = -1} we have $Z=Z(\gamma_1(G))=Z(\Z_p^t\rtimes_\rho \gamma_1(K))\neq 1$. Then $0_Q\neq \delta=\mathcal{O}_Z\leq \beta$. Thus $\delta=\beta$ by the minimality of $\beta$ and so $\dis_{\beta}\leq Z$. Consequently, $\gamma_1(G)$ is abelian and it has size $2p^t$. By Lemma \ref{Sylow}, the Sylow subgroups of $\gamma_1(G)$ provide trivially intersecting congruences, leading to a contradiction. \qedhere
%
%\end{itemize}
%
%\end{proof}


%\begin{theorem}
%There are no latin quandles of order $16p$ with congruence lattice given by a chain of height $4$.
%\end{theorem}

%\begin{proof}
 %   Let $G=\Z_p^t\rtimes_\rho \Z_4^2$ where $t=2,3$, $Z(G)=1$. According to Proposition \ref{group 4-chain}, latin quandles of order $16p$ with congruence lattice given by a chain of length $4$ are therefore given by $\mathcal{Q}(G,Fix(f),f)$, $[G:Fix(f)]=16p$ and $Fix(f_{\Z_p^t})=1$. We can assume that the action $\rho$ is diagonal, according to Proposition \ref{action is diagonal}, and so $G=G_{M,N}(t)$ as defined in Section \ref{Sec: G_{M,N}}. According to Lemma \ref{determinant} we can assume that $\det(N)=\det(M)=1$. Corollary \ref{t=2 forbidden} implies that $t=3$ and so we can also assume that $G\cong G_{A,B}(3)$ where $A,B$ are defined in \eqref{def A,B}. Therefore $|Fix(f)|=p^2$ and $Fix(f_{\Z_p^3})=1$. We can finally use Proposition \ref{chain 4 then p=3} in order to conclude that there are no such automorphisms and so there are no such quandles.
%\end{proof}

%\comment{Viene $p=1\pmod{12}$.}

%\comment{I generatori di $\Z_4^2$ vengono mappati in matrici con sulla diagonali elementi tali che $x^4=1$. Se $a$ è il generatore del gruppo di ordine $4$ in $\Z_p^*$ then $x\in \{1,a,-1,a^3\}$, i.e. $x\in \langle a\rangle$. Se $t=2$ vengono proprio tutte le matrici possibili.}
%
%\begin{itemize}
%\item[(1)] \comment{The blocks of $\gamma_Q$ are of the form $\aff(\Z_p,k)$ with $|k|=3$ perchè i blocchi di $\alpha$ sono latini di ordine $4p$ e non abeliani! \label{ecco!}}
%
%    \item[(2)] \comment{$Q/\gamma_Q$ is an S.I. latin quandle of size $16$. Then either it is SmallQuandle(16,4) or it is affine over $\Z_4^2$. Can the second option happen? }
%\item[(2)]  \comment{Let $G=\Z_p^t\rtimes_\rho \Z_4^2$. and $f\in \aut{F}$. Then $F=f|_{\Z_p^t}$ is a diagonalizable matrix with eigenvalues $1$ and $w$ and so $\det(F)=w$. Since $F\rho_x F^{-1}=\rho_{f_{\gamma_1(G)(x)}}$ and $f_{\gamma_1(G)}$ has no fixed point then $F$ does not centralize any $\rho_x$.}
%
%\item[(3)] \comment{Let $Q=\mathcal{Q}(G_{A,B}(t),Fix(f),f)$ be a latin quandle. Then $Fix(f)\leq \Z_p^t$.}
%\end{itemize}

Finally, we deal with quandles with congruence lattice given by a chain of length $3$ and we obtain further constrains of the values of $p$ and on the structure of the displacement group.
\begin{proposition}\label{LSS 16p}
Let $Q$ be a latin quandle of size $16p$ such that $\con(Q)$ is a chain of length $3$. Then: 
\begin{itemize}
%\item[(i)] $\con(Q)$ is
%\begin{center}
%		$\xymatrixrowsep{0.15in}
%		\xymatrixcolsep{0.15in}
%		\xymatrix{ 
%			 1_Q  \ar@{-}[d]\\
%			\gamma_Q\ar@{-}[d] \\
%			0_Q 
%		}$
%		\end{center}
\item[(i)]  $Z(\dis(Q))=1$.

\item[(ii)] $\dis(Q)=\Z_2^{4+k}\rtimes_\rho \Z_p$ for $k\leq 4$, and $\rho$ is a faithful action.



\item[(iii)] $|Q/\gamma_Q|=p\in \{3,5,7,17,31,127\}$.


%\item[(iii)] If $|Q/\gamma_Q|=16$ then $Q\in \LSS(p,16,\sigma_Q)$.

%\item[(iii)] If $|Q/\gamma_Q|=p$ then $p\in \{3,5,7,17,31,127\}$.     \comment{Can I have subquandles of size $4p$? }% and $[x]_\alpha\cong \mathcal{Q}_4\times \mathcal{Q}_4$.



%$Q$ is $2$-step solvable and $\zeta_Q=0_Q$. 

%\item[(iv)] $Q=Sg(x,y,z)$ for every triple $x,y,z$ such that $x\in [y]_{\gamma_Q}\neq [z]_{\gamma_Q}$ and $x\neq y$.
\end{itemize}


\end{proposition}


 
 




 


\begin{proof}
%
%(i) The quandle $Q$ is solvable but not abelian. Thus $0_Q\neq \gamma_Q\neq 1_Q$ and so $\gamma_Q$ is the unique non-trivial congruence of $Q$. The quandle $Q$ is not nilpotent, then according to Lemma \ref{dis_gamma} we have that $Z(\dis(Q))=1$.

Let $G=\dis(Q)$ and $\alpha=\gamma_Q$. Note that $Q$ is solvable, $\alpha\in \ma_Q\cap\Ma_Q$ and so it is abelian. Therefore, both $Q/\alpha$ and $[x]_\alpha$ have size a power of a prime.

(i) Clearly $\alpha$ is the only proper congruence, and so it is minimal. Therefore, $Z(G)=1$ according to Lemma \ref{dis_gamma}.

(ii) We separate the discussion into some cases.

\begin{itemize}
    \item Assume that $|Q/\alpha|$ has size $4$. Then, the blocks of $\alpha$ have size $4p$, which is a contradiction. % and $\alpha$ is abelian. Thus, $\dis_\alpha$ is abelian and $4p$ divides $|\dis_\alpha|$. According to Lemma \ref{Sylow}, $Q$ has two congruences below $\alpha$, contradiction.

%The congruence $\alpha$ is minimal, since it has blocks of size $p$. Then according to Lemma \ref{cong of directly irred} $\alpha$ is the unique congruence of $Q$. Thus, $Q\in \LSS(p,16,\sigma_Q)$. According to \cite[]{LSS} $\sigma_Q=0$. 

\item If $|Q/\alpha|=16$, then the blocks of $\alpha$ are simple, so $Q\in \LSS(p,16,\sigma_Q)$. But according to Proposition \ref{no LSS}, the class $\LSS(p,16,\sigma_Q)$ is empty.


\item Assume that $|Q/\alpha|$ has size $p$. The blocks of $\alpha$ are affine over $\Z_2^4$. 
 %Let $G=\dis(Q)$. According to \cite{GB} and \cite{CP} we have that
%$$[\dis^{\gamma_Q},\dis(Q)]=\gamma_2(G)\leq \dis_{\gamma_Q}\leq \dis^{\gamma_Q}= \gamma_1(G).$$
%The group $G$ is not nilpotent and $\dis_{\gamma_Q}$ is a minimal subgroup of $\N(Q)$, hence $\dis_{\gamma_Q}=\gamma_2(G)$. 
The group $G/\gamma_1(G)$ is cyclic and $\gamma_1(G)/\gamma_2(G)$ is generated by the elements $\setof{[x,y]\gamma_2(G)}{x,y\in X}$, where $X$ generates $G/\gamma_1(G)$ (see \cite[Proposition 9.2.5]{Sims}). We conclude that $\gamma_2(G)=\gamma_1(G)$. Since $\dis_{\alpha}$ embeds into $\Z_2^{4+k}$ for $k\leq 4$ (see Corollary \ref{embedding as quandle}), according to Lemma \ref{dis for 3-chains}, we have that $G\cong \Z_2^{4+k}\rtimes_\rho \Z_p$ and $\rho$ is faithful .% (see Lemma \ref{no center}).
\end{itemize}


(iii) The prime $p$ divides $|GL_2(4+k)|$ for $0\leq k\leq 4$, therefore $p\in \{3,5,7,17,31,127\}$. %
%
%The congruence $\gamma_Q$ is minimal and it is abelian. Hence, $Q$ is $2$-step solvable. If $\zeta_Q\neq 0_Q$ then $\gamma_Q\leq \zeta_Q$ and so $Q$ is nilpotent and then abelian, contradiction.
%
%(iv) Let $x,y,z\in Q$ such that $x\in [y]_{\gamma_Q}\neq [z]_{\gamma_Q}$ and $S=Sg(x,y,z)$. Then $S$ maps onto $Q/\gamma_Q$, but $S$ is not isomorphic to $Q/\gamma_Q$. 
%Then $|S|\in \{4p,16p\}$. In the first case then $S=Sg(x,y,z)\in \LSS(4,p,\sigma_Q)$ and accordingly $p=7$ and $S$ is not latin. Thus, $S=Q$. \comment{Can $S$ be affine of size $4p$?}
\end{proof}


\subsection{Computational Results}\label{computational}

To classify all subdirectly irreducible latin quandles of size $16p$, we can proceed computationally. Latin quandles are, in particular, connected. Consequently, they admit a minimal coset representation over their displacement group.
    By Lemma~\ref{SI 16p case 4p} and Lemma~\ref{LSS 16p}, such quandles are coset quandles over only finitely many groups. For these reasons, the classification problem reduces to a finite computation. Moreover, once we identify the suitable groups, we can easily obtain isomorphism classes using Theorem \ref{iso theo}.
    
In this section, we describe the algorithm we performed in GAP in order to deal with the cases described above.

    Let $Q$ be a latin quandle of size $16p$. Then $Q$ has a minimal coset representation as $\mathcal{Q}(G,Fix(f),f)$ where $G\cong \dis(Q)$.



\begin{enumerate}
    \item 
    If the congruence lattice of $Q$ not a chain, then, by Lemma~\ref{SI 16p case 4p}, $Q$ is a coset quandle over a centerless group whose order divides $2^6 p$ for $p \neq 3$, or $2^6 \cdot 9$ when $p = 3$ (note that $16p$ must divide $|\dis(Q)|$ in each case since $|\dis(Q)|=|Q||\dis(Q)_x|$ for every $x\in Q$). %By direct computation, we enumerate all the paris $(G,f)$ where $G$ groups of admissible order as above, and discard those lacking suitable automorphisms, namely, automorphisms $f$ such that $\mathcal{Q}(G_{p,k},\fix(f),f)$ is a latin quandle of size $16p$.
    If the congruence lattice of $Q$ is a chain, then by Lemma~\ref{LSS 16p}, $Q$ is a coset quandles over a group of the form $G_{p,k} = \mathbb{Z}_2^{4+k} \rtimes_\rho \mathbb{Z}_p$, where $0 \leq k \leq 4$ and $p \in \{3, 5, 7, 17, 31, 127\}$, with $Z(G_{p,k}) = 1$. This leads to a similar investigation of potential groups of the correct order, satisfying the appropriate structural conditions.


    \item The second step is to construct all the coset quandles of the form $\mathcal{Q}(G, \fix(f), f)$ over the groups  identified above with the right properties.
    
    \item We proceed by filtering by isomorphism classes using Theorem \ref{iso theo}. 
    \end{enumerate}
    
   The algorithm computes that there are only two such quandles, namely $Q_p=\mathcal{Q}(G_p,Fix(f_p),f_p)$ for $p=3,5$ where the presentations of the groups used for the coset representation are
   \begin{align*}
       G_3 &= \texttt{SmallGroup(48,50)} = \langle a, b, c \ \colon\ 
        a^{3} ,\,
        c^{2} ,\,
        b^{2} ,\,
        (c b)^{2} ,\,
        (b a)^{3} ,\,
        (c a)^{3} ,\,
        (a^{-1} c a b)^{2} ,\,
        (c a^{-1} b a)^{2} 
        \rangle \cong \Z_2^4\rtimes \Z_3,\\ G_5 &= \texttt{SmallGroup(80,49)}=\langle a, b \ \colon \ 
        a^{5} ,\,
        b^{2} ,\,
        (b a^{-1} b a)^{2} ,\, 
        (b a^{-1})^{5} ,\, 
        (b a^{-2} b a^{2})^{2} 
        \rangle\cong \Z_2^4\rtimes \Z_5,
        \end{align*}
    and the automorphisms are defined as
\[
    \begin{aligned}
    f_3 = \begin{cases}
    a \mapsto a^2ba^{-1}ba \\
    b \mapsto ba^{-1}ca \\
    c \mapsto bc
    \end{cases}
    \end{aligned}
    \hspace{2cm}
    \begin{aligned}
    f_5 = \begin{cases}
    a \mapsto a^{2} (a^{2} b)^{2} a^{-3} b a \\
    b \mapsto b a^{2} b a^{-2}
    \end{cases}
    \end{aligned}
    \]  
   Both of these quandles are subdirectly irreducible, and their congruence lattice is a 3-element chain. 
    



\section{Appendix}  \label{sec:appendix}



%\comment{not used}
%\begin{lemma}
%Let $K$ be a group and $\rho:K\longrightarrow GL_n(p)$ and $G=\Z_p^n\rtimes_\rho K$. If $\rho$ has a $1$-dimensional invariant subspace then $Z(\Z_p^n\rtimes_\rho [K,K])\neq 1$.
%\end{lemma}

%\begin{proof}
 %   If $V$ has dimension $1$ and it is $\rho$-invariant, then we have a morphism of group $\rho|_V:K\longrightarrow \aut{\Z_p}$ and $\aut{\Z_p}$ is abelian. Thus, $[K,K]\leq \ker{\rho|_V}$. In particular $V\leq Z(\Z_p^n\rtimes_\rho [K,K])$.
%\end{proof}





%\subsection{Faithful representation of SmallGroup(32,50)}

%The group SmallGroup(32,50) is an extraspecial $2$-group of size $32$ given by the central product of the quaternion group $Q_8$ and the dihedral group $D_8$. In this section we are denoting such group by $G$. The group $G$ contains a copy of the groups $D_8$ and $Q_8$ for which we are going to use the following presentations:
%\begin{align*}
 %   D_8&=\langle a,x\,|\, a^4=x^2=axax=1\rangle,\\
  %      Q_8&=\langle b,y\, |\, b^4=b^2y^2=yby^{-1}b=1\rangle.
%\end{align*}

%A presentation of $G$ is the following:
%\begin{align*}
%    G=\langle a,b,x,y\,|\,a^4=x^2=axax=b^4=b^2y^2=yby^{-1}b=1\rangle.
%\end{align*}
%\comment{Is the presentation correct? According to GAP it is NOT!!} Moreover we have that $Z(G)=[G,G]$ has size $2$ and $u^2\in Z(G)$ for every $u\in G$.



%We are interested in studying faithful representation of $G$ of dimension $2$ over finite fields of prime order. In particular we focus on faithful representation defining a semidirect product $\Z_p^2\rtimes_\rho G$ such that the subgroup $\Z_p^2\rtimes_\rho Z(G)$ has trivial center.

%\begin{lemma}
 %   Let $\rho:G\longrightarrow GL_2(p)$ be a faithful representation and $H=\Z_p^2\rtimes_\rho G$. If $H'=\Z_p^2\rtimes_\rho Z(G)$ has trivial center then $\rho(z)=-Id_2$ for $1\neq z\in Z(G)$.
%\end{lemma}

%\begin{proof}
 %   The group $G$ has a center of size $2$. Then $\rho(z)$ is diagonalizable since it has order $2$ (see Lemma 7.1 \cite{LSS}). The eigenvalues of $\rho(z)$ belong to $\{1,-1\}$. If there exists $v\in \Z_p^2$ such that $\rho_z(v)=v$ then $v\in Z(H')=1$. Thus, $\rho_z=-Id_2$.
%\end{proof}


%\begin{lemma}
 %   Let $\rho:G\longrightarrow GL_2(p)$ be a faithful representation such that $\rho_z=-Id_2$ for $1\neq z\in Z(D_8)$. Then, up to isomorphims, we have that
  %  \begin{align}\label{rho D_8}
   %     \rho_x=\begin{bmatrix}
    %        1 &0 \\
     %       0 & -1
      %  \end{bmatrix},\quad \rho_a=\begin{bmatrix}
       %     0 & -1 \\
        %    1 & 0
       % \end{bmatrix}.
   % \end{align}
%\end{lemma}

%\begin{proof}
 %   The order of $\rho_x$ is $2$ and so $\rho_x$ is diagonalizable (see Lemma 7.1 \cite{LSS}). Then we can assume that $$\rho_x=\begin{bmatrix}
 %       1 & 0\\
  %      0 &-1
   % \end{bmatrix}.$$
%Using that $\rho_a^2=\rho_z=-Id_2$ and $\rho_x \rho_a \rho_x\rho_a=1$ we have that
%\begin{align*}
 %   \rho_a=\begin{bmatrix}
  %      0 & -y\\
   %     y^{-1} & 0
   % \end{bmatrix}.
   % \end{align*}
%Then we can change the basis as follows:
%\begin{align*}
%   \begin{bmatrix}
 %       1 & 0\\
  %      0 & y
   % \end{bmatrix} \begin{bmatrix}
    %    0 & -y\\
     %   y^{-1} & 0
  %  \end{bmatrix}\begin{bmatrix}
   %     1 & 0\\
    %    0 & y
   % \end{bmatrix}^{-1}=\begin{bmatrix}
   %     0 & -1\\
    %    1 & 0
   % \end{bmatrix}.
    %    \end{align*}
%Thus, the unique representation (up to isomorphisms) with the desired property, is the representation defined in the statement. 
%\end{proof}
 



%\begin{lemma}
 %   Let $\rho:G\longrightarrow GL_2(p)$ be a faithful representation such that $\rho_z=-Id_2$ for $1\neq z\in Z(G)$ and \eqref{rho D_8} holds. Up to isomorphism, one of the following holds:

%\begin{itemize}
 %   \item[(i)] $\rho_y=\beta(s)$, $\rho_b=\alpha_\pm (t)$.
 %   \item[(ii)] $\rho_y=\alpha_+(s)$, $\rho_b=\alpha_- (t)$.
  %  \item[(iii)] $\rho_y=\alpha_-(s)$, $\rho_b=\beta(t)$.
   %     \item[(iv)] $\rho_y=\alpha_-(s)$, $\rho_b=\alpha_+(t)$.
%\end{itemize}
%Moreover, $p=1\pmod{(4)}$.
%\end{lemma}

%\begin{proof}
 %   Let $\rho_b$ be 
  % $$ \rho_b=\begin{bmatrix}
   %    s & u \\
    %   t & v
  % \end{bmatrix}.$$
 %  Since $\rho_b^2=\rho_z$ and $\rho_b\notin Z(G)$ then $s=-v$ and $s^2+tu=-1$.
   
  % Moreover, $[\rho_b,\rho_x],[\rho_b,\rho_a]\in \langle \rho_z\rangle$ then $u^2=t^2$, i.e. $u=\pm t$ and $ts=0$. Then $\rho_b\in \{\alpha_\pm (s),\beta(t),\, s,t\in \Z_p\}$. The same argument shows that $\rho_y\in  \{\alpha_\pm (s),\beta(t),\, s,t\in \Z_p\}$. The only choices such that the relation $\rho_y \rho_b \rho_y^{-1}\rho_b=1$ holds are those listed in (i)-(iv).  
%\end{proof}


\subsection{Groups theoretical results}

In this subsection we collect some results we are going to use to describe the displacement groups of latin quandles of size $16p$. In particular, we are going to prove some results on actions of $2$-groups on elementary abelian $p$-groups.

%\begin{lemma}
 %   Let $K$ be a nilpotent group of nilpotency length $2$. Then $exp([K,K])$ divides the exponent of $K/[K,K]$.
%\end{lemma}
%\comment{not used}

\begin{lemma}\label{derived si elemetary}
    Let $K$ be a $p$-group with nilpotency class $2$. If $K/\gamma_1(K)$ is elementary $p$-abelian group then $\gamma_1(K)$ is elementary $p$-abelian group.
\end{lemma}

\begin{proof}
    Let $x\in G$. The map $[x,-]:y\longrightarrow [x,y]$ is a group homomorphism with values in $\gamma_1(K)$. Moreover, the map factor through the center of $K$. Thus, we have a group homomorphism from $K/Z(K)$ to $\gamma_1(K)$. Since $\gamma_1(K)\leq Z(K)$ then we also have a group homomorphism from $K/\gamma_1(K)$ onto $K/Z(K)$. Hence $K/Z(K)$ is an elementary $p$-abelian group. Therefore, every element of the form $[x,y]$ has order $p$ and so $\gamma_1(K)$ is an elementary $p$-abelian group.
\end{proof}


\begin{lemma}\label{no center}
    Let $H$ be an abelian group, $K$ be a nilpotent group, and $G=H\rtimes_{\rho} K$. If $Z(G)=1$ then $\ker{\rho}=1$.
\end{lemma}

\begin{proof}
%If $K$ is abelian then $\ker{\rho}\leq Z(G)=1$.
%\end{proof}


%\begin{lemma}\label{no center 2}
 %   Let $H$ be an abelian group, $K$ be an extraspecial $p$-group and $G=H\rtimes_{\rho} K$. If $Z(G)=1$ then $\ker{\rho}=1$.
%\end{lemma}

%\begin{proof}

We have $\ker{\rho}\cap Z(K)\leq Z(G)=1$. Then $\ker{\rho}=1$ since $K$ is nilpotent.
\end{proof}



Let us show that matrices of order $2$ in $GL_2(p)$ for $p>2$ are diagonalizable.
%Recall that an invertible matrix in $GL_2(p)$ of order $2$ is diagonalizable (see \cite[Lemma 7.1]{LSS}).

\begin{lemma}\label{diag nxn}
    Let $A\in GL_n(p)$ of order $k$ and $p$ be an odd prime. Then: 
    
    \begin{itemize}
        \item[(i)] the $A$-irreducible representations has dimension $\leq k$.
%\end{lemma}

%\begin{corollary}
 \item[(ii)]    If $k=2$ then $A$ is diagonalizable.
%\end{corollary}
    \end{itemize}
\end{lemma}
\begin{proof}

(i)    Let $v\in \Z_p^n$. Then $\setof{f^i(v)}{i\in \mathbb{Z}}=\{v, f(v),\ldots,f^{k-1}(v)\}$ generates an $A$-invariant subspace of dimension less than or equal to $k$. So, if $V$ is irreducible, then $V$ has dimension less than or equal to $k$.  

    
(ii) The order of $A$ and of $\Z_p^n$ are relatively prime, so according to Maschke theorem, the space $V=\Z_p^n$ decomposes as $V=\oplus_i V_i$ where $V_i$ are irreducible representations of $A$. According to (i), the irreducible representations of $A$ have dimension $1$ or $2$. Let $V_j$ be of dimension $2$ and $B=A|_{V_j}$. According to \cite[Lemma 7.1]{LSS} $B$ is diagonalizable. Thus, also $A$ is diagonalizable.
%
%The order of $B$ is $2$ and so either it $B$ is diagonal or $Tr(B)=0$ and $\det(B)=-1$. Thus $B$ has characteristic polynomial $\lambda^2-1=(\lambda-1)(\lambda+1)$ and so $B$ is diagonalizable.
\end{proof}




\begin{corollary} \label{ro faithful on Z_p}
    Let $p$ be an odd prime, $G=\Z_p^n\rtimes_\rho \Z_2^m$. If $\rho$ is faithful then $m\leq n$.
\end{corollary}

\begin{proof}
    According to Lemma \ref{diag nxn} the image of $H$ under $\rho$ is a subgroup of the diagonal matrices of order $2$ (see Lemma \ref{diag nxn}(ii)), which has size $2^n$. Thus, $m\leq n$. 
\end{proof}

%\begin{corollary}
 %   Let $G=\Z_p^m\rtimes_\rho \Z_2^n$. If $\rho$ is faithful then $n\leq m$.
%\end{corollary}

%\begin{proof}
 %   The image of $\rho$ is given by diagonalizable commuting matrices of order $2$ (see Lemma \ref{diag nxn}(ii)). Then, the image of $\rho$ is in the subgroup of the diagonal matrices of order $2$ that has size $2^m$.
%\end{proof}

 

%\begin{lemma}
%Let $A\in GL_n(p)$. Then $C_{GL_n(p)}(A)$ permutes the $A$-invariant subspaces. In particular, the eigenspaces of $A$ are $C_{GL_n(p)}(A)$-invariant.. 
%\end{lemma}
%\comment{not used}
%\begin{proof}
%Let $A(V)=V$ and $[B,A]=1$. Then $AB(V)=BA(V)=V$, so $B(V)$ is $A$-invariant. If $V$ is the eigenspace relative to $\lambda$, i.e. $A(V)=\lambda V$ then $AB(V)=BA(V)=\lambda B(V)$, i.e. $B(V)$ is the eigenspace relative to $\lambda$.
%\end{proof}



%\begin{corollary}\label{center = -1}
%Let $K$ be an extra-special $2$-group, $p>2$ and $G=\Z_p^n\rtimes_\rho K$. If $Z(G)=Z(\Z_p^n\rtimes \gamma_1(K))= 1$ then $\rho_z=-Id_n$ for $1\neq z\in Z(K)$ and $n$ is even.   %
%\end{corollary}
%\comment{and $p=1\pmod{4}$ since $e$ divides $n$.}

%\comment{if we take a group with $\gamma_1(K)\leq Z(K)\cong \Z_2^m$. If we assume that $Z(G)=Z(Z_p^n\rtimes Z(K))=1$ then the same conclusion holds - To check}

%\begin{proof}
%According to Lemma \ref{no center} $\ker{\rho}=1$. Let $1\neq z\in Z(K)$. Then $\rho_z\neq 1$ and it is diagonalizable with eigenvalues $\lambda\in \{1,-1\}$ according to Lemma \ref{diag nxn}. If $V$ is the eigenspace relative to $\lambda=1$ then $V\leq Z(\Z_p^n\rtimes_\rho Z(K))$. Therefore $\rho_z=-Id_n$. The map $\det:K\longrightarrow \Z_p^*$ mapping $x\in K$ to $\det(\rho_x)$ is a group homomorphisms into an abelian group. Then $Z(K)=[K,K]\leq \ker{\det}$ and so $\det(\rho_z)=(-1)^n=1$. Therefore $n$ is even.
%\end{proof}


\begin{corollary}\label{center = -1}
Let $p$ be an odd prime, $K$ be a $2$-group of nilpotency class $2$ and $G=\Z_p^n\rtimes_\rho K$ where $\rho$ is faithful. If $\gamma_1(K)=\langle z\rangle \cong \Z_2$ and $Z(\Z_p^n\rtimes \gamma_1(K))= 1$ then $\rho_z=-Id_n$ and $n$ is even.    
\end{corollary}
%\comment{and $p=1\pmod{4}$ since $e$ divides $n$.}


\begin{proof}
The order of $z$ is $2$ and so $\rho_z\neq 1$ and it is diagonalizable with eigenvalues $\lambda\in \{1,-1\}$ according to Lemma \ref{diag nxn}(ii). If $V$ is the eigenspace relative to $\lambda=1$ then $V\leq Z(\Z_p^n\rtimes_\rho \gamma_1(K))=1$. Therefore, $\rho_z=-Id_n$. The map $\det:K\longrightarrow \Z_p^*$ mapping $x\in K$ to $\det(\rho_x)$ is a group homomorphisms into an abelian group. Then $\gamma_1(K)\leq \ker{\det}$ and so $\det(\rho_z)=(-1)^n=1$. Therefore, $n$ is even.
\end{proof}





%\subsection{Representations of Extraspecial $2$-groups of order $32$}

%Let $K$ be an exstraspecial $2$-group of order $32$. Then $Z(K)=[K,K]=\{1,z\}$ has size $2$ and $u^2\in Z(K)$ for every $u\in K$. In particular, either $u^2=1$ or $u^2=z.$ There are two such groups, namely $K_i=$SmallGroup(32,i) for $i=49,50$
%The group SmallGroup(32,50) is an extraspecial $2$-group of size $32$.% given by the central product of the quaternion group $Q_8$ and the dihedral group $D_8$. 

%\comment{This implies that $GL_2(p)$ has no subgroup isomorphic to SmallGroup(32,50).}

%\subsection{Representations of SmallGroup(32,49)}

%A presentation of the group $K_{49}$ is
%\begin{align*}
%    K_{49}=\langle a,x,b,y\, |\, a^2=x^2=b^2=y^4=(ba)^2=(yx)^2=(yb)^4=(xay)^2=bxybaxay=(ya)^4=1\rangle.
%\end{align*}
%Moreover, $Z(K)=[K,K]=\{1,z\}$ has size $2$, $|b|=|c|=4$ and $u^2\in Z(K)$ for every $u\in K$. 
%In particular, either $u^2=1$ or $u^2=z$. 


%In particular $ba=(ba)^{-1}=a^{-1}b^{-1}=ab$, and so $[a,b]=[x,y]=1$. 


%\begin{lemma}\label{no center 3}
 %   Let $\rho:K\longrightarrow GL_2(p)$ be a faithful representation and $H=\Z_p^2\rtimes_\rho K$. If $H'=\Z_p^2\rtimes_\rho Z(K)$ has trivial center then $\rho(z)=-Id_2$ for $1\neq z\in Z(K)$.
%\end{lemma}

%\comment{questo ora segue da Corollary \ref{center = -1}}
%\begin{proof}
 %   The group $K$ has a center of size $2$. Then $\rho(z)$ is diagonalizable since it has order $2$ (see Lemma 7.1 \cite{LSS}). The eigenvalues of $\rho(z)$ belong to $\{1,-1\}$. If there exists $v\in \Z_p^2$ such that $\rho_z(v)=v$ then $v\in Z(H')=1$. Thus, $\rho_z=-Id_2$.
%\end{proof}




%\begin{proposition} \label{no reps}
 %   Let $p$ be an odd prime, $G=\Z_p^2\rtimes_\rho K_{49}$ where $\rho$ is a faithful action. Then $Z(\Z_p^n\rtimes_\rho Z(K_{49}))\neq 1$.
%\end{proposition}



%\begin{proof}
%Assume that $Z(\Z_p^n\rtimes Z(K_{49}))= 1$ then $\rho_z=-Id_2$ according to Lemma \ref{center = -1}. Since $a^2=b^2=1$ then $abab=[a,b]=1$. Therefore the matrices $\rho_a$ and $\rho_b$ are simultaneously diagonalizable according to Lemma \ref{diag nxn}. The action $\rho$ is faithful, and so $\rho_a$ and $\rho_b$ has eigenvalues $\pm 1.$ Then $\rho_{ab}=\rho_a \rho_b=\rho_z$ and so $z=ab$. The element $ab$ is not central, contradiction.
%  Therefore we have that $b=xax$. Therefore
%
%\begin{align*}
 %   1&=bxybaxay=xaxxyxaxaxay=y^2,
%\end{align*}
 %   contradiction.
%    The matrix $\rho_a$ has order $2$ and so it is diagonalizable according to Lemma \ref{diag nxn}. The action $\rho$ is faithful and $\rho_z=-Id_2$ according to Corollary \ref{center = -1}. So, we can assume that:
%\begin{align*}
 %   \rho_a=\begin{bmatrix}
  %  1 & 0\\
  %  0 & -1
%\end{bmatrix}.%,\quad \rho_b=\begin{bmatrix}
  %  -1 & 0 \\
  %  0& 1
%\end{bmatrix}.
%\end{align*}
%Moreover, we have $\rho_{(ay)^2}=(\rho_a\rho_y)^2=\rho_a \rho_y\rho_a\rho_y=\rho_z$, since $ay$ has order $4$. On the other hand, $\rho_{[a,y]}=[\rho_a,\rho_y]=\rho_a\rho_y\rho_a\rho_y^{-1}=\rho_z$ since 
%
%$=\rho_z=-Id_2$. Then, necessarily
%\begin{align*}
 %   \rho_y=\begin{bmatrix}
  %      0 & y^{-1}\\
   %     y & 0
   % \end{bmatrix}.
%\end{align*}
%In this case, $\rho_y^2=1$, contradiction.
%\end{proof}
Let us conclude the section by showing some applications of the previous results to some concrete semidirect product, where the normal summand is an elementary abelian $p$-group and the other summand is a $2$-group.

Let $K_{50}$ be extraspecial $2$-group labeled by {\tt SmallGroup(32,50)} in the {\tt SmallGroup} library of GAP. A presentation of $K_{50}$ is the following:
\begin{align*}
    K_{50}=\langle a,b,c,d\,|\,a^2,\,b^4,\,c^4 d^2,\, cbcb^{-1},\,db^{-1}db,\,c^{-1}aca,\,dc^{-1}dc,\,dadb^2a\rangle.
\end{align*}



%We are interested in studying faithful representation of $G$ of dimension $2$ over finite fields of prime order. In particular we focus on faithful representation defining a semidirect product $\Z_p^2\rtimes_\rho K$ such that the subgroup $\Z_p^2\rtimes_\rho Z(K)$ has trivial center.



%\begin{lemma}\label{no center}
 %   Let $\rho:K\longrightarrow GL_2(p)$ be a faithful representation and $H=\Z_p^2\rtimes_\rho K$. If $H'=\Z_p^2\rtimes_\rho Z(K)$ has trivial center then $\rho(z)=-Id_2$ for $1\neq z\in Z(K)$.
%\end{lemma}
%\comment{questo ora segue da Corollary \ref{center = -1}}
%\begin{proof}
 %   The group $K$ has a center of size $2$. Then $\rho(z)$ is diagonalizable since it has order $2$ (see Lemma 7.1 \cite{LSS}). The eigenvalues of $\rho(z)$ belong to $\{1,-1\}$. If there exists $v\in \Z_p^2$ such that $\rho_z(v)=v$ then $v\in Z(H')=1$. Thus, $\rho_z=-Id_2$.
%\end{proof}
%Let us define:
%\begin{itemize}
  %  \item[(i)] 
%$    \alpha_\pm (t)=\begin{bmatrix}
 %       0 & \pm t\\
  %      t & 0
  %  \end{bmatrix}$, 
   % for $t^2=\mp 1$. 
   % \item[(ii)]
  %  $$\beta (s)=\begin{bmatrix}
   %     s & 0\\
 %       0 & -s
 %   \end{bmatrix}, \quad s^2=-1$$.
%\end{itemize}

%Note that $\beta(s)$ and $\alpha_+(t)$ are defined only if $p=1\pmod{4}$. Moreover, if $\rho_x$ and $\rho_a$ are defined as in \eqref{rho D_8} we have:
%\begin{itemize}
 %   \item[(i)] $[\rho_x,\alpha_{+}(t)]=-Id_2$ and $[\rho_a,\alpha_+(s)]=-Id_2$.
  %  \item[(ii)] $[\rho_x,\alpha_{-}(t)]=-Id_2$ and $[\rho_a,\alpha_-(s)]=1$.
   % \item[(ii)] $[\rho_x,\beta(s)]=1$ and $[\rho_a,\beta(s)]=-Id_2$.
%\end{itemize}


\begin{proposition}\label{no rep}
    Let $p$ be an odd prime, $G=\Z_p^2\rtimes_\rho K_{50}$ where $\rho$ is a faithful action. Then $Z(\Z_p^2\rtimes_\rho Z(K_{50}))\neq 1$.
\end{proposition}

    \begin{proof}
 Note that $Z(K_{50})=\gamma_1(K_{50})=\langle b^2\rangle$ and that it has order $2$. Assume that $Z(\Z_p^n\rtimes Z(K_{50}))= 1$ then $\rho_{b^2}=-Id_2$ according to Corollary \ref{center = -1}. Since $|a|=2$, $\rho_a$ is diagonalizable (see Lemma \ref{diag nxn}) with eigenvalues equal to $\pm 1$ and $\rho_a\neq \pm Id_2$. In particular, we can assume that 
$$\rho_a=\begin{bmatrix}
1 & 0\\
0 & -1
\end{bmatrix}.$$ 
Since we have that $[a,c]=1$, $c^2=b^2\in Z(K_{50})$ and $c\notin Z(G)$ we have that $[\rho_a,\rho_c]=1$ and $\rho_a^2=\rho_b^2=\rho_z=-Id$. So necessarily  $$\rho_c=\begin{bmatrix}
    x &0 \\
    0 & -x
\end{bmatrix}$$
for $x^2=-1$. Since $[c,d]=1$ we have $[\rho_c,\rho_d]=1$ and so $\rho_d$ must be a diagonal matrix. Therefore, $[\rho_d,\rho_a]=\rho_{[a,d]}=1$, but $[a,d]\neq 1$, which is a contradiction. 
\end{proof}

\subsection{The Group \texorpdfstring{$\mathcal{G}_k\times \Z_2^2$}{GkxZ2xZ2}}\label{grup for sr}

Let $p=1\pmod{3}$ be a prime number and let $k\in \Z_p$ such that $k^3=1$ and $k\neq 1$. Let us consider the following presentation of the group $\mathcal{Q}_8$:
$$\mathcal{Q}_8=\langle x,y,z\,|\, x^2,\, y^2,\, [x,y]z^{-1},\, z^2,\,[z,y],\,[z,x]\rangle.$$

As defined in \cite[Section 7]{LSS}, we have the group $\mathcal{G}_k=\mathbb{Z}_p^2\rtimes_{\rho} \mathcal{Q}_8$, where the action $\rho$ is defined by setting
\begin{displaymath}
\rho_x=\begin{bmatrix}
0 & -1\\
1 & 0
\end{bmatrix},\quad
\rho_{y}=\begin{bmatrix}
k^2 & k\\
k & -k^2
\end{bmatrix},\quad \rho_z=\begin{bmatrix}
-1 & 0\\
0 & -1
\end{bmatrix}.
\end{displaymath}
%Faitfhul representation of degree $2$ of non-abelian groups are irreducible. According to \cite{Mayr}, $\rho$ is the unique irreducible action of $\mathcal{Q}_8$ of degree $2$, and moreover it is a fixed-point-free action, i.e. $Fix(\rho_h)=\setof{v\in \mathbb{Z}_p^2}{\rho_h(v)=v}=0$ for every $h\in \mathcal{Q}_8$. Equivalently $1-\rho_h\in GL_2(p)$ for every $h\in \mathcal{Q}_8$. In particular, $\gamma_1(\mathcal{G}_k)=\mathbb{Z}_p^2\rtimes Z(\mathcal{Q}_8)$ and $\gamma_2(\mathcal{G}_k)=\mathbb{Z}_p^2\times \{1\}$ is the $p$-Sylow subgroup of $\mathcal{G}_k$. The action $\rho$ has no invariant subspaces, so $\gamma_2(\mathcal{G}_k)$ is a minimal normal subgroup of $\mathcal{G}_k$.

Note that $Z(\mathcal{G}_k)=1$, so the center of the group $\mathcal{G}_k\times A$ where $A$ is an abelian group is the subgroup $1\times A$ and $\gamma_n(\mathcal{G}_k\times A)=\gamma_n(\mathcal{G}_k)\times 1$ for every $n\geq 1$.
\begin{lemma}
    The automorphisms of the group $G=\mathcal{G}_k\times \Z_2^2$ are of the form
    \begin{align}\label{aut of G_k x Klein}
        f=\begin{cases}
          x\mapsto \widetilde{f}(x)z_1,\\
                  y\mapsto \widetilde{f}(y)z_2,\\
                  z\mapsto \widetilde{f}(z)\\   f|_{\gamma_2(G)}=\widetilde{f}|_{\gamma_2(\mathcal{G}_k)},\\
                  f|_{Z(G)}=F,
        \end{cases}
    \end{align}
    where $\widetilde{f}\in \aut{\mathcal{G}_k}$, $z_1,z_2\in Z(G)$ and $F\in \aut{Z(G)}$. In particular $|\aut{G}|= 2^8 \cdot 3^2\cdot p^2(p-1)$.
    %In particular, the maps defined by setting:
    %\begin{align*}
     %   \aut{G}&\longrightarrow \aut{G_k},\quad f\mapsto f_{Z(G)},\\
     %   \aut{G}&\longrightarrow \aut{Z(G)},\quad f\mapsto f|_{Z(G)},
   % \end{align*}
   % are surjective group morphisms.
\end{lemma}

\begin{proof}
  Let $f\in \aut{G}$ then clearly $f$ is as in \eqref{aut of G_k x Klein} since $Z(G)$, $\gamma_1(G)$ and $\gamma_2(G)$ are characteristic subgroups. On the other hand, let $f$ be defined as in \eqref{aut of G_k x Klein}. We must check that all relations in the presentation of $G$ are preserved by $f$. 
  \begin{align*}
      [f(x),f(y)]&=[\widetilde{f}(x)z_1,\widetilde{f}(x)z_2]=[\widetilde{f}(x),\widetilde{f}(y)]=\widetilde{f}([x,y])=\widetilde{f}(z)=f(z)=f([x,y]).\\
      f(x)^4&=\widetilde{f}(x)^4 z_1^4=\widetilde{f}(x)^4=1=f(x^4)\\
      f(y)^4&=\widetilde{f}(y)^4 z_2^4=\widetilde{f}(y)^4=1=f(y^4)\\
%      f(z)^2&=\widetilde{f}^2=1,\\
      f(v^p)&=f(v)^p=f(z^2)=f(z)^2=1, \\
      f(x) f(v) f(x)^{-1}&=\widetilde{f}(x)^{-1}z_1\widetilde{f}(v)z_1^{-1}\widetilde{f}(x)^{-1}=\widetilde{f}(xvx^{-1})=f(xvx^{-1}),\\
      f(y) f(v) f(y)^{-1}&=\widetilde{f}(y)^{-1}z_2\widetilde{f}(v)z_2^{-1}\widetilde{f}(y)^{-1}=\widetilde{f}(yvy^{-1})=f(yvy^{-1}),\\
    [f(z),f(t)]&=1=f([x,t]),
  \end{align*}
for every $v\in \gamma_2(G)$, $z\in Z(G)$ and $t\in G$. Therefore, $f\in \aut{G}$. The second statement is straightforward, using that $|\aut{\mathcal{G}_k}|=24p^2(p-1)$ (see Proposition 7.5 of \cite{LSS}) and $|\aut{\Z_2^2}|=6$.
%  given the pair $(g,h)\in \aut{G_k}\times \aut{Z(G)}$ then clearly the map defined by setting $f(xz)=g(x)h(z)$ for every $x\in G_k$ and $z\in Z(G)$ defines an automorphism of $G$. 
\end{proof}


Let us define $\mathcal{F}=\setof{h\in \aut{\mathcal{G}_k}}{|Fix(h)|=2p,\, |h|=3}$ and $\mathcal{H}=\setof{h\in \aut{G}}{h_{Z(G)}\in \mathcal{F},\, |h|_{Z(G)}|=3}$. According to \cite[Lemma 7.7]{LSS} $|\mathcal{F}|=16p$ and consequently, $|\mathcal{H}|=16p\cdot 16\cdot 2=2^9p$. The representatives of conjugacy classes of the elements in $\mathcal{F}$ are 
\begin{align} 
        f(j)=\begin{cases}
          x\mapsto y,\\
                  y\mapsto xy,\\
                  z\mapsto z\\   f|_{\gamma_2(\mathcal{G}_k)}=F_j=\begin{bmatrix}
                      -k(1+k^j) & -(1+k^j)\\
                      0 & -k^2(1+k^j)
                  \end{bmatrix},
       \end{cases}
    \end{align}
for $j=1,2$, see the proof of \cite[Theorem 5.5]{LSS}. Let us also define
\begin{align*}
    M=\begin{bmatrix}
        0& 1\\
        1 & 1
    \end{bmatrix}\in GL_2(2).
\end{align*}
    
\begin{lemma}\label{iso classes of sr}
   The representatives of the conjugacy classes of elements of $\mathcal{H}$ are of the form
   \begin{align*}%\label{aut of G_k x Klein_2}
        f(j,a)=\begin{cases}
          x\mapsto ya,\\
                  y\mapsto xy ,\\
                  z\mapsto z\\   f|_{\gamma_2(G)}=F_j,\\
                  f|_{Z(G)}=M,
        \end{cases}
    \end{align*}
   for $j=1,2$ and $a\in \{(0,0),(1,0)\}$.
\end{lemma}

\begin{proof}
Note that $f(j,a)_{Z(G)}=f(j)$ as defined in \eqref{fj} for $j=1,2$. Clearly the representatives listed in the statement are not conjugate. Indeed if $f(l,a)$ and $f(k,a')$ are conjugate then $f(l,a)_{Z(G)}=f(l)$ and $f(k,a')_{Z(G)}=f(k)$ are conjugate in $\aut{\mathcal{G}_k}$ and so $l=k$. Moreover $|f(l,(0,0))|=3$ and $|f(l,(1,0))|=6$, so they are not conjugate.

Let us compute the centralizer of $f(j,a)$ in $\aut{G}$. If $h\in C(f(j,a))$ then $h|_{Z(G)}$ centralizes $M$ and so $h|_{Z(G)}\in \langle M\rangle$. Moreover $h_{Z(G)}$ centralizes $f(j)$ in $\aut{\mathcal{G}_k}$ and therefore we have that $h$ is one of the following:
\begin{align*}\label{aut of G_k x Klein}
        h_1&=\begin{cases}
          x\mapsto xv_1 z_1,\\
                  y\mapsto y v_2 z_2',\\
                  z\mapsto z w\\   h_1|_{\gamma_2(G)}=\begin{bmatrix}
                      x & 0\\
                      0 & x
                  \end{bmatrix},\\
                  h_2|_{Z(G)}=M^s
       \end{cases},\qquad
       h_2&=\begin{cases}
          x\mapsto xy v_1 z_1',\\
                  y\mapsto x v_2 z_2',\\
                  z\mapsto z w\\   h_2|_{\gamma_2(G)}=\begin{bmatrix}
                      x & -xk\\
                      0 & xk^2
                  \end{bmatrix},\\
                  h_2|_{Z(G)}=M^s
       \end{cases},\qquad
        h_3&=\begin{cases}
          x\mapsto y v_1 z_1',\\
                  y\mapsto xy v_2 z_2',\\
                  z\mapsto z w\\   h_3|_{\gamma_2(G)}=\begin{bmatrix}
                      x & xk\\
                      0 & xk^2
                  \end{bmatrix},\\
                  h_2|_{Z(G)}=M^s
       \end{cases}
    \end{align*}
    for $x\neq 0$, $v_1,v_2,w\in \Z_p^2$, $z_1',z_2'\in \Z_2^2$ and $s=0,1,2$. It is easy to check that if $a=(0,0)$ then $s$ and $z_1'$ can be chosen freely and $z_2'$ is uniquely determined in all three cases. Therefore, $|C(f(j,(0,0),(0,0)))|=|C(f_j)|\cdot 3\cdot 4$. If $z_1=(1,0)$ then in all three cases $s$ and $z_2'$ are uniquely determined and $z_1'$ is free. Therefore $|C(f(j,(1,0),(0,0)))|=|C(f_j)|\cdot 4$. The set $\mathcal{H}$ is invariant under the action by $\aut(\mathcal{G}_k \times \Z_2^2)$ and so the conjugacy class of $f(j,a)$ lies in $\mathcal{H}$. So, given that $|C(f(j,a)_{Z(G)})|=3p(p-1)$ (see \cite[Lemma 7.8]{LSS}) we have that
    \begin{align*}
    2\left(\frac{|\aut{G}|}{|C(f(j,(0,0),(0,0)))|}+\frac{|\aut{G}|}{|C(f(j,(1,0),(0,0)))|}\right)=2( 2^6p+3\cdot 2^6 p)=2^9p=|\mathcal{H}|.    
    \end{align*}
    Therefore the automorphisms in the statement are a set of representatives of conjugacy classes of elements of $\mathcal{H   }$.
    \end{proof}


%\comment{From this can I deduce that some class of $\LSS(p,16,\sigma_Q)$ is empty?}



%
%\begin{proposition}\label{action is diagonal}
 %   Let $G=\Z_p^t\rtimes_\rho \Z_4^2$ where $t\in \{2,3\}$. If $Z(G)=1$ then $\rho_x$ is diagonalizable for every $x\in \Z_4^2$.% and $p=1\pmod{4}$.
%\end{proposition}
 
        %\begin{proof}
%According to Lemma \ref{no center} $\ker{\rho}=1$. Let $K=\Z_4^2$ and $x\in K$. Then $\rho_x^2$ has order $2$ and so it is diagonalizable by Lemma \ref{diag nxn}.
% The group $K$ is abelian, then all the squares of elements of $K$ are diagonal with respect to some basis. %If $\rho_x$ is irreducible then $\rho_x$  $\rho_x^2$ is not diagonal. 

% If $t=2$ then $\rho_x$ commutes with a diagonal matrix with eigenvalues $\pm 1$. Then $\rho_x$ is diagonal. 
 
% Let $t=3$. The set $U=\setof{\rho_x^2}{x\in K}$ is a subgroup of the diagonal matrices with eigenvalues $\pm 1$ and $|U|=4$. We have that $\rho_x^2\in U$, $\rho_x$ commutes with the elements of $U$ and according to Corollary \ref{case q=2 then n even} we have that $p=1\pmod{4}$. By a direct computation we have that for every choice of the subgroup $U$ the matrix $\rho_x$ is diagonalizable. 
 %
% Therefore since $|\rho_x|=4$ the eigenvalues of $\rho_x$ has order $4$. Hence, $p=1\pmod{4}$.
%\end{proof}

\section*{Acknowledgments}
\noindent The author Filippo Spaggiari was supported by GAUK (301-10/252012, project number 166122) and the GAČR grant no. 22-19073S.




\bibliographystyle{amsalpha}
\bibliography{references} % References file


\end{document}


\subsection{Sylow subgroups of \texorpdfstring{$GL_n(p)$}{GLnp}} \label{app:syl_of_GL}



In this section, we explicitly describe the $q-$ Sylow subgroups of $GL_n(p)$ following the classical results of \cite{carter1964sylow} and \cite{weir1955sylow} summarized in \cite{schmidtsylow}. We denote by $\syl_q(G)$ the $q$-Sylow of a group $G$ and by $\mathbb{F}_{p^u}$ the finite field of $p^u$ elements. We set 
\begin{align*}
e(p,q)&=\min\setof{k\in \mathbb{N}} {p^k=1\pmod{q}} \text{ if } q>2,\\
e(p,2)&=\min\setof{k\in \mathbb{N}} {p^k=1\pmod{4}}.
\end{align*}
We can assume that $e(p,q)<n$, otherwise $q$ does not divide $|GL_n(p)|$. Thus, we can set $n=e(p,q)s+r$ for some $0\leq r<e(p,q)$. 

We have the following cases.

\begin{itemize}
    \item[(1)] Assume that $q>n$ and $q\neq 2$. Then $\syl_q(GL_n(p))$ is isomorphic to the following group: 
\begin{align*}
T_{p,q,n}=
    \begin{bmatrix}
        \boxed{\boldsymbol{\zeta_1}} &  \mathbf{0}_{e(p,q)} & \mathbf{0}_{e(p,q)} & \mathbf{0}_{e(p,q)} \\
        \mathbf{0}_{e(p,q)} & \ddots &  \mathbf{0}_{e(p,q)} & \mathbf{0}_{e(p,q)} \\
         \mathbf{0}_{e(p,q)} &  \mathbf{0}_{e(p,q)} & \boxed{\boldsymbol{\zeta_s}} & \mathbf{0}_{e(p,q)} \\
        \mathbf{0}_{e(p,q)} & \mathbf{0}_{e(p,q)} & \mathbf{0}_{e(p,q)} & \boxed{\mathbf{1}_r}
    \end{bmatrix}
\end{align*}



    where $\boldsymbol{\zeta_i}$ denotes the $\Z_p-$linear map given by the multiplications by $\zeta_i\in \syl_q(\mathbb{F}_{p^{e(p,q)}}^*)$ in $\mathbb{F}_{p^{e(p,q)}}$. Therefore, $T_{p,q,n}$ is given by matrices with $s$ blocks of dimension $e(p,q)$ and $r$ blocks of dimension $1$ that are fixed by all elements $T_{p,q,n}$. %\comment{Note that in this case $\syl_q(GL_n(p))$ is abelian and $|T_{n,e}|=t^s$ where $t=\max\setof{m\in \mathbb{N}}{p^e=1\pmod{q^m}}$. - SERVE?}

 \item[(2)] Assume that $q\leq n$ and $q\neq 2$. In  this case we have:
 \begin{align*}
\syl_q(GL_n(p))=T_{p,q,n}\rtimes \syl_q(\sym_{s})    
\end{align*}
where $\syl_q(\sym_{s})$ acts on the first $e(p,q)s$ columns by permuting the sets of $s$ adjacent columns.


 
   \item[(3)] Assume that $q=2$ and $n\geq 2$. Note that $e(p,2)\in \{1,2\}$. Then
    \begin{align*}
\syl_2(GL_n(p))=(\syl_2(\mathbb{F}_{p^e})\rtimes \syl_2(\langle \varphi_p\rangle))\rtimes \syl_2(\sym_{s})    
\end{align*}
where $\varphi_p$ denotes the $\mathbb{Z}_p$-linear map on $\mathbb{F}_{p^e}$ given by the Frobenius automorphism $x\mapsto x^p$. In particular:

\begin{align*}
\syl_2(GL_n(p))\ni M=\left(
    \begin{bmatrix}
        \boxed{\boldsymbol{\zeta_1}\boldsymbol{\varphi}_p^{k_1}} &  \mathbf{0}_{e(p,q)} & \mathbf{0}_{e(p,q)} & \mathbf{0}_{e(p,q)} \\
        \mathbf{0}_{e(p,q)} & \ddots &  \mathbf{0}_{e(p,q)} & \mathbf{0}_{e(p,q)} \\
         \mathbf{0}_{e(p,q)} &  \mathbf{0}_{e(p,q)} & \boxed{\boldsymbol{\zeta_s}\boldsymbol{\varphi}_p^{k_s}} & \mathbf{0}_{e(p,q)} \\
        \mathbf{0}_{e(p,q)} & \mathbf{0}_{e(p,q)} & \mathbf{0}_{e(p,q)} & \boxed{\mathbf{1}_r}
    \end{bmatrix},\sigma\right)
\end{align*}
where $\zeta_i\in \syl_2(\mathbb{F}_{p^e}^*)$, $\varphi_p^{k_j}\in \syl_2(\langle \varphi\rangle)$ and $\sigma\in \syl_2(\sym_s)$.
\end{itemize}



%\comment{Case $q=2$ and $p>2$: then $p=1\pmod{4}$ or $p^2=1\pmod{4}$. So, if $q=2$ and $p>2$ we have that $e=1$ or $e=2$. How does the $2$-Sylow subgroups look like in the two cases? If $e=1$ then $\varphi_p$ is trivial and the matrices of $T_e$ are diagonal.}

This description allows us to have some restrictions on $n$ adding some further assumptions on the group of the form $\Z_p^n\rtimes_\rho K$ where $K$ is a $q$-group.

\begin{proposition}\label{trivial center}
    Let $p,q$ be primes, $K$ is a $q$-group and $G=\Z_p^n\rtimes_\rho K$. If $Z(G)=1$, then $e(p,q)$ divides $n$.
\end{proposition}

\begin{proof}
Let $n=e(p,q)s+r$.  The group $K$ embeds into $T_{p,q,n}$ or into the group $T_{p,q,n}\rtimes Syl_q(\sym_s)$. If $r>0$ then the action of $K$ is trivial on a subspace of dimension $r$. Then $K$ centralizes a subspace of dimension $r$ of $\Z_p^n$. Such subspace commute also with the subgroup $\Z_p^n$ and so it is contained in $Z(G)=1$. Therefore, $r=0$. 
\end{proof}







\begin{corollary}\label{case q=2 then n even}
    Let $p>2$ be a prime such that $e(p,2)=2$, $K$ be a $2$-group and $G=\Z_p^n\rtimes_\rho K$. If $Z(G)=1$ then $n$ is even.
\end{corollary}

\begin{proof}
    It is enough to apply Proposition \ref{trivial center}.
\end{proof}


\subsection{The group $G_{M,N}(t)$}\label{Sec: G_{M,N}}
\comment{non serve più tutta la section}
Let $M,N\in GL_t(\Z_p)$ be a pair of commuting and diagonalizable matrices of order $4$. We denote by $G_{M,N}(t)$ the group $\Z_p^t\rtimes_\rho \Z_4^2$ where $\rho$ is a group action such that $\rho(1,0)=M$ and $\rho(0,1)=N$. %For $v=(r,s)\in\Z_4^2$ we denote by $T_v=T_{(r,s)}=M^rN^s$, so that $(u,v)\ast(x,y)=(u+T_v(x),v+y)=(u+M^rN^s(x),v+y)$.


\begin{theorem}\label{iso G_{A,B}(t)}
    The map $f:G_{M,N}(t)\longrightarrow G_{A,B}(t)$ defined by setting:
    \begin{align*}
        f(0,(1,0))=(w,F(1,0)),\quad 
        f(0,(0,1))=(v,F(0,1)),\quad f(u,(0,0))=(G(u),(0,0))
    \end{align*}
    where $u,w,v\in \Z_p^{t}$, $F=\begin{bsmallmatrix}r&h\\s&k\end{bsmallmatrix}\in \aut{\Z_4^{2}}$ and $G\in GL_t(\Z_p)$, can be extended to a group isomorphisms between $G_{M,N}(t)$ and $G_{A,B}(t)$ if and only if $GMG^{-1}=A^rB^s$ and $GNG^{-1}=A^hB^k$. %In particular, $G_{M,N}\cong G_{A,B}$ if and only if the subgroups $\langle M,N\rangle$ and  $\langle A,B\rangle$ are conjugate in $GL_t(\Z_p)$.
\end{theorem}

\comment{Data una mappa definita sui generatori. La sua estensione "lineare" è ben definita se mappa relazioni in relazioni.}




\begin{proof}
    Let $\setof{(u_i,0), (1,0), (0,1)}{i=0,\ldots, t}$ be a set of generators of $G_M$. 

    Assume that the map $f:G_{M,N}(t)\longrightarrow G_{A,B}(t)$ can be extended to an isomorphism. Then $f$ maps the unique $p$-Sylow subgroup of $G_{M,N}(t)$ onto the unique $p$-Sylow subgroup of $G_{A,B}(t)$. In particular the restriction of $f$ to the $p$-Sylow subgroup is $G\in GL_{t}(\Z_p)$. Assume that the image of $(0,(1,0))$ under $f$ is $(w,(r,s))$. Then
    \begin{align*}
        f((0,(1,0)) (u,(0,0))(0,(1,0))^{-1}) 
        &=(G M(u),0)=f((0,(1,0)) f(u,(0,0)) f((0,(1,0))^{-1} \\
        &=(w,(r,s))(G(u),(0,0))(w,(r,s))^{-1}=(T_{(r,s)}G(u),(0,0))\\
        &=(A^rB^sG(u),(0,0))
    \end{align*}
    therefore $GMG^{-1}=A^rB^s$. In the same way, if the image of $(0,(0,1))$ under $f$ is $(v,(h,k))$. Then we must have $GNG^{-1}=A^hB^k$. The map $f$ induces an automorphism on the factor group $G_{M,N}(t)/\Z_p^t\cong G_{A,B}(t)/\Z_p^t\cong \Z_4^2$ mapping $(1,0)$ to $(r,s)$ and $(0,1)$ to $(h,k)$.
    

    On the other hand, note that $\setof{(G(u_i),0), (0,F(1,0))=(0,(r,s)), (0,F(0,1))=(0,(h,k))}{i=1,\ldots,t}$ is a set of generators of $G_{A,B}(t)$. \comment{In the same way as above, if $GMG^{-1}=A^rB^s$ and $GNG^{-1}=A^hB^k$, then $f$ maps the relations of $G_{N,M}(t)$ to relations of $G_{A,B}(t)$ with respect to these sets of generators. Thus, $f$ can be extended to a well defined group isomorphism.}
\end{proof}


%\begin{corollary}
%Let $t=2$.    There is a unique group of the form $G_{N,M}$ where $N,M$ are diagonal matrices. In particular, if $a$ is the generator of the unique subgroup of order $4$ of $\Z_p^*$ then such a unique group is $G_{N_a, M_a}$ where
%    \begin{align*}
%        N_a=\begin{bmatrix}
 %           a & 0\\
  %          0& 1
   %     \end{bmatrix}, \quad         M_a=\begin{bmatrix}
    %        1 & 0\\
     %       0& a
   % \end{align*}
%\end{corollary}

%\begin{proof}
%Let $N,M$ be diagonal matrix of order $4$. Then $N=N_a^r M_a^s$ and $M=N_a^h M_a^k$. According to ... the map defined by setting 
%\begin{align*}
 %   (0,(1,0))&\mapsto (0,(r,s))\\
  %  (0,(0,1))&\mapsto (0,(h,k))\\
   % (u,(0,0))&\mapsto (u,(0,0)).
%\end{align*}
 %   define an isomorphism between $G_{N,M}$ and $G_{N_a, M_a}$.
%\end{proof}

%\comment{I think that such a group has no good automorphisms! Indeed, if $f\in \aut{G_{N_a,M_a}}$ then if $F$ is either diagonal, and in this case the induced automorphism on $\Z_4^2$ is the identity, or it has non-zero entries on the odd diagonal. In this case $FN_a F^{-1}=M_a$ and $FM_a F^{-1}=N_a$, and so the induced automorphism on $\Z_4^2$ has fixed points different from $1$. Hence necessarily $t=3$.}



\begin{corollary}\label{aut of G_{M,N}}
    The automorphism group of $G_{M,N}(t)$ are the maps defined by setting:
\begin{align*}
        f(0,(1,0))=(w,(r,s)),\quad 
        f(0,(0,1))=(v,(h,k)),\quad f(u,(0,0))=(G(u),(0,0))
    \end{align*}
where $\{(r,s),(h,k)\}$ generates $\Z_4^2$, $G\in GL_t(p)$ such that $GMG^{-1}=A^rB^s$ and $GNG^{-1}=A^hB^k$.% In particular $|\aut{G_M}|=p^{2t}|N_{GL_t(\Z_p)}(\langle M,N\rangle)|$.
\end{corollary}


We are interested in the automorphism of $G_{M,N}(t)$, such that the induced automorphism on the abelianization of $G_{M,N}(t)$ fix only the identity. The existence of such automorphisms provide some restrictions on the structure of the group.


\begin{lemma}\label{determinant}
   Let $G_{N,M}(t)$ where $M$ and $N$ are diagonal matrices of order $4$. If there exists $f\in \aut{G_{M,N}(t)}$ such that $Fix(f_{\Z_p^t})=1$ then $\det(N)=\det(M)=1$.   
\end{lemma}


\begin{proof}
    Let $f\in \aut{G_{M,N}(t)}$, $g=f_{\Z_p^t}$ and $F=f|_{\Z_p^t}$. According to Corollary \ref{aut of G_{M,N}} we have that
    $FMF^{-1}=N^r M^s$, $FNF^{-1}=N^h M^k$ where $g(1,0)=(r,s)$ and $g(0,1)=(h,k)$. In particular $\det(N)=\det(N)^r\det(M)^s$, $\det(M)=\det(N)^h\det(M)^k$ and $\det(N)^4=\det(M)^4=1$. Assume that $Fix(g)=0$. Taking al the possible automorphisms of $\Z_4^2$ with such property, the equations for the determinant imply that $\det(N)=\det(M)=1$.
\end{proof}

\begin{corollary}\label{t=2 forbidden}
   Let $G_{M,N}(2)$ where $M$ and $N$ are diagonal matrices of order $4$. If $\det(M)=\det(N)=1$ then $\rho$ is not faithful.% $$Fix(f_{\Z_p^2})\neq 0$ for every $f\in \aut{G_{M,N}(2)}$.   
\end{corollary}

\begin{proof}
    Let $a$ be an element of order $4$ in $\Z_p^*$. Let $f\in \aut{G_{M,N}(t)}$ and $F=f|_{\Z_p^2}$. The matrices $N,M$ are diagonal. According to Lemma \ref{determinant} $\det(N)=\det(M)=1$, so the matrices have eigenvalues $a,a^{-1}$. Then $N=M^{-1}$ and so the image under $\rho$ has size $4$.
%%   
  %  Then, also $FNF^{-1}$ is diagonal and has eigenvalues $a,a^{-1}$. Thus, either $FNF^{-1}=N$ and $FMF^{-1}=M$ and so the induced automorphism on $\Z_4^2$ is the identity or $FNF^{-1}=M$ and $FMF^{-1}=N$. In this case, we have that $(x,x)\in \Z_4^2$ is fixed by the induced automorphism on $\Z_4^2$.
\end{proof}


%\begin{corollary}\label{t=2 forbidden}
 %  Let $G_{M,N}(2)$ where $M$ and $N$ are diagonal matrices of order $4$. Then $Fix(f_{\Z_p^2})\neq 0$ for every $f\in \aut{G_{M,N}(2)}$.   
%\end{corollary}
%
%\begin{proof}
 %   Let $a$ be an element of order $4$ in $\Z_p^*$. Let $f\in \aut{G_{M,N}(t)}$ and $F=f|_{\Z_p^2}$. The matrices $N,M$ are diagonal. According to Lemma \ref{determinant} $\det(N)=\det(M)=1$, so the matrices have eigenvalues $a,a^{-1}$. Then, also $FNF^{-1}$ is diagonal and has eigenvalues $a,a^{-1}$. Thus, either $FNF^{-1}=N$ and $FMF^{-1}=M$ and so the induced automorphism on $\Z_4^2$ is the identity or $FNF^{-1}=M$ and $FMF^{-1}=N$. In this case, we have that $(x,x)\in \Z_4^2$ is fixed by the induced automorphism on $\Z_4^2$.
%\end{proof}


Let $p$ be an odd prime such that $p=1\pmod{4}$ and $a\in \Z_p^*$ an element of order $4$. Let us define the following matrices:
 \begin{align}\label{def A,B}
       A=\begin{bmatrix}
           a & 0 & 0\\
           0 & 1& 0\\
           0 & 0 & a^{-1}
       \end{bmatrix} , \quad 
       B=\begin{bmatrix}
           1 & 0 & 0\\
           0 & a& 0\\
           0 & 0 & a^{-1}
       \end{bmatrix} 
    \end{align}
\begin{proposition}\label{G_{A,B}}
    Let $G_{M,N}(3)$ where $M$ and $N$ are diagonal matrices of order $4$ with $\det(N)=\det(M)=1$ such that $\langle M,N\rangle\cong \Z_4^2$. Then $G_{M,N}(3)\cong G_{A,B}(3)$ where $A,B$ are defined as in \eqref{def A,B}.
\end{proposition}

%\comment{Credo serva anche che le azioni siano fedeli altrimenti non viene un automorfismo la mappa $(1,0)\mapsto (l,r)$, $(0,1)\mapsto (t,u)$ CHECK!}
\begin{proof}
The eigenvalues of $M$ and $N$ has order dividing $4$, so they are a power of $a$. Thus 
 \begin{align*}
       N=\begin{bmatrix}
           a^l & 0 & 0\\
           0 & a^r& 0\\
           0 & 0 & a^{-l-r}
       \end{bmatrix}=A^l B^r , \quad 
       M=\begin{bmatrix}
           a^t & 0 & 0\\
           0 & a^u& 0\\
           0 & 0 & a^{-t-u}
       \end{bmatrix} =A^t B^u
    \end{align*}
The subgroup generated by $M$ and $N$ has order $16$, so the assignment $(1,0)\mapsto (l,t)$, $(0,1)\mapsto (h,k)$ is an automorphism of $\Z_4^2$ since $M,N$ generate a group isomorphic to $\Z_4^2$.

Setting $f(0,(1,0))=(0,(l,r))$, $f(0,(0,1))=(0,(t,u))$ and $f(u,(0,0))=(u,(0,0))$, we have that $f$ is an isomorphism between $G_{M,N}(3)$ and $G_{A,B}(3)$ according to Lemma \ref{iso G_{A,B}(t)}.
    \end{proof}

%\comment{Se parto dal fatto che i blocchi di $\alpha$ hanno left multiplication di ordine $3$ allora $f|_{\Z_p^3}$ ha ordine $3$. Quindi $FAF^{-1}$ permuta gli automorfismi di $A$ con una permutazione di ordine $3$. Probabilmente questo implica che l'automorfismo indotto su $\Z_4^2$ ha punti fissi.     Alla fine viene che $F$ ha autovalore doppio $1$ e autovalore $w$. Si fa vedere che $w=-1$, ma questo contraddice \ref{ecco!}}




\begin{proposition}\label{chain 4 then p=3}
    Let $G_{A,B}(3)$ where $A,B$ are defined as in \eqref{def A,B}. There is no $f\in \aut{G_{A,B}(3)}$ such that $|Fix(f)|=p^2$ and $Fix(f_{\Z_p^3})=1$. 
\end{proposition}

\begin{proof}
Let $G=G_{A,B}(t)$. According to Corollary \ref{aut of G_{M,N}} then $FAF^{-1}=A^r B^s$ and $FBF^{-1}=A^h B^k$ and the induced map $f_{\Z_p^3}$ defined as $(1,0)\mapsto (r,s)$, $(0,1)\mapsto (h,k)$ is an automorphism of $\Z_4^2$. In particular, $FAF^{-1},FBF^{-1}$ are diagonal matrices with eigenvalues $a,1,a^{-1}$. Accordingly $FAF^{-1},FBF^{-1}\in \{A^{\pm 1}, B^{\pm 1}, (AB^{-1})^{\pm 1}\}$. Moreover, $Fix(f_{\Z_p^3})=1$, so the only choices allowed are $FAF^{-1}=B^{-1}$, $FBF^{-1}=AB^{-1}$ or $FAF^{-1}=AB^{-1}, FBF^{-1}=A^{-1}$ and accordingly $F$ is 
\begin{align}\label{F}
    F=\begin{bmatrix}
        0 & x & 0\\
        y & 0 & 0\\
    0& 0 & z
    \end{bmatrix}
\end{align}
where $\det(F)=-xyz\neq 0$. On the other hand, the matrix $F$ has an eigenspace relative to $1$ of dimension $2$ and so $F$ is similar to a matrix of the form
\begin{align}\label{F'}
    F'=\begin{bmatrix}
        1 & 0 & t\\
        0 & 1 & u\\
    0& 0 & w
    \end{bmatrix}
\end{align}
for some $w\neq 0$.
The characteristic polynomials of $F$ and $F'$ coincide, so we have
\begin{align*}
  k^3- xyz=(k-1)^2(k-w)=k^3-(2+w)k^2+(2w+1)k-w.
\end{align*}
Accordingly $w=-2$ and so $Tr(F')=0=Tr(F)=z$. Hence $\det(F)=0$, contradiction.
\end{proof}

\section{Computational results}



\subsection{The group $G_M$}


%\comment{questo elenco non mi pare serva più.  }

%\begin{enumerate}
%    \item If $G=\langle X\,|\,R\rangle$ and $H=\langle Y\,|\,S\rangle$ are two presentations for groups $G$ and $H$, then a presentation for the semidirect product $G\rtimes_\varphi H$ is 
%    \[
       % G\rtimes_\varphi H=\langle X,Y \,|\, R,S, yxy^{-1}=\varphi_y(x)\ \forall x\in X\,\forall y\in Y\rangle.
%    \]

%    \item If $G=\langle X\,|\,R\rangle$ and $H=\langle Y\,|\,S\rangle$ are two presentations for groups $G$ and $H$, then an isomorphism between $G$ and $H$ is a map sending generators of $G$ in generators of $H$ (not necessarily elements of $X$ in elements of $Y$) and preserving the relations of both $G$ and $H$. \comment{magari farei un lemma con tutti i dettagli.} \fs{d'accordo, per il momento avevo solo elencato le cose che mi servivano}

%    \item It is known that $\aut{\mathbb{Z}_p^k}=GL(k,\mathbb{Z}_p)$.

 %   \item The matrix $I-N$ is invertible if and only if $\lambda=1$ is not an eigenvalue of $N$. \comment{Note that $1-N$ is invertible iff $1-N^h$ is invertible for $h\neq 0$.}

%    \item The eigenvalues of $N^t$ are the $t$-th power of the eigenvalues of $N$.
%\end{enumerate}

Let $G_{M}=\mathbb{Z}_2^{4+k}\rtimes_M\mathbb{Z}_p$ where the operation on $G_M$ is given by $(x,y)\ast (u,v)=(x+M^y u,y+v)$, for $x,u\in\mathbb{Z}_2^{4+k}$ and $y,v\in\mathbb{Z}_p$. We can consider $G_M$ as a group with the following presentation
    \[
        G_M=\langle u_1,\dots u_{4+k},a\,|\, u_i^2=1,\, u_iu_j=u_ju_i,\, a^p=1,\ au_ia^{-1}=Mu_i\rangle
    \]
More precisely, we will see generators of the semidirect product as $u_i=(u_i,0)$ and $a=(0,1)$, where the operation on the first component is the sum of vectors of $\mathbb{Z}_2^{4+k}$ and the operation on the second component is the sum in $\mathbb{Z}_p$.




\begin{theorem} \label{thm:GM_char}
The map $f:G_M\longrightarrow G_N$ defined by setting:
    \begin{align*}
        f(0,1)=(w,h),\quad f(u,0)=(F(u),0)
    \end{align*}
    where $u,w\in \Z_2^{4+k}$ and $F\in GL_{4+k}(\Z_2)$, can be extended to a group isomorphisms between $G_M$ and $G_N$ if and only if $FMF^{-1}=N^h$ and $h\neq 0$. In particular, $G_M\cong G_N$ if and only if $\langle N\rangle$ and $\langle M\rangle$ are conjugate in $GL_{4+k}(\Z_2)$.
\end{theorem}

%\comment{Data una mappa definita sui generatori. La sua estensione "lineare" è ben definita se mappa relazioni in relazioni.}

\begin{proof}
    Let $\setof{(u_i,0), (0,1)}{i=0,\ldots 4+k}$ be a set of generators of $G_M$. 

    Assume that the map $f:G_M\longrightarrow G_N$ can be extended to an isomorphism. Then $f$ maps the unique $2$-Sylow subgroup of $G_M$ onto the unique $2$-Sylow subgroup of $G_N$. In particular the restriction of $F$ to the $2$-Sylow subgroup is $F\in GL_{4+k}(\Z_2)$. Assume that the image of $(0,1)$ under $f$ is $(w,h)$. Then $f((0,1) (u,0)(0,1)^{-1})=(F M(u),0)=(w,h) f(u,0) (w,h)^{-1}=(N^hF(u),0)$. 

    On the other hand, note that $\setof{(F(u_i),0), (w,h)}{i=1,\ldots,4+k}$ is a set of generators of $G_N$ for every $h\neq 0$. If $f((0,1)(u_i,0)(0,1)^{-1})=(FM(u_i),0)=f(0,1) f(u_i,0)f(0,1)^{-1}=(N^hF(u_i),0)$, namely $FMF^{-1}=N^h$ then $f$ maps the relations of $G_N$ to relations of $G_M$ with respect to these sets of generators. Thus, $f$ can be extended to a group isomorphism.

    \comment{add relation for (w,h) to have right order} \fs{The condition for $(w,h)^p=(0,0)$ in $G_N$ is $(I+N^h+N^{2h}+\dots+N^{(p-1)h})w=0$.} 

    \fs{QUESTO CI SERVE ANCORA?}
\end{proof}




%\begin{proof}
  %  Let $f\colon G_M\to G_N$ be an isomorphism between groups $G_M$ and $G_N$, then it is a map sending generators in generators and respecting the relations. Note that $f$ restricted to $\mathbb{Z}_2^{4+k}$ must be an automorphism of $\mathbb{Z}_2^{4+k}$, and similarly $f$ restricted to $\mathbb{Z}_p$ must be an automorphism of $\mathbb{Z}_p$, therefore there is $F\in GL(4+k,\mathbb{Z}_2)$ and $s\in\mathbb{Z}_p$ such that $f(u_i,0)=(Fu_i,h_i)$ \comment{Il $2$-Sylow è caratteristico per cui sappiamo già che $h_i=0$} and $f(0,1)=(w,h)$ for some $h_i,h\in\mathbb{Z}_p$ and $w\in \mathbb{Z}_2^{4+k}$. By imposing that $f$ respect the relations, we obtain necessary conditions on  $w,h,h_i$.
  %  \begin{itemize}
   %     \item $f(2u_i,0)=(0,0)$ is equivalent to $(Fu_i+N^{h_i}Fu_i,2h_i)=(0,0)$, that is $h_i=0$ (as $p$ is odd) and the first condition turns out to be trivially true always. 
    %    \item $f$ trivially respects the commutativity of the $u_i$ as the operation involves only the first component, thanks to the previous step.
     %   \item $f((0,1)^p)=(w,h)^p=(0,0)$ implies that $((I+N^h+N^{2h}+\dots +N^{k(p-1)})w,ph)=(0,0)$, which is equivalent to $(I+N^h+N^{2h}+\dots +N^{h(p-1)})w=0$, as the second condition always holds. Note that the matrix 
      %  \[
       %     I+N^h+N^{2h}+\dots +N^{h(p-1)}=(I-N^h)^{-1}(I-N^{hp})
       % \]
        %if and only if $I-N^h$ is invertible, hence if and only if $\lambda=1$ is not an eigenvalue of $N$, which is always the case because $N$ is a block diagonal matrix and each block is irreducible. Therefore $I+N^h+N^{2h}+\dots +N^{h(p-1)}$ is again invertible, being a product of invertible matrices, and this pushes to $w=0$. \comment{Se $h\neq 0$ e $N^h$ ha autovalore $1$. Dato che $N=N^{hl}$ allora $1$ è anche autovalore di $N$. Quindi $1-N^h$ è  invertibile, quindi viene che $I+N^h+N^{2h}+\dots +N^{h(p-1)}=(I-N^h)^{-1}(I-N^{hp})=0$ perchè $N^h$ ha ordine $p$. Quindi hai un automorfismo per ogni scelta di $w$ che non per forza è $0$. Anzi hai $2^{4+k}$ scelte per $w$.}
    %    \item $f((0,1)\ast (u_i,0)\ast (0,p-1))=f(Mu_i,0)=(FMu_i,0)$ implies that $(N^hFu_i,p)=(FMu_i,0)$ hence $FMF^{-1}=N^h$.
   %     In conclusion, the map $f$ acts on the generators as $f(u_i,0)=(Fu_i,0)$ and $f(0,1)=(0,h)$, therefore, by linearity in the first component and by the operation in the semidirect product, we obtain that $f(u,t)=(Fu,ht)$.
%    \end{itemize}
%\end{proof}




\begin{corollary}\label{aut of G_M}
    The automorphism group of $G_M$ are the maps defined by setting:
\begin{align*}
    f(0,1)=(w,h),\, \quad f(u,0)=(F(u),0),
\end{align*}
where $w\in \Z_2^{4+k}$, $h\in \{1,\ldots,p-1\}$ and $F=f|_{\Z_2^{4+k}}\in GL_{4+k}(\Z_2)$ such that $FMF^{-1}=M^h$. %In particular $|\aut{G_M}|=2^{4+k}|N_{GL_{4+k}(\Z_2)}(M)|$. 
\end{corollary}

%\begin{lemma} \label{latin_computational}
 %   Let $f\in\aut{G_M}$ be constructed with a matrix $F\in N_{GL_{4+k}(2)}(M)$, as in Corollary \ref{aut of G_M}, and let $H=\fix(f)$. Then the coset quandle $\mathcal{Q}(G_M,H,f)$ is latin if and only if $h\neq 1$ and the condition 
  %  \[
   % M^{(h-1)x}(I - M^x F M^{-x})v \in H \implies v\in H
   % \]
   % holds every $x\in \{2,\dots, p-1\}$, and every vector $v\in \Z_p^{4+k}$.
%\end{lemma}

%\begin{proof}
 %   The condition $h \neq 1$ is equivalent to saying that $F$ belongs to the normalizer $N_{GL_{4+k}(2)}(M)$ but not to the centralizer $C_{GL_{4+k}(2)}(M)$. The second condition follows by expanding and directly computing the latin statement concerning the right multiplication maps: $R_H(xH) = R_H(yH)$ implies $xH = yH$.
%\end{proof}


In order to compute all the latin quandles of size $16p$ with congruence lattice given by a chain of length $3$ we cam proceed as follows: 

\begin{itemize}
    \item[(i)]  according to Lemma \ref{LSS 16p} such quandles are coset quandles over a group of the form $G_{p,k}=\Z_2^{4+k}\rtimes_\rho \Z_p$ for $0\leq k\leq 4$ and some $p\in \{3,5,7,17,31,127\}$ such that $Z(G_{p,k})=1$. Then, according to Proposition \ref{trivial center}, $e(p,2)$ need to divide $4+k$. This condition, allow us to identify the groups in the SmallGroup library of GAP as follows:   
    

\begin{table}[h!]
    \centering
\begin{tabular}{|c|c|c|c|c|c|}
\hline
$p\ $\textbackslash$\ k$ & 0 &1 & 2 & 3 & 4 \\ 
\hline
3 & \text{(48,50)} & - & \text{(192,1541)} & - & \text{(768,1085321)} \\ 
\hline
5 & \text{(80,49)} & - & - & - & \text{(1280,1116356)} \\ 
\hline
7 & - & - & \begin{tabular}{c} (448,1393) \\ (448,1394) \end{tabular} & - & - \\ 
\hline
17 & - & - & - & - & \text{(4352,1118911)} \\ 
\hline
31 & - & \text{(992,194)} & - & - & - \\ 
\hline
127 & - & - & - & \text{(16256,19325)} & - \\ 
\hline
\end{tabular}
    \caption{Groups of the form $G_{p,k}$ with trivial center identified by their label in the SmallGroup library of GAP.}
\end{table}



    \item[(ii)] The second step is to construct all the coset quandles of the form $\mathcal{Q}(G_{p,k},Fix(f),f)$ over the groups above. Then we consider which one of them are latin of size $16p$ and we filter the isomorphism classes as is Theorem \ref{iso theo}. The following table display the number of such quandles:
    

\begin{table}[h!]
    \centering
\begin{tabular}{|c|c|c|c|c|c|}
\hline
$p\ $\textbackslash$\ k$ & 0 & 1 & 2 & 3 & 4 \\ 
\hline
3 & 1 & - & 0 & - & 0 \\ 
\hline
5 & 1 & - & - & - & 0 \\ 
\hline
7 & - & - & 0 & - & - \\ 
\hline
17 & - & - & - & - & 0\\ 
\hline
31 & - & 0 & - & - & - \\ 
\hline
127 & - & - & - & 0 & - \\ 
\hline
\end{tabular}
    \caption{Number of quandles of size $16p$.}
\end{table}
\end{itemize}




%Moreover we can use Theorem \ref{iso theo} to identify the isomorphism classes of such quandles.

%The previous theorems allow us to proceed with classification using only a computational algorithm in GAP. The implementation details are provided in the appendix. For each pair $(p, 4 + k)$ identified above, the steps of the algorithm are as follows:

%\begin{enumerate}
 %   \item The structure of the $2-$Sylow subgroups of $GL(4 + k, 2)$ is known and described in \cite{carter1964sylow}. The latter conditions, together with the fact that the group $G_M$ has trivial center, allow us to set some initial constraints, exclude all pairs $(p, 4+k)$ for which $k$ is not a multiple of $e(p,2)$ (which is $|2|$ modulo $p$, as in Section \ref{app:syl_of_GL} of the Appendix). 

  %  \item Generate candidate matrices $M$ by filtering for conjugacy class of generated cyclic group, order (equal to $p$), and size (equal to $4 + k$), while ensuring that the eigenvalues are different from 1 (otherwise, the center would not be trivial). This step is justified by Theorem \ref{thm:GM_char}.

   % \item Construct the groups $G_M$ by taking semidirect products with the candidate matrices $M$ identified in the previous step.

    %\item If the order of $G_M$ is computationally small enough, find all automorphisms $\alpha$ and directly test whether the coset quandle $\mathcal{Q}(G_M, \fix(\alpha), \alpha)$ is latin of size $16p$, where $\fix(\alpha)$ denotes the fixed-point subgroup under $\alpha$. 

   % \item Find explicitly the matrices $F$ as described in Corollary \ref{aut of G_M}. Specifically, consider all matrices in $GL(4 + k, 2)$ that belong to the normalizer of $M$ but not its centralizer and for which the subspace of fixed points has dimension $k$ (and thus size $2^k$). For each such matrix, test directly whether the coset quandle it defines is latin using the condition determined in Lemma \ref{latin_computational}. Note that, by general quandle theory, it is known that $\dis(Q)$ is isomorphic to $G_M$.
%\end{enumerate}

%In the following, we present tables listing all candidate groups $G_M$, all candidate matrices $M$, and the existence of latin quandles of size $16p$ (in GAP notation) constructed in this way.


%\begin{center}
%\textbf{Groups $G_M$}
%\begin{tabular}{|c|c|c|c|c|c|}
%\hline
%$p\ $\textbackslash$\ 4+k$ & 4 & 5 & 6 & 7 & 8 \\ 
%\hline
%3 & \text{(48,50)} & - & \text{(192,1541)} & - & \text{(768,1085321)} \\ 
%\hline
%5 & \text{(80,49)} & - & - & - & \text{(1280,1116356)} \\ 
%\hline
%7 & - & - & \begin{tabular}{c} (448,1393) \\ (448,1394) \end{tabular} & - & - \\ 
%\hline
%17 & - & - & - & - & \text{(4352,1118911)} \\ 
%\hline
%31 & - & \text{(992,194)} & - & - & - \\ 
%\hline
%127 & - & - & - & \text{(16256,19325)} & - \\ 
%\hline
%\end{tabular}
%\end{center}


%\begin{center}
%\textbf{Quandles of size $16p$}

%\begin{tabular}{|c|c|c|c|c|c|}
%\hline
%$p\ $\textbackslash$\ 4+k$ & 4 & 5 & 6 & 7 & 8 \\ 
%\hline
%3 & 1 & X & 0 & X & 0 \\ 
%\hline
%5 & 1 & X & X & X & 0 \\ 
%\hline
%7 & X & X & 0 & X & X \\ 
%\hline
%17 & X & X & X & X & 0\\ 
%\hline
%31 & X & 0 & X & X & X \\ 
%\hline
%127 & X & X & X & 0 & X \\ 
%\hline
%\end{tabular}
%\end{center}

%\begin{sidewaystable}
%\centering
%\textbf{Matrices $M$}

%\begin{tabular}{|c|c|c|c|c|c|}
%\hline
%$p\ $\textbackslash$\ 4+k$ & 4 & 5 & 6 & 7 & 8 \\ 
%\hline
%3 & $\begin{bmatrix}
%0 & 1 & 0 & 0 \\
%1 & 1 & 0 & 0 \\
%0 & 0 & 0 & 1 \\
%0 & 0 & 1 & 1
%\end{bmatrix}$ 
%& X &
%$\begin{bmatrix}
%0 & 1 & 0 & 0 & 0 & 0 \\
%1 & 1 & 0 & 0 & 0 & 0 \\
%0 & 0 & 0 & 1 & 0 & 0 \\
%0 & 0 & 1 & 1 & 0 & 0 \\
%0 & 0 & 0 & 0 & 0 & 1 \\
%0 & 0 & 0 & 0 & 1 & 1
%\end{bmatrix}$
%& X & 
%$\begin{bmatrix}
%0 & 1 & 0 & 0 & 0 & 0 & 0 & 0 \\
%1 & 1 & 0 & 0 & 0 & 0 & 0 & 0 \\
%0 & 0 & 0 & 1 & 0 & 0 & 0 & 0 \\
%0 & 0 & 1 & 1 & 0 & 0 & 0 & 0 \\
%0 & 0 & 0 & 0 & 0 & 1 & 0 & 0 \\
%0 & 0 & 0 & 0 & 1 & 1& 0 & 0 \\
%0 & 0 & 0 & 0 & 0 & 0 & 0 & 1 \\
%0 & 0 & 0 & 0 & 0 & 0 & 1 & 1 
%\end{bmatrix}$
%\\[10pt]
%\hline
%5 & 
%$\begin{bmatrix}
%0 & 0 & 0 & 1 \\
%1 & 0 & 0 & 1 \\
%0 & 1 & 0 & 1 \\
%0 & 0 & 1 & 1
%\end{bmatrix}$
%& X & X & X & $\begin{bmatrix}
%0 & 0 & 0 & 1 & 0 & 0 & 0 & 0 \\
%1 & 0 & 0 & 1 & 0 & 0 & 0 & 0 \\
%0 & 1 & 0 & 1 & 0 & 0 & 0 & 0 \\
%0 & 0 & 1 & 1 & 0 & 0 & 0 & 0 \\
%0 & 0 & 0 & 0 & 0 & 0 & 0 & 1 \\
%0 & 0 & 0 & 0 & 1 & 0 & 0 & 1 \\
%0 & 0 & 0 & 0 & 0 & 1 & 0 & 1 \\
%0 & 0 & 0 & 0 & 0 & 0 & 1 & 1
%\end{bmatrix}$ 
%\\[10pt]
%\hline
%7 & X & X & 
%$\begin{bmatrix}
%0 & 0 & 1 & 0 & 0 & 0 \\
%1 & 0 & 0 & 0 & 0 & 0 \\
%0 & 1 & 1 & 0 & 0 & 0 \\
%0 & 0 & 0 & 0 & 0 & 1 \\
%0 & 0 & 0 & 1 & 0 & 0 \\
%0 & 0 & 0 & 0 & 1 & 1
%\end{bmatrix}
%\begin{bmatrix}
%1 & 1 & 0 & 1 & 0 & 1 \\
%0 & 1 & 1 & 0 & 0 & 1 \\
%1 & 0 & 1 & 1 & 1 & 1 \\
%0 & 0 & 1 & 0 & 0 & 1 \\
%1 & 1 & 0 & 1 & 1 & 1 \\
%1 & 1 & 0 & 0 & 1 & 0
%\end{bmatrix}
%\begin{bmatrix}
%0 & 1 & 1 & 1 & 0 & 0 \\
%1 & 1 & 1 & 0 & 1 & 1 \\
%0 & 0 & 1 & 0 & 1 & 0 \\
%1 & 0 & 0 & 0 & 0 & 1 \\
%1 & 0 & 1 & 1 & 1 & 0 \\
%1 & 1 & 0 & 0 & 0 & 0
%\end{bmatrix}$
%& X & X \\
%\hline
%17 & X & X & X & X & 
%$\begin{bmatrix}
%1 & 0 & 0 & 0 & 1 & 1 & 1 & 0 \\
%0 & 1 & 1 & 0 & 0 & 1 & 1 & 1 \\
%0 & 0 & 1 & 1 & 1 & 0 & 1 & 1 \\
%0 & 0 & 0 & 0 & 0 & 0 & 1 & 1 \\
%0 & 1 & 0 & 0 & 0 & 0 & 0 & 0 \\
%1 & 0 & 1 & 1 & 0 & 1 & 1 & 1 \\
%1 & 1 & 0 & 1 & 1 & 0 & 1 & 1 \\
%0 & 0 & 1 & 0 & 1 & 0 & 0 & 1
%\end{bmatrix}
%\begin{bmatrix}
%1 & 1 & 1 & 0 & 1 & 1 & 0 & 0 \\
%0 & 0 & 1 & 1 & 0 & 0 & 0 & 1 \\
%1 & 1 & 0 & 1 & 0 & 0 & 0 & 1 \\
%0 & 0 & 0 & 0 & 1 & 1 & 1 & 0 \\
%1 & 0 & 1 & 0 & 0 & 0 & 1 & 1 \\
%1 & 0 & 0 & 1 & 1 & 0 & 1 & 1 \\
%1 & 1 & 0 & 1 & 1 & 1 & 1 & 0 \\
%1 & 0 & 1 & 1 & 1 & 0 & 0 & 1
%\end{bmatrix}$\\
%\hline
%31 & X & 6 matrices & X & X & X\\
%\hline
%127 & X & X & X & 18 matrices & X\\
%\hline
%\end{tabular}
%\end{sidewaystable}%








\end{document}


    \subsection{Gap algorithm} In this subsection, we describe the code used to obtain the table.

\begin{lstlisting}
matrixCandidates := function(G,p)

local CC, x;

CC := ConjugacyClasses(G);
CC := List(CC,x->Random(x));
CC := Filtered(CC,x->Order(x)=p);
CC := Filtered(CC,x-> not (Z(2)^0 in Eigenvalues(GF(2),x)));

return CC;

end;
\end{lstlisting}

The function \texttt{matrixCandidates} takes as input a matrix group G (in this case G=GL(4+k,2) ) and a prime $p$, and it returns one representative per conjugacy class that has order $p$ and for which $1$ is not an eigenvalue.


\begin{lstlisting}
islatin16p := function(G,p)

local A, CC, repr, goodAut, a, Q;

A := AutomorphismGroup(G);
CC := ConjugacyClasses(A);
repr := List(CC,x->Random(x));
goodAut := [];

for a in repr do
    Q := cosetQuandle(G,fixedPoints(G,a),a);
    if Size(Q)=16*p and Islatin(Q) then
        Add(goodAut,a);
    fi;
od;

return goodAut;

end;
\end{lstlisting}

The function \texttt{islatin16p} takes as input a group $G$ and a prime $p$ and it returns all the automorphisms $\alpha$ of $G$ such that $\mathcal{Q}(G,\fix(\alpha),\alpha)$ is a coset latin quandle of size $16p$.

\begin{lstlisting}
normalizerMatrixCandidates := function(M,k)

local C,matrices,N, eigen,dimensions,x,y,fix, check,goodMatrices;

N := Normalizer(GL(Size(M),2),M);
C := Centralizer(GL(Size(M),2),M);
matrices := Difference(N,C);
matrices := Filtered(matrices,x -> Size(Eigenspaces(GF(2),x))>0);
matrices := Filtered(matrices,x -> Dimension((Eigenspaces(GF(2),x))[1])=k);

I := IdentityMat(4+k)*One(GF(2));
V := GF(2)^(4+k);
goodMatrices := [];

for F in matrices do
    
    check := true;
    fix := Eigenspaces(GF(2),F)[1];

    for x in V do
        for y in V do
            if (y-x)^(I-F) in fix and not y-x in fix then
                check := false;
            fi;
        od;
    od;
    
    if check = true then
        Add(goodMatrices,F);
    fi;
od;

return goodMatrices;
end;
\end{lstlisting}


The function \texttt{normalizerMatrixCandidates} takes as input a matrix $M$ and an integer $k$ and returns all the matrices $F$ normalizing $M$, according to the statement of Theorem \ref{thm:GM_char}.This has been used to test manually if a group $G_M$ gives rise to a latin coset quandle of size $16p$ when the size of $G_M$ was computationally too large to apply the previous method.
We denote an element of $f\in\aut{G_M}$ by $f=f_{F,h}$ meaning that $f(u,t)=(Fu,ht)$ and $FMF^{-1}=M^h$, as in the previous proof.

\begin{proposition}[Group structure of $\aut{G_M}$] Let $f_{F,h}$ and $f_{G,r}$ be elements of $\aut{G_M}$. Then
\begin{enumerate}
    \item $f_{I,1}$ is the identity of $\aut{G_M}$.
    \item $f_{F,h}^{-1}=f_{F^{-1},h^{-1}}$.
    \item $f_{F,h}f_{G,r}=f_{FG,hr}$.
    \item $f_{F,h}f_{G,r}f_{F,h}^{-1}=f_{FGF^{-1},r}$
\end{enumerate}
\end{proposition}

\begin{proof}
    Proofs are all straightforward. We have to check the compositions and the property $FMF^{-1}=M^h$.
\end{proof}

\begin{corollary}
    A conjugation class in $\aut{G_M}$ is completely determined by a conjugation class of $GL(4+k,\mathbb{Z}_2)$ and an element $s\in\mathbb{Z}_p$.
\end{corollary}

Let $Q_{M,k,F,h}$ be the coset quandle $\mathcal{Q}(G_M,Fix(f_{F,h}),f_{F,h})$.

\begin{proposition}
    $Fix(f_{F,h})=1$ for every choice of $F$ and $h$, consequently, $Q_{M,k,F,h}$ is always a principal quandle.
\end{proposition}


\begin{proof}
    For $h\neq 1$, the automorphism $f_{F,h}$ has no fixed points, while for $h=1$, the automorphism $f_{F,1}=F$ is simply an automorphism of $GL(4+k,\mathbb{Z}_2)$. In this second case, $Fu=u$ if and only if $u$ is an eigenvector associated with the eigenvalue 1, which does not exist for the same reasons as above. \comment{Noi sappiamo che $M$ non ha autovalori uguali a $1$ perchè non c'è centro. Gli automorfismi de gruppo $G_M$ ristretti a $\Z_2^{4+k}$ possono avere autovalori pari a $1$. Chi sono le matrici che normalizzano $\langle N\rangle$?}
\end{proof}

\begin{corollary}
    $|Q_{M,k,F,h}|=16p$ only if $k=0$.
\end{corollary}

\begin{proof}
    As the quandle $\mathcal{Q}(G_M,Fix(f_{F,h}),f_{F,h})$ is principal for every choice of the automorphism, we have $|Q_{M,k,F,h}|=|G_M|=2^{4+k}p$ and this is $16p$ only when $k=0$.
\end{proof}


\begin{proposition}
    The quandle $Q_{M,k,F,h}$ is latin if and only if $I+F$ is invertible and $h\neq -1$.
\end{proposition}

\begin{proof}
    As the quandle $Q_{M,k,F,h}$ is homogeneous, it is latin if and only if $R_{(0,0)}$ is invertible. By an explicit computation, we obtain that the map $R_{(0,0)}(u,t)=((I+F)u,(h+1)t)$, hence condition in the statement.
\end{proof}


Let $G=\mathbb{Z}_p^{t}\rtimes_\rho\mathbb{Z}_4^2$, where $\rho\colon\mathbb{Z}_4^2\to\aut{\mathbb{Z}_p^{t}}\cong GL(t,p), \rho(v)=T_v$, and operation on $G$ is given by $(u,v)\ast (w,z)=(u+T_v(w),v+z)$. We can consider $G$ as a group with the following presentation
    \[
        G_M=\langle u_1,\dots u_t, v_1,v_2\,|\, u_i^p=1,\, u_iu_j=u_ju_i,\, v_j^4=1,\, v_1v_2=v_2v_1,\, v_1u_iv_1^{-1}=T_{v_1}(u_i),\, v_2u_iv_2^{-1}=T_{v_2}(u_i)\rangle
    \]
More precisely, we will see generators of the semidirect product as $u_i=(u_i,0)$ and $v_j=(0,v_j)$, where the operation on the first component is the sum of vectors of $\mathbb{Z}_p^{t}$ and the operation on the second component is the sum in $\mathbb{Z}_4^2$. With a technique similar to the proof of the previous theorem, we can prove the following

\fs{the following holds assuming that $\lambda=1$ is not an eigenvalue of $T_{Gv_j}$ for $j=1,2$}.

\begin{theorem}
    Let $H=\mathbb{Z}_p^t\rtimes_\rho\mathbb{Z}_2^4$. Then $\aut{H}$ is made by all the maps $f\colon G\to G$ such that there are matrices $F\in GL(t,p)$ and $G\in GL(2,4)$ such that $FT_vF^{-1}=T_{Gv}$ for all $v\in\mathbb{Z}_4^2$. In this case, the automorphisms have the following form $f(u,v)=(Fu,Gv)$.
\end{theorem}
\comment{Secondo me il $2$-Sylow non viene per forza caratteristico come si deduce dalla forma che ottieni tu degli automorfismi. C'è lo stesso errore dovuto a "$w$" come nell'altro caso secondo me.}


\newpage


\begin{lemma}\comment{check and finish!!!}
Let $Q$ be a subdirectly reducible latin quandle of size $4p$ and let $|Q/\alpha|=p$. If $\alpha$ is not the unique maximal congruence of $Q$ then $Q$ is affine.
\end{lemma}

\begin{proof}
If $Q$ has a maximal congruence of size $16$ then it is the direct product. Assume that $Q/\beta$ has size $4$. Then $Q/\alpha\wedge/\beta$ is affine and it has size $4p$ and the blocks have size $4$. Therefore $\alpha\wedge\beta$ is a minimal congruence. Hence $\alpha\wedge\beta\wedge\gamma=0_Q$ for every $\gamma\neq \alpha,\beta$. If $\gamma\leq \alpha$ then it has blocks of size $4$ and then $Q\cong Q/\gamma\times Q/\alpha$ 
\end{proof}



\begin{lemma}
Let $Q$ be a non-affine latin quandle of size $16p$, $p>5$. Then either $Q$ has a unique maximal congruence or the simple factors of $Q$ have size $4$ or $p$.
\end{lemma}
\begin{proof}
The quandle $Q$ has no factor of size $2$ and $8$. If $Q$ has two maximal congruences $\alpha, \beta$, then $|Q/\alpha\wedge\beta|=|Q/\alpha||Q/\beta|$ and it divides $16p$. If the size of $Q/\alpha$ is $16$ then the size of $|Q/\beta|=p$ and then $Q$ is the direct product. 
\end{proof}



\begin{lemma}
Let $Q$ be a non-affine latin quandle of size $16p$. If $Q$ has a unique maximal congruence $\alpha$	with factors of size $p>5$. Then $\alpha$ is the unique proper congruence of $Q$
\end{lemma}

\begin{proof}
It $Q$ has a congruence $\beta\leq \alpha$ Then $Q/\beta$ has size $4p$ and it has a subdirectly irreducible factor of size $4p$. There are no subdirectly irreducible latin quandles of size $4p$ with a factor of size $p$ \cite[]{LSS}. 
\end{proof}

